% DER_I-b: Grothendieck's text with the commentary set in, then the overview, inside a group that applies the chunk's own definitions
\dgrpart{1}{I-b}{Part 2 of 2, manuscript pages 67--126}{I-b\enspace\textperiodcentered\enspace Chapter I, part 2 of 2\enspace\textperiodcentered\enspace manuscript pages 67--126}{Chapter I \textperiodcentered\ \textup{I-b}}
\dgrchunkbegin
% DER_I-b: the definitions of the chunk's preamble, applied inside its group
\csname tagsleft@true\endcsname
\let\printedpage\relax\newcommand{\printedpage}[1]{\marginnote{\scriptsize\sffamily ed.\ #1}}
\let\msp\relax\newcommand{\msp}[1]{%
  \begingroup\def\thismsp{#1}\def\blankmsp{page 114 vide}%
  \ifx\thismsp\blankmsp
    \marginnote{\scriptsize\sffamily [114: vide]}[1.3\baselineskip]%
  \else\def\followingmsp{page 115}%
    \ifx\thismsp\followingmsp
      \marginnote{\scriptsize\sffamily [page 115]}[2.6\baselineskip]%
    \else\marginnote{\scriptsize\sffamily [#1]}[1.3\baselineskip]\fi
  \fi\endgroup}
\let\DD\relax\newcommand{\DD}{\ensuremath{\mathbb{D}}}
\let\cA\relax\newcommand{\cA}{\ensuremath{\mathcal{A}}}
\let\Dia\relax\newcommand{\Dia}{\ensuremath{\mathsf{Dia}}}
\let\Ob\relax\newcommand{\Ob}{\ensuremath{\mathsf{Ob}}}
\let\Hom\relax\newcommand{\Hom}{\ensuremath{\mathsf{Hom}}}
\let\hot\relax\newcommand{\hot}{\ensuremath{\mathsf{hot}}}
\let\Hot\relax\newcommand{\Hot}{\ensuremath{\mathsf{Hot}}}
\let\HOT\relax\newcommand{\HOT}{\ensuremath{\mathsf{HOT}}}
\let\isoto\relax\newcommand{\isoto}{\overset{\sim}{\to}}
\let\isolongto\relax\newcommand{\isolongto}{\overset{\sim}{\longrightarrow}}
\let\D\relax\newcommand{\D}{\ensuremath{\mathbb{D}}}
\let\Fl\relax\newcommand{\Fl}{\ensuremath{\mathsf{Fl}}}
\let\Cat\relax\newcommand{\Cat}{\ensuremath{\mathcal{C}at}}
\let\bs\relax\newcommand{\bs}{\backslash}
\let\RR\relax\newcommand{\RR}{\ensuremath{\mathbb{R}}}
\let\sfup\relax\newcommand{\sfup}[1]{\textsf{\upshape #1}}

% ---- from der1_p041-045.tex
\printedpage{39}
% head: A. Grothendieck : Dérivateurs, Ch. I, Édition des Universités de Montpellier II et Paris VII
% foot: Version 8 juillet 2001
\msp{page 67}

\noindent Donc la flèche envisagée $(**)$ s'interprète comme
\[
H_*^{\D}(y\backslash X, M_X) \longrightarrow M
\]
et c'est l'homomorphisme d'adjonction $p_!(p^*(M)) \to M$.
\[
\begin{tikzcd}
X \arrow[rr, "f"] \arrow[dr, "p"{swap}] & & Y \arrow[dl, "q"] \\
& e &
\end{tikzcd}
\]

\noindent Donc c'est un isomorphisme pour tout $M$ si et seulement si $y\backslash X$ est $\D$-contractile.

\noindent On va énoncer ce qu'on a prouvé.

\noindent\textbf{Théorème 1.} \textit{Soit}
\[
\begin{tikzcd}
X \arrow[rr, "f"] \arrow[dr, "p"{swap}] & & Y \arrow[dl, "q"] \\
& e &
\end{tikzcd}
\]
\textit{dans} $\Dia$. \textit{On suppose que} $\D$ \textit{est un prédérivateur satisfaisant} \textbf{Der 2}, \textit{que} $p_*, q_*$ \textit{et} $f_!$ \textit{existent, que} $f_!$ \textit{satisfait à} \textbf{Der 4} (d) \textit{pour le calcul des fibres. Conditions équivalentes :}

a) $f$ \textit{est une} $\D$\textit{-équivalence costricte, i.e.} (prop. 2) \textit{pour tout} $y \in \Ob\, Y$, $y\backslash X$ \textit{est} $\D$\textit{-contractile ;}

b) \textit{pour tout} $F \in \Ob\, \D(Y)$, \textit{l'homomorphisme suivant est} \texttt{[un]} \textit{isomorphisme}
\[
H^*_{\D}(Y, F) \isoto H^*_{\D}(X, f^*(F)) \ ;
\]

c) \textit{pour tout} $M \in \Ob(\mathcal{A} = \D(e))$, \textit{l'homomorphisme d'adjonction est} \texttt{[un]} \textit{isomorphisme}
\[
f_!(M_X) = f_!(f^*(M_Y)) \isoto M_Y \ .
\]

\dgrblock{dgrb37}{The conclusion reached on manuscript page 67}
\begin{dcomment}{exposition}{p67:f_!(M_X) = f_!(f^*(M_Y))}{The conclusion reached on manuscript page 67}
Manuscript pages~60--66 are commented with the first part of the chapter; its blocks ``Equivalence measured by constant coefficients'', ``How much constant coefficients can detect'', ``Local isomorphisms and local constancy'', and ``Cofinality and the final Yoneda argument'' give the definitions and the argument now being completed. The distinction needed below is quite exact. A $\D$-equivalence preserves the morphisms between constant objects; a coaspherical functor preserves cohomology with \emph{every} coefficient object on its target. Neither assertion, by itself, says that inverse image is fully faithful on the entire coefficient category.

For the theorem on manuscript page~67, the relevant test is the counit
\[
f_!p^*M\longrightarrow q^*M,\qquad p=qf.
\]
Applying $\operatorname{Hom}(-,F)$ and then the two adjunctions identifies this test with cohomological invariance for all $F$. Evaluation at $y$ identifies it instead with the augmentation of homology of $(y\backslash f)$ with constant value $M$. These are two readings of the \emph{same natural transformation}. This observation explains the appearance of homology in a theorem about cohomology; it also prevents a spurious gap from being introduced into an otherwise complete Yoneda argument. The theorem explicitly supplies the right adjoints to the terminal projections, whose existence was less explicit in the proposition on page~66.
\end{dcomment}


Je vais renuméroter les axiomes sur le prédérivateurs. % [sic] le prédérivateurs -> les prédérivateurs, article singulier devant un pluriel

\noindent\textbf{Der 1}, \textbf{Der 2} comme avant (p. 45).

\noindent\textbf{Der 3}\footnote[28]{Axiome des catégories d'indices asphériques, \textit{cf.} variantes \textbf{Der 3} a) b) c) page 97, et surtout pages 111-112.} : Soit $X$ dans $\Dia$, $p : X \to e$. Supposons $X$ contractile.\footnote[29]{N.B. Il suffit de supposer que $X$ a un objet final. Mais ce n'est pas un axiome autodual.} Alors
\[
p^* : M \longmapsto M_X \quad , \qquad p^* : \mathcal{A} = \D(e) \longrightarrow \D(X)
\]
est pleinement fidèle.

\printedpage{40}
% head: A. Grothendieck : Dérivateurs, Ch. I, Édition des Universités de Montpellier II et Paris VII
% foot: Version 8 juillet 2001
\msp{page 68}

\noindent\textbf{Der 4} (g) : existence des $f_*$.

\noindent\textbf{Der 4} (d) : existence des $f_!$.

\noindent\textbf{Der 5} (g) Calcul des fibres de $f_*(F)$.

\noindent\textbf{Der 5} (d) Calcul des fibres de $f_!(F)$.

\noindent L'ancien axiome \textbf{Der 5} (g) résulte de \textbf{Der 2}, \textbf{Der 3}, \textbf{Der 4} (g,d), \textbf{Der 5} (d), grâce au théorème précédent.

\dgrblock{dgrb38}{Keeping the axiom numbers separate}
\begin{dcomment}{exposition}{p68:grâce au théorème précédent.}{Keeping the axiom numbers separate}
Grothendieck is reorganising the assumptions after the failure diagnosed on manuscript page~60. The new numbering must be kept separate from both the preceding pages and present-day derivator terminology. Here \textbf{Der 3} asserts that a contractible indexing category has fully faithful constant-diagram functor. The two parts of \textbf{Der 4} assert existence of right and left Kan extensions; the two parts of \textbf{Der 5} give their pointwise comma-category formulas. In the now common numbering, existence is usually (Der3), and the pointwise formulas are (Der4); compare \citet[Definition~1.5]{Groth2013}. The extra strongness axiom used to lift incoherent arrows is a different matter and is not being imposed by this renumbering.

Throughout these comments we retain the manuscript's presheaf convention: a coefficient object over $X$ has underlying diagram $X^{\mathrm{op}}\to\D(e)$. Thus
\[
 (f_*F)_y\simeq H^*_{\D}(X/y,F),\qquad
 (f_!F)_y\simeq H^{\D}_*(y\backslash X,F).
\]
A modern reference using covariant diagrams reverses the comma categories. The right adjoint is still the right adjoint; the difference comes from replacing every indexing category by its opposite.

The new contractibility axiom is not an assertion that the whole category $\D(X)$ collapses to $\D(e)$. It concerns only constant objects and their morphisms. For example, over the walking arrow a coefficient is a coherent arrow, and its two values need not be isomorphic even though the walking arrow has contractible nerve. What the axiom guarantees is that taking a constant object over this shape introduces no new endomorphisms. This is precisely the amount of invariance used when the lifting categories in the following theorem are contractible.

The comment that both homological and cohomological axioms seem indispensable should consequently be read as a description of this proof strategy, not as a proved independence theorem. Modern treatments may begin with a one-sided derivator and formulate the corresponding one-sided conclusions. The argument here deliberately uses the opposite adjoint to make a universal-coefficient test accessible to pointwise calculation.
\end{dcomment}


Un aspect intéressant du théorème précédent, si on l'envisage comme énonçant l'équivalence de a) et b), c'est que cet énoncé ne fait pas appel à des notions homologiques, mais que sa démonstration y fait appel, sous une forme relativement forte (calcul des fibres de $f_!(F)$). Pour les théorèmes de changement de base pour la $\D$-cohomologie relative, \textit{i.e.} pour les $f_*^{\D}$, on a besoin aussi constamment du calcul des fibres des $f_*(F)$. Donc il semble qu'on ne peut se passer d'aucun des axiomes précédents, y compris \textbf{Der 4} (d) et \textbf{Der 5} (d), même si on ne veut s'intéresser qu'au formalisme des $f^*, f_*$ (aspect cohomologique). Par la suite, on va donc supposer ces axiomes satisfaits.

\dgrblock{dgrb39}{The category of lifts behind properness}
\begin{dcomment}{exposition}{p68:Par la suite, on va donc supposer ces axiomes satisfaits.}{The category of lifts behind properness}
There is a useful concrete form of the first two conditions of Theorem~2. Given $x\in X$ and $\alpha:f(x)\to y$, let $L_f(x,\alpha)$ have objects
\[
 h:x\longrightarrow z,\qquad f(z)=y,\qquad f(h)=\alpha.
\]
A morphism from $h$ to $h'$ is a map $k:z\to z'$ with $f(k)=1_y$ and $kh=h'$. These are strict lifting data; replacing equality by an unspecified isomorphism would change the condition. The theorem says that $f$ is $\D$-proper precisely when every $L_f(x,\alpha)$ is $\D$-contractible.

Indeed, if $i_y:X_y\to X/y$ is the canonical inclusion, an object of $X/y$ is exactly a pair $(x,\alpha)$ as above, and
\[
 ((x,\alpha)\backslash i_y)\cong L_f(x,\alpha).
\]
Condition~(a), applied to the definition of coasphericity, is therefore condition~(b). This also proves its stability under arbitrary strict base change. For a pullback $f':X\times_Y Y'\to Y'$, the lift of an arrow $\alpha':y'_0\to y'_1$ from $(x,y'_0)$ is the lift of its image in $Y$, together with the already fixed second coordinate $y'_1$. The resulting lifting categories are isomorphic.

The significance of these categories is stronger than the existence of lifts. An opfibration gives a preferred cocartesian lift, which is an initial object of $L_f(x,\alpha)$. Properness permits a whole category of choices, requiring only that constant coefficients cannot distinguish its homotopy from a point. In particular, requiring merely that each strict fibre $X_y$ be contractible misses the lifting condition altogether.

To obtain the base-change statement, first replace $Y$ by any $Y'\to Y$. At an object $y'$, the right Kan extension is calculated over $X'/y'$. By coasphericity of $X'_{y'}\to X'/y'$ and Theorem~1, restriction induces
\[
 H^*_{\D}(X'/y',F)\xrightarrow{\sim}
 H^*_{\D}(X'_{y'},F|_{X'_{y'}})
\]
for \emph{every} coefficient object $F$ on $X'/y'$. Consequently right direct image is calculated by the actual fibres. After a further base change, the actual fibre is unchanged. The two evaluations of the Beck--Chevalley comparison are therefore identified by their canonical restriction maps; \textbf{Der 2} makes the comparison an isomorphism.

Conversely, suppose the fibre formula holds after localisation $Y/y\to Y$. The category $Y/y$ has terminal object $1_y$, so evaluating the right direct image there identifies its global cohomology with that of $X/y$. The asserted fibre formula then states that restriction along $i_y$ preserves cohomology for every object of $\D(X/y)$. Theorem~1 gives the coasphericity of $i_y$. This completes the equivalence of (a), (b), (c), (c$'$), and (c$''$), without substituting constant coefficients where the argument needs arbitrary ones.
\end{dcomment}


\noindent\textbf{Théorème 2.} \textit{Soit} $f : X \to Y$ \textit{dans} $\Dia$. \textit{Conditions équivalentes :}

a) \textit{Pour tout} $y_0 \in \Ob\, Y$, \textit{le foncteur canonique}
\[
X_{y_0} \overset{\varphi}{\longrightarrow} X/y_0 \quad (\simeq X \times_Y Y/y_0) \ ,
\]
\textit{où} $X_{y_0}$ \textit{désigne la fibre en} $y_0$, \textit{est une} $\D$\textit{-équivalence co-stricte.}

b) \textit{Pour tout} $x_0 \in \Ob\, X$, \textit{le foncteur induit sur les catégories colocalisées en} $x_0$ \textit{et en} $y_0 = f(x_0)$
\[
x_0\backslash X \longrightarrow y_0\backslash Y
\]
\textit{est à fibres} $\D$\textit{-contractiles.}

\msp{page 69}

c) \textit{La formation des} $f_*F$ \textit{(}$F$ \textit{dans} $\D(X)$\textit{) commute à tout changement de base} $Y' \to Y$, \textit{et ça reste vrai après tout changement de base -- i.e. si on a} $Y'' \to Y' \to Y$, \textit{d'où}
\[
\begin{tikzcd}
X'' \arrow[r] \arrow[d, "f''"{swap}] & X' \arrow[r] \arrow[d, "f'"] & X \arrow[d, "f"] \\
Y'' \arrow[r, "p'"{swap}] & Y' \arrow[r, "p"{swap}] & Y
\end{tikzcd}
\]
\textit{le “théorème de changement de base” est vrai pour} $f'_*$ \textit{et} $p'^{*}$.

c$'$) \textit{La formation des} $f_*(F)$ \textit{(}$F \in \Ob\, \D(X)$\textit{) commute au passage aux fibres et ceci reste vrai après tout changement de base} $p : Y' \to Y$.

c$''$) \textit{Comme} c$'$) \textit{avec} $Y' \to Y$ \textit{un morphisme de localisation} $Y/y \to Y$.

d) \textit{Si} $Y'' \overset{p}{\to} Y' \to Y$, \textit{avec} $p$ \textit{coasphérique, alors} $q : X'' = X \times_Y Y'' \to X' = X \times_Y Y'$ \textit{est coasphérique.}

\printedpage{41}
% head: A. Grothendieck : Dérivateurs, Ch. I, Édition des Universités de Montpellier II et Paris VII
% foot: Version 8 juillet 2001

d$'$) \textit{Comme} d) \textit{avec} $Y'' = e$ (donc $Y'$ est catégorie avec objet final $e$).

d$''$) \textit{Comme} d) \textit{avec} $Y'' = \{1\}$, $Y' = \boldsymbol{\Delta}_1$.

.

\noindent\textsc{Démonstration}. a) $\Leftrightarrow$ b) cf. proposition page 53. Montrons que ces conditions impliquent c). Comme elles sont stables par changement de base (\textit{cf.} démonstration immédiate page 55) donc par \textbf{Der 2} on est ramené à prouver que $H^*_{\D}(X/y, F) \overset{\sim}{\to} H^*_{\D}(Y_y, F)$ $(F \in \Ob\, \D(Y))$, ce qui résulte du théorème 1.\footnote[30]{Donc on utilise \textbf{Der 5} (d).} (Cf. démonstration du théorème page 56, qui marche avec des hypothèses plus générales.)

On a c) $\Leftrightarrow$ c$'$) $\Rightarrow$ c$''$) (l'équivalence c) $\Leftrightarrow$ c$'$) par \textbf{Der 2}).

Prouvons c$''$) $\Rightarrow$ a).

\msp{page 70}

\noindent Considérons le diagramme à carrés cartésiens
\[
\begin{tikzcd}
X_{y_0} \arrow[r, "q"] \arrow[d, "p"{swap}] & X/y_0 \arrow[r] \arrow[d, "f_0"] & X \arrow[d, "f"] \\
\{y_0\} \arrow[r] & Y/y_0 \arrow[r] & Y
\end{tikzcd}
\qquad\qquad .
\]

\noindent L'hypothèse c$''$) prouve que pour tout $F$ dans $\D(X/y_0)$, on a $H^*_{\D}(X/y_0, F) \overset{\sim}{\to} H^*_{\D}(X_{y_0}, F)$. Par le théorème 1, ceci implique que $f$ est un $\D$-équivalence co-stricte. % [sic] un D-équivalence -> une D-équivalence, accord du genre

\dgrblock{dgrb40}{Completing the tests in the proper base-change theorem}
\begin{dcomment}{exposition}{p70:Par le théorème 1, ceci implique que}{Completing the tests in the proper base-change theorem}
The three conditions labelled (d) are not proved in the displayed argument. They have a short verification which also explains why a single walking arrow is a sufficient test. Write a pullback square as
\[
\begin{tikzcd}
 X''\arrow[r,"q"]\arrow[d,"f''"']&X'\arrow[d,"f'"]\\
 Y''\arrow[r,"p"']&Y'.
\end{tikzcd}
\]
If $f'$ is proper, the isomorphism $p^*f'_*\simeq f''_*q^*$ has, on passing to its left adjoint, the canonical isomorphism
\[
 q_!f''{}^*\simeq f'{}^*p_!.
\]
Here we take the left adjoints of the two \emph{composite functors}; no general assertion that every mate of every invertible cell is invertible is needed. If $p$ is coaspherical, its counit is an isomorphism on constants by Theorem~1. Applying the displayed comparison to a constant object on $Y''$ identifies the counit $q_!M_{X''}\to M_{X'}$ with its pullback. Theorem~1 now makes $q$ coaspherical. This proves (c)$\Rightarrow$(d).

Condition (d) implies its tests on inclusions of terminal objects, and these imply the test $\{1\}\to[1]$. Conversely, pull $f$ back along the functor $[1]\to Y$ determined by $\alpha:f(x)\to y$. Apply (d$''$) to the terminal-object inclusion $\{1\}\to[1]$. In the resulting coaspherical functor $X_y\to X\times_Y[1]$, the comma category under $(x,0)$ is exactly $L_f(x,\alpha)$. Thus (d$''$) implies (b). All the asserted conditions are now accounted for.

The parenthesis in (d$'$), ``donc $Y'$ est catégorie avec objet final $e$'', should be understood as specifying the terminal-object tests. A $\D$-coaspherical functor from the point need not literally select a terminal object for every possible coefficient theory. For example, take the represented derivator of the complete lattice $\{0<1\}$. Every nonempty discrete category is $\D$-contractible, since its diagonal constant-diagram functor is fully faithful. The inclusion $e\to BC_2$ is then $\D$-coaspherical: its only comma category is the discrete two-element set $C_2$. The unique object of $BC_2$ is not terminal. Reading (d$\prime$) as the restricted terminal-object test avoids attributing this false converse to the proof.

Several other slips here are already in the edition. The proof on page~69 requires $X_y$, not $Y_y$, and a coefficient on $X$, not on $Y$. At the end of the argument on page~70 the coaspherical functor is the displayed $q:X_{y_0}\to X/y_0$, not $f$. Definition~4 refers to Theorem~2, not Theorem~1. These corrections repair the typing and references; they do not alter the theorem.
\end{dcomment}


\noindent\textbf{Corollaire 1.} Énoncé dual du théorème, pour les foncteurs $X_{y_0} \to y_0\backslash X$, $X/x_0 \to Y/y_0$, et le changement de base des $j_!(F)$ \ldots

\noindent\textbf{Définition 4.} Sous les conditions équivalentes du théorème 1, on dit que $f$ est $\D$-\textit{propre.} La condition duale s'appelle $f$ $\D$-\textit{lisse.}

\noindent\textbf{Corollaire 2.} \textit{Soit} $f : X \to Y$ \textit{dans} $\Dia$. \textit{Si} $f$ \textit{est} $\D$\textit{-propre ou} $\D$\textit{-lisse, et à fibres} $\D$\textit{-contractiles, alors} $f$ \textit{est une} $\D$\textit{-équivalence stricte et costricte. Mieux,} $f$ \textit{est une} $\D$\textit{-équivalence} universelle \textit{i.e. est une} $\D$\textit{-équivalence et le reste après tout changement de base.}\footnote[31]{Se généralise : \textbf{Corollaire 2$'$.} \textit{Soit}
\[
\begin{tikzcd}[ampersand replacement=\&]
X \arrow[rr, "f"] \arrow[dr] \& \& Y \arrow[dl] \\
\& S \&
\end{tikzcd}
\]
\textit{avec} $X, Y$, $\D$\textit{-propres sur} $S$. \textit{Si} $f$ \textit{induit une} $\D$\textit{-équivalence sur les fibres c'est une} $\D$\textit{-équivalence. Itou dans le cas lisse.}\par
C'est la proposition 6 ci-dessous.}

\printedpage{42}
% head: A. Grothendieck : Dérivateurs, Ch. I, Édition des Universités de Montpellier II et Paris VII
% foot: Version 8 juillet 2001
\msp{page 71}

\noindent Il suffit de prouver que $f$ est une $\D$-équivalence. Supposons par exemple $f$ $\D$-propre. Il suffit de prouver que si $Y$ est $\D$-contractile, alors $X$ l'est (quitte à l'appliquer aux $Y/y_0$ après changement de base, pour montrer que les $X/y_0$ sont contractiles). Pour ceci, il suffit de prouver que $f$ est une $\D$-équivalence, et on sait que ceci est le cas si pour tout $F$ \textit{constant} sur $Y$, l'homomorphisme $F \overset{\text{adj}}{\longrightarrow} f_*f^*F$ est \texttt{[un]} isomorphisme. En fait, c'est vrai pour \textit{tout} $F$.

\noindent\textbf{Corollaire 3.} \textit{Si} $f : X \to Y$ \textit{est} $\D$\textit{-propre ou} $\D$\textit{-lisse, à fibres} $\D$\textit{-contractiles, alors} $f^* : \D(Y) \to \D(X)$ \textit{est pleinement fidèle, i.e. les homomorphismes d'adjonction}
\[
F \longrightarrow f_*f^*F \quad , \qquad f_!f^*F \longrightarrow F
\]
\textit{dans} $\D(Y)$ \textit{sont des isomorphismes.}

\noindent Il suffit de le voir pour $f$ propre, et cela résulte du calcul des fibres de $f_*f^*F$ \ldots

\noindent\textbf{Proposition 5.} \textit{La notion de foncteur} $\D$\textit{-propre} (resp. $\D$\textit{-lisse}) \textit{est stable par changement de base,} (trivial), \textit{et par composition} (utiliser \texttt{[le]} corollaire 2). \textit{Toute immersion fermée est} $\D$\textit{-propre, toute immersion ouverte} $\D$\textit{-lisse. Plus généralement, tout isomorphisme local est} $\D$\textit{-lisse, tout co-isomorphisme local est} $\D$\textit{-propre.}

\noindent\textbf{Proposition 6.} \textit{Soient} $X, Y$ \textit{sur} $S$, $f : X \to Y$ \textit{un} $S$\textit{-morphisme. Supposons} $X, Y$ \textit{propres sur} $S$, \textit{ou} $X$ \textit{et} $Y$ \textit{lisses sur} $S$, \textit{et que pour tout} $s \in \Ob\, S$, $f_s : X_s \to Y_s$ \textit{soit une} $\D$\textit{-équivalence, alors} $f$ \textit{aussi.}\footnote[32]{\textsc{Démonstration}. On prouvera même que sur la sous-catégorie des coefficients sur $Y$ qui proviennent de $\D(Y)$ \texttt{[plutôt de }$\D(S)$\texttt{]}, $f^* : \D(Y) \to \D(X)$ induit un foncteur pleinement fidèle. Si $f$ \texttt{[est]} propre, je l'interprète comme $q_*q^*F \overset{\sim}{\to} p_*p^*F$ pour tout $F$, et cela résulte du calcul de $p_*$, $q_*$ par passage aux fibres (propre), et de l'hypothèse sur les $f_s$. OK.}

\dgrblock{dgrb41}{Properness does not require constant fibres}
\begin{dcomment}{examples}{p71:et de l'hypothèse sur les}{Properness does not require constant fibres}
Let $A:[1]\to\Cat$ send $0$ to the point and $1$ to the discrete category $\{a,b\}$, with the transition selecting $a$. Its Grothendieck construction $X$ has three objects $u,a,b$, one nonidentity arrow $u\to a$, and projection $f:X\to[1]$ sending $u$ to $0$ and $a,b$ to $1$. The map $u\to a$ is a cocartesian lift of $0\to1$. The corresponding lift category is a point; the lift categories over identities have their identity lifts as initial objects. Thus $f$ is proper for every derivator under consideration.

For a contravariant coefficient diagram $F$ on $X$, the right direct image is
\[
 (f_*F)_0\simeq F_u,\qquad
 (f_*F)_1\simeq F_a\times F_b,
\]
and its transition is the composite $F_a\times F_b\to F_a\to F_u$. In the derivator of complexes, these are derived products, which for a finite pair agree with the ordinary finite product. For constant value $M$, the transition is the projection $M\times M\to M$. Taking $M=\mathbb Q$ in degree zero proves that it is not locally constant. The map is therefore proper but not a cohomological fibration in the definition about to be introduced.

This example separates two uses of a fibre calculation. Properness identifies a direct image with the cohomology of the fibres. A fibration condition additionally requires that these cohomologies be transported by isomorphisms along the base. There is no contradiction in a proper map having fibres with different homotopy types.

By contrast, if a proper map has $\D$-contractible fibres, then for every $F\in\D(Y)$ the component at $y$ of its unit is
\[
 F_y\longrightarrow H^*_{\D}(X_y,(F_y)_{X_y}),
\]
which is an isomorphism. Pointwise detection proves that $f^*$ is fully faithful. Since $f^*$ has both adjoints, this is equivalent both to invertibility of the unit of $f^*\dashv f_*$ and to invertibility of the counit of $f_!\dashv f^*$. This proves Corollary~3 and its smooth counterpart. It also gives universal $\D$-equivalence: properness and the fibres are preserved by base change. The claims of strict and costrict equivalence follow by applying that universal assertion to both kinds of local-isomorphism base changes.
\end{dcomment}


\msp{page 72}

\noindent\textbf{Théorème 3.} \textit{Soit}
\[
\begin{tikzcd}
X \arrow[d, "f"{swap}] & X' \arrow[l, "q"{swap}] \arrow[d, "f'"] \\
Y & Y' \arrow[l, "p"{swap}]
\end{tikzcd}
\]
\textit{cartésien dans} $\Dia$. \textit{Alors le morphisme de changement de base}
\[
p^*f_* \longrightarrow f'_*q^*
\]
\textit{est un isomorphisme si} $f$ \textit{est} $\D$\textit{-propre, ou} $p$ $\D$\textit{-lisse.}

\noindent Le cas $\D$-propre déjà prouvé. Le cas $\D$-lisse : la démonstration p. 56 est OK.

\dgrblock{dgrb42}{The geometric pattern and its categorical content}
\begin{dcomment}{context}{p72:la démonstration p. 56 est OK.}{The geometric pattern and its categorical content}
The vocabulary ``proper'' and ``smooth'' records a specific geometric analogy. In étale cohomology, proper direct image commutes with base change, and a smooth change of base gives a second base-change theorem under the appropriate torsion hypotheses. The comparison map has the same shape as the one displayed here. \citet[Introduction and Theorem~1.2]{Maltsiniotis2005Asphericite} explicitly traces the categorical definitions to this pattern in \emph{Pursuing Stacks} and identifies the corresponding SGA~4 exposés. His Theorems~4.13 and~5.4 develop the relation between lifting categories and base change in the framework of asphericity structures.

There is no implication here that a proper functor between categories is a proper morphism of spaces in the topological sense. The definition asks for a Beck--Chevalley isomorphism in a coefficient theory. The lifting categories supply its combinatorial content. The smooth half of Theorem~3 is obtained by applying that same calculation with the two adjunctions interchanged; it is not obtained merely by observing that the underlying square is a pullback. In the covariant-diagram convention, \citet[Propositions~1.24 and~1.26]{Groth2013} formulate the corresponding exactness of pullback squares for Grothendieck fibrations and opfibrations and of comma squares. Opposing all four indexing categories returns the convention of this manuscript.

For Proposition~6 one needs only to compare $q_*q^*M$ and $p_*p^*M$ over the common base $S$. Their components are constant-coefficient cohomologies of $Y_s$ and $X_s$. The hypothesis identifies those components, and \textbf{Der 2} finishes the proof. The footnote's stronger statement for coefficients pulled back from $S$ follows by the same calculation with $M$ replaced pointwise by $F_s$, followed by adjunction. This argument is valid without requiring the comparison map $X\to Y$ itself to be proper.
\end{dcomment}


\printedpage{43}
% head: A. Grothendieck : Dérivateurs, Ch. I, Édition des Universités de Montpellier II et Paris VII
% foot: Version 8 juillet 2001

\subsection*{Comparaison de $X$ et $X^{\circ}$.}

Soit $X$ \texttt{[une]} catégorie, on définit une catégorie $\Phi(X)$ et des foncteurs
\[
\begin{tikzcd}[column sep=small]
& & \Phi(X) \arrow[ddl, "p_X"{swap, name=P}] \arrow[ddr, "q_X"{name=Q}] & & \\
|[alias=L]| \text{\footnotesize\begin{tabular}{c}$\Cat$-cofibrant\\(donc propre)\end{tabular}} & & & & |[alias=R]| \text{\footnotesize\begin{tabular}{c}$\Cat$-fibrant\\(donc lisse)\\ \texttt{[plutôt }$\Cat$\texttt{-cofibrant}\\ \texttt{(donc propre)]}\end{tabular}} \\
& X & & X^{\circ} &
\arrow[from=P, to=L]
\arrow[from=Q, to=R]
\end{tikzcd}
\]
qui seront des $\hot$-équivalences, et des $\D$-équivalences si $\D$ est un dérivateur dont le domaine contient $X$ et $\Phi(X)$. On pose
\[
\Ob\, \Phi(X) = \Fl(X) \quad ,
\]
donc mêmes objets que $\mathcal{F}\ell(X) = \underline{\Hom}(\boldsymbol{\Delta}_1, X)$, mais les morphismes sont différents :

\msp{page 73}
\[
\begin{tikzcd}
Y_1 \arrow[r, "u"] & Y_0 \arrow[d, "v_0"] \\
Y'_1 \arrow[u, "v_1"] \arrow[r, "u'"{swap}] & Y'_0
\end{tikzcd}
\]
un morphisme de $u$ dans $u'$ est un couple $v_0 : Y_0 \to Y'_0$, $v_1 : Y'_1 \to Y_1$ tel que $v_0uv_1 = u'$. Ainsi source et but $s(u)$ et $b(u)$ sont l'une contravariante, l'autre covariant en $u$. On prend pour $q$ le foncteur source, pour $p$ le foncteur but. % [sic] l'autre covariant -> l'autre covariante, accord avec « l'une contravariante »

Si on applique ceci à $X^{\circ}$, on trouve
\[
\begin{tikzcd}
& \Phi(X^{\circ}) \arrow[dl, "p_{X^{\circ}}"{swap}] \arrow[dr, "q_{X^{\circ}}"] & \\
X^{\circ} & & (X^{\circ})^{\circ} = X
\end{tikzcd}
\]
On a $\Ob\, \Phi(X^{\circ}) = \Fl(X^{\circ}) \simeq \Fl(X)$
\[
\begin{tikzcd}
Y_1 \arrow[r, "u"] & Y_0 \arrow[d, "v_1"] & \text{dans}\ \ X \\
Y'_1 \arrow[u, "v_0"] \arrow[r, "u'"{swap}] & Y'_0 & \text{dans}\ \ X
\end{tikzcd}
\]
et mêmes homomorphismes au sens de $\Phi(X)$ et au sens de $\Phi(X^{\circ})$. Donc
\[
\left\{
\begin{aligned}
&\Phi(X^{\circ}) = \Phi(X) \\
&p_{X^{\circ}} = q_X \\
&q_{X^{\circ}} = p_X \quad ,
\end{aligned}
\right.
\]

% ---- from der1_p046-050.tex
\printedpage{44}
% head: A. Grothendieck : Dérivateurs, Ch. I, Édition des Universités de Montpellier II et Paris VII
le passage de $X$ à $X^\circ$ revient à échanger $p_X$ et $q_X$. La fibre de $\Phi(X)$ en $x \in \Ob X$ est $\simeq (X/x)^\circ$, celle en $y^\circ \in \Ob X^\circ$ est $\simeq y\bs X$
\[
\left\{
\begin{aligned}
&p_X^{-1}(\{x\}) \simeq (X/x)^\circ\\
&q_X^{-1}(\{y^\circ\}) \simeq y\bs X
\end{aligned}
\right.
\]

{\raggedright
N.B. : $p_X$ est $W$-lisse \texttt{[plutôt $W$-propre]} à fibres $W$-asphériques, donc dans $W^u$ \texttt{[$W^u =$ ensemble des morphismes qui sont dans $W$ et le restent après tout changement de base]}, $q_X$ est $W$-propre à fibres $W$-asphériques, donc dans $W^u$.\par}

\dgrblock{dgrb43}{The twisted-arrow category and its two projections}
\begin{dcomment}{exposition}{p73:et le restent après tout changement de base}{The twisted-arrow category and its two projections}
The category $\Phi(X)$ is the twisted-arrow category, with the orientation explicitly fixed by the equation $u'=v_0uv_1$. Keeping that equation in view avoids confusing it with the ordinary arrow category. An object $u:a\to b$ maps to $b$ under $p_X$ and to $a$, considered in $X^{\mathrm{op}}$, under $q_X$. Composition is componentwise in the appropriate orders. Replacing $u$ by $u^{\mathrm{op}}$ interchanges the two components of a twisted morphism and gives a canonical isomorphism
\[
 \Phi(X^{\mathrm{op}})\cong\Phi(X),
\]
not the opposite category on the right. The two projection formulas follow under this identification.

Both projections are opfibrations. For $p_X$, an arrow $\beta:b\to b'$ is lifted from $u$ by the twisted arrow $u\to\beta u$, with source component $1_a$ and target component $\beta$. If a twisted arrow from $u$ to $w:c\to d$ has target component $\gamma\beta$, it factors uniquely through this chosen lift, using target component $\gamma$ and the same source component. This is the cocartesian universal property. For $q_X$, an arrow $a\to a'$ in $X^{\mathrm{op}}$ is an arrow $\alpha:a'\to a$ in $X$; its cocartesian lift is $u\to u\alpha$. The same factorisation argument applies.

The fibre of $p_X$ over $b$ is $(X/b)^{\mathrm{op}}$: its arrows are the reversed triangles between maps to $b$. Its object $1_b$ is initial. The fibre of $q_X$ over $a$ is $a\backslash X$, whose object $1_a$ is initial as well. Consequently both maps are proper with contractible fibres. The preceding corollaries prove directly that they are universal $\D$-equivalences and that inverse image along either is fully faithful. This gives a shorter proof of the assertion announced on page~74; it uses exactly the corrected ``cofibrant, donc propre'' label in the edition.

The displayed slice identities are also literal categorical isomorphisms, and not merely equivalences of homotopy types. An object of $\Phi(X)/x$ is a composable pair $a\xrightarrow{u}b\xrightarrow{h}x$. It determines the arrow $(a,hu)\to(b,h)$ in $X/x$, hence an object of $\Phi(X/x)$. A twisted morphism consists of $v_1:a'\to a$, $v_0:b\to b'$ with $u'=v_0uv_1$ and $h'v_0=h$. The remaining slice equation $h'u'=huv_1$ then follows. These are precisely the morphisms in $\Phi(X/x)$. Reversing the construction is immediate. The second identity follows by applying this argument to $X^{\mathrm{op}}$ and interchanging the endpoints.

The last comma notation in the succeeding proof requires a correction. The strict-equivalence test for $q_Z:\Phi(Z)\to Z^{\mathrm{op}}$ uses $\Phi(Z)/z_0^{\mathrm{op}}$, whose additional arrow in $Z$ is $z_0\to s(u)$. It is this category that is isomorphic to $\Phi(z_0\backslash Z)$. The edition instead prints $z_0\backslash\Phi(Z)$; interpreted literally as a comma category under $q_Z$, that symbol would require $s(u)\to z_0$ in $Z$ and reverse the intended arrow. With the intended comma category restored, the proof continues correctly. If $Z$ has a terminal object, $z_0\backslash Z$ has both an initial and a terminal object. The unique arrow between these is terminal in its twisted-arrow category, exactly as Grothendieck observes.
\end{dcomment}


\msp{page 74}

On vérifie
\[
\left\{
\begin{aligned}
&\Phi(X)/x \simeq \Phi(X/x) &\qquad& x \in \Ob X\\
&\Phi(X)/x^\circ \simeq \Phi(x\bs X) &\qquad& x^\circ \in \Ob X^\circ
\end{aligned}
\right.
\]

Je dis que $p_X$, $q_X$ sont des $\D$-\textit{équivalences strictes} et \textit{costrictes.} Par les formules précédentes, cela signifie que les $\Phi(X/x)$, $\Phi(x\bs X)$ sont $\D$-contractiles. Donc que si $Z$ a un objet final, $\Phi(Z)$ est $\D$-contractile. Mais considérons
\[
\begin{tikzcd}
\Phi(Z) \arrow[dr, "q_Z"] & \\
& Z^\circ
\end{tikzcd}
\quad ;
\]
il suffit de prouver que c'est une $\D$-équivalence stricte (puisque $Z^\circ$, ayant un objet initial, est $\D$-contractil), donc que les $z_0\bs\Phi(Z)$ sont $\D$-contractiles. On a par les formules précédentes % [sic] contractil -> contractile, adjectif ailleurs écrit « contractile »
\[
z_0\bs\Phi(Z) \simeq \Phi(z_0\bs Z) = \Phi(T) \quad ,
\]
où cette fois $T$ a un objet initial $\emptyset$ et un objet final $e$. Considérons alors l'objet
\[
\emptyset \longrightarrow e
\]
de $\Phi(T)$, c'est un objet final, car pour tout $x \to y$ dans $\Phi(T)$, il y a une unique flèche
\[
\begin{tikzcd}
x \arrow[r] & y \arrow[d] \\
\emptyset \arrow[u] \arrow[r] & e
\end{tikzcd}
\]
% foot: Version 8 juillet 2001

\dgrblock{dgrb44}{Intervals in a finite ordered set}
\begin{dcomment}{examples}{p74:il y a une unique flèche}{Intervals in a finite ordered set}
For $X=[1]$, the twisted-arrow category has objects $1_0$, $u:0\to1$, and $1_1$, with precisely the nonidentity arrows
\[
 1_0\longrightarrow u\longleftarrow1_1.
\]
Thus $u$ is terminal. Under the target projection, the fibre at $0$ is the point $1_0$, whereas the fibre at $1$ is the arrow $1_1\to u$. Under the source projection, the roles are reversed. Each fibre is contractible, although the projections are not isomorphisms of categories. Their nerves subdivide the interval into two edges.

For the chain $[n]$, objects are intervals $(i,j)$ with $0\leq i\leq j\leq n$. There is exactly one morphism
\[
 (i,j)\longrightarrow(i',j')
 \quad\Longleftrightarrow\quad i'\leq i\leq j\leq j'.
\]
The largest interval $(0,n)$ is terminal. More generally this description works for comparable pairs in any poset, showing immediately that $\Phi$ preserves finite posets. This is the domain-closure observation at the end of manuscript page~75; it is not an assertion that every finite category is a finite poset.

The comparison with the opposite category is useful because it is natural. For a functor $f:X\to Y$, applying $f$ to arrows and to both components gives $\Phi(f)$, and both endpoint squares commute. Since their vertical projections are $\D$-equivalences, two-out-of-three first gives
\[
 f\in W_{\D}\ \Longleftrightarrow\ \Phi(f)\in W_{\D},
\]
and then the corresponding equivalence with $f^{\mathrm{op}}\in W_{\D}$. One must have $X^{\mathrm{op}},Y^{\mathrm{op}},\Phi(X),\Phi(Y)$ in the domain. There is no claim that $X$ and $X^{\mathrm{op}}$ are equivalent as ordinary categories. The natural span compares precisely the information detected by the chosen coefficient theory.
\end{dcomment}


\printedpage{45}
% head: A. Grothendieck : Dérivateurs, Ch. I, Édition des Universités de Montpellier II et Paris VII
\msp{page 75}

Pour résumer

\textbf{Proposition 1 :} \textit{$p_X$ et $q_X$ sont des} \textsf{hot}\textit{-équivalences strictes,} i.e. \textsf{Hot}\textit{-asphériques (donc aussi des $\D$-équivalences strictes si $\Phi(X)$ est dans le domaine de $\D$) à fibres les $(X/x)^\circ$ et $y\bs X$ respectivement (donc à fibres contractiles).}

\textbf{Corollaire 1.} \textit{Pour que $X$ soit $\D$-contractile, il faut et il suffit que $X^\circ$ le soit.}

\hspace*{\parindent}Plus généralement :

\textbf{Corollaire 2.} \textit{Soit $f : X \to Y$. Pour que $f$ soit une $\D$-équivalence, il faut et il suffit que $f^\circ$ le soit.}

\hspace*{\parindent}On utilisera que le diagramme
\[
\begin{tikzcd}
& \Phi(X) \arrow[dl, "p_X"'] \arrow[dr, "q_X"] & \\
X & & X^\circ
\end{tikzcd}
\]
est fonctoriel en $X$. On a donc
\[
\begin{tikzcd}
& \Phi(X) \arrow[rr, "\Phi(f)"] \arrow[dl, "p_X"'] & & \Phi(Y) \arrow[dl, "p_Y"] \arrow[dd, "q_Y"] \\
X \arrow[rr, "f" near end] \arrow[from=1-2, to=3-2, "q_X"' near end, crossing over] & & Y & \\
& X^\circ \arrow[rr, "f^\circ"'] & & Y^\circ
\end{tikzcd}
\]
les deux carrés commutatifs. Comme $p_X$ et $p_Y$ sont des $\D$-équivalences, $f$ est une $\D$-équivalence si et seulement si $\Phi(f)$ l'est, et de même pour $f^\circ$. Ces raisonnements supposent que si $X$ est dans le domaine de $\D$, $\Phi(X)$ aussi. C'est le cas si \Dia{} est formé des catégories ordonnées finies.
% foot: Version 8 juillet 2001

\dgrblock{dgrb45}{Twisted arrows after the opposite-category comparison}
\begin{dcomment}{later}{p75:Ces raisonnements supposent que si}{Twisted arrows after the opposite-category comparison}
The construction has become a standard indexing device for homotopy ends and coends. \citet[Example~7.2.9 and the footnote in Section~7.8]{Riehl2014} describe the twisted-arrow category and the identification of its nerve with edgewise subdivision. This supplies a simplicial interpretation of the span $X\leftarrow\Phi(X)\to X^{\mathrm{op}}$. Its realization compares the two classifying spaces without choosing an equivalence of categories.

For a derivator written with the covariant convention, \citet[Section~5]{PontoShulman2016} use the same category to define the homotopy coend of $A\in\mathbb V(X^{\mathrm{op}}\times X)$ by
\[
 \int^{X}A=(\pi_{\Phi(X)^{\mathrm{op}}})_!
       (t^{\mathrm{op}},s^{\mathrm{op}})^*A.
\]
The opposite on the indexing category is essential in this formula. Ends use the other variance. The passage from this formula to composition of homotopy profunctors explains why the two endpoint projections occur together so persistently in later derivator calculus. Grothendieck uses them here for a narrower purpose: to show that the weak-equivalence class obtained from constant coefficients is unchanged by taking opposites. The later coend formalism extends the utility of the construction; it is not an extra assumption in his proof.
\end{dcomment}


\printedpage{46}
% head: A. Grothendieck : Dérivateurs, Ch. I, Édition des Universités de Montpellier II et Paris VII
\msp{page 76}

\subsection*{$\D$-fibrations}

\textbf{Proposition 1.} \textit{Soit $f : X \to Y$ dans} \Dia\textit{, conditions équivalentes :}

\hspace*{\parindent}a) \textit{Pour tout $y \in Y$, l'inclusion $X_y \to X/y$ est une $\D$-équivalence, et pour toute flèche $y_1 \to y_0$ de $Y$, $X/y_1 \to X/y_0$ est une $\D$-équivalence.}

\hspace*{\parindent}b) \textit{Pour tout coefficient} constant \textit{$F$ sur $Y$, $f_*f^*F$ est localement constant,}\footnote[33]{Par là on entend que pour tout $u : y \to y'$ dans $Y$, $G_{y'} \to G_y$ est \texttt{[un]} isomorphisme. Sans axiomes supplémentaires sur $\D$, il n'est pas clair que cela implique que les ``restrictions'' de $G$ aux $Y/y$ ou $y\bs Y$, soient isomorphes à des coefficients constants. Voir cependant l'axiome \textbf{Der 3} a$'$), page 97, qui assure que c'est pareil.} \textit{et ``se calcule fibre par fibre''} i.e. \textit{l'homomorphisme canonique}
\[
f_*f^*(F)_y \overset{\sim}{\longrightarrow} H^*_{\D}(X_y, F|X_y)
\]
\texttt{[est un isomorphisme]}.

Ces conditions sont équivalentes, et impliquent les conditions de (b) pour tout $F$ \textit{localement} constant sur $Y$.

La première condition dans (a) signifie que pour tout $y \in Y$, et tout $M \in \Ob \mathcal{A}$,
\[
H^*_{\D}(X/y, M) \longrightarrow H^*_{\D}(X_y, M)
\]
est un isomorphisme. Or cette flèche s'interprète comme
\[
f_*(M_X)_y \longrightarrow H^*_{\D}(X_y, M)
\]
est un isomorphisme, \textit{i.e.} que $f_*(M_X)$ se calcule fibre par fibre. Comme $M_X \simeq f^*(M_Y)$, c'est la deuxième condition (b). La deuxième condition dans (a), pour $y_1 \to y_0$ signifie que pour $M$ dans $\mathcal{A}$,
\[
H^*_{\D}(X/y_0, M) \longrightarrow H^*_{\D}(X/y_1, M)
\]
est \texttt{[un]} isomorphisme, or cette flèche s'interprète comme
\[
f_*(M_X)_{y_0} \longrightarrow f_*(M_X)_{y_1} \quad ,
\]
que ce soit toujours un isomorphisme, c'est la première condition (b). Que (b) soit valable aussi pour $F$ localement constant sur $Y$ est immédiat, car c'est une condition de nature locale par rapport à $Y$.

\dgrblock{dgrb46}{Local constancy and the first fibration condition}
\begin{dcomment}{exposition}{p76:que ce soit toujours un isomorphisme,}{Local constancy and the first fibration condition}
The new fibration condition retains two separate tests. First, the canonical inclusion of an actual fibre into its comma category must be a $\D$-equivalence. Second, the transition functors between comma categories must be $\D$-equivalences. The first identifies the values of $f_*M_X$ with fibre cohomology. The second says that their transport maps are isomorphisms. Since $\D$-equivalence tests all absolute coefficients $M$, this is an equivalence of conditions, not an inference from the cohomology of a single distinguished coefficient.

The extension from constant to locally constant coefficients is valid with the usual two-dimensional functoriality and pointwise detection axioms. There is a direct proof of the local issue raised in footnote~33. If a category $A$ has terminal object $a$, the natural transformation $1_A\to ap$ induces, in the manuscript's presheaf convention, a morphism
\[
 p^*a^*F\longrightarrow F.
\]
For arrowwise locally constant $F$, every component is an isomorphism; \textbf{Der 2} makes the whole morphism invertible. If $A$ has an initial object, the corresponding natural transformation gives the comparison in the other direction, with the same conclusion. Applying these observations to $Y/y$ and $y\backslash Y$ proves that restrictions of a locally constant object are constant there. This only uses categories possessing the indicated distinguished object. The stronger assertion for \emph{every} weakly contractible category is the separate issue addressed by the later variants of \textbf{Der 3}.

Now fix an arrow $y_1\to y_0$. Pull a locally constant $F$ back to $Y/y_0$, where it is isomorphic to the constant object with value $F_{y_0}$. Pulling this isomorphism further to $X/y_0$, $X/y_1$, and the two fibres turns the desired comparison into the one already established for constant coefficients. The isomorphism $F_{y_0}\to F_{y_1}$ supplies the change of value at the other endpoint. This proves the statement for locally constant $F$, including the compatibility of transport, rather than relying only on the phrase ``de nature locale''.

For model-category derivators, a systematic treatment is \citet[Theorem~1.3 and Section~3]{Cisinski2009LocallyConstant}: weak equivalences of indexing nerves induce equivalences on the homotopy categories of locally constant diagrams. The distinction between constant and locally constant objects remains necessary when the base has nontrivial fundamental group. It is exactly this distinction that will invalidate the proposed implication on manuscript pages~80--81.
\end{dcomment}


\msp{page 77}

La proposition duale fait appel aux foncteurs
\[
y_0\bs X \longrightarrow y_1\bs X \qquad \text{et} \qquad X_y \longrightarrow y\bs X \quad ,
\]
et aux faisceaux $f_!f^*M$, pour $M$ constant, ou localement constant sur $Y$.
% foot: Version 8 juillet 2001

\printedpage{47}
% head: A. Grothendieck : Dérivateurs, Ch. I, Édition des Universités de Montpellier II et Paris VII
\textbf{Définition 1.} Si les conditions équivalentes de la proposition 1 sont satisfaites, on dit que $f$ est une $\D$-\textit{fibration cohomologique} (ou $\D$-\textit{fibration gauche}). Si les conditions duales sont satisfaites, on dit que c'est une $\D$-\textit{fibration homologique} (ou $\D$-\textit{fibration droite}).

N.B. Les $\D$-fibrations cohomologiques sont stables par tout changement de base qui est un isomorphisme local. Les $\D$-fibrations homologiques par tout changement de base qui est un co-isomorphisme local. Mais ce ne sont pas des notions stables par \textit{tout} changement de base\footnote[34]{? donner exemple!}. Mais la notion de $\D$-fibration cohomologique est stable par changement de base $\D$-lisse, et celle de $\D$-fibration homologique par changement de base $\D$-propre.

\dgrblock{dgrb47}{A two-object failure of arbitrary base change}
\begin{dcomment}{development}{p77:Mais ce ne sont pas des notions stables par}{A two-object failure of arbitrary base change}
The footnote asks for an example showing that the first fibration condition need not survive arbitrary base change. We can supply one before the more elaborate counterexamples later in the manuscript. We checked the treatments of smoothness and asphericity in \citet[Sections~1 and~4]{Maltsiniotis2005Asphericite} and of locally constant diagrams in \citet{Cisinski2009LocallyConstant}; we did not find this precise small example there. The construction below is an elementary verification for this commentary, without a claim of priority.

\textbf{A non-universal cohomological fibration.}\ours
Let $J$ be the groupoid with two objects $a,b$ and exactly one arrow between every ordered pair of objects. Let $f:[1]\to J$ send $0$ to $a$, $1$ to $b$, and the nonidentity arrow to the unique isomorphism $a\to b$. For the space derivator, $f$ is both a cohomological and a homological fibration in Definition~1, but one strict base change of $f$ has neither property.

\emph{Proof.} Each strict fibre of $f$ is a point. For either $y\in J$, an object of $[1]/y$ consists of $i\in[1]$ and an arrow $f(i)\to y$. The latter is unique. Projection therefore gives an isomorphism $[1]/y\cong[1]$. Under these isomorphisms every transition functor is the identity. The inclusion of either fibre is an endpoint inclusion $e\to[1]$, hence a weak equivalence of nerves. This proves the cohomological condition. The same argument with $y\backslash[1]$ proves the homological condition. More generally these tests succeed for any coefficient theory in which the point and the walking arrow have the usual constant-coefficient invariance.

Define $g:[1]\to J$ in the reverse order: $g(0)=b$, $g(1)=a$. Form the strict pullback
\[
 P=[1]\times_J[1]\xrightarrow{\,f'\,}[1],
\]
with $f$ on the first factor and $g$ on the second. Its objects are exactly $(1,0)$ and $(0,1)$. There is no arrow from the first to the second, because there is no arrow $1\to0$ in the first factor. There is none in the reverse direction, because there is no arrow $1\to0$ in the second factor. Thus $P$ is the discrete two-object category, with one object in each fibre of $f'$.

The comma category $P/1$ is discrete with two objects, whereas $P_1$ is a point. Hence $P_1\to P/1$ is not a weak equivalence, and $f'$ is not a cohomological fibration. Dually $0\backslash P$ is discrete with two objects while $P_0$ is a point, so the homological condition fails too. For complexes of rational vector spaces the same failure is visible in degree zero: the fibre comparison is $\mathbb Q\times\mathbb Q\to\mathbb Q$, a projection with nonzero kernel.\hfill$\square$

Both $[1]$ and $J$ in the original pair of maps are contractible, and every original strict fibre is a point. The defect is the incompatible choice of direction for a strict lift of the base isomorphism. This example therefore isolates the need for a universal condition without invoking monodromy or a failure of constant coefficients to detect it. The later counterexamples address a different, stronger assertion: whether a universal fibration preserves all $\D$-equivalences under pullback.
\end{dcomment}


\hspace*{\parindent}Je vais essayer de construire un exemple où le changement de base canule, en posant $Y = \boldsymbol{\Delta}_1$. Se donner $X$ sur $\boldsymbol{\Delta}_1$ revient à se donner une sous-catégorie ouverte $X_0$, et son complémentaire fermé $X_1$
\[
\begin{tikzcd}
& X & \\
X_1 \arrow[ur, hook] \arrow[d, dash, dashed] & & X_0 \arrow[ul, hook'] \arrow[d, dash, dashed] \\
1 & & 0 \arrow[ll]
\end{tikzcd}
\quad .
\]
Dire que $X \to Y = \boldsymbol{\Delta}_1$ est une $\D$-fibration cohomologique signifie simplement
\msp{page 78}
que les restrictions
\[
X_0 \overset{j}{\longrightarrow} X \quad , \qquad X_1 \overset{i}{\longrightarrow} X
\]
sont des $\D$-équivalences. Prenons un changement de base $Y' \to Y = \boldsymbol{\Delta}_1$, \textit{i.e.} un $Y'$ avec $Y'_0, Y'_1$
\[
\begin{tikzcd}
& Y' & \\
Y'_1 \arrow[ur, hook] \arrow[d, dash, dashed] & & Y'_0 \arrow[ul, hook'] \arrow[d, dash, dashed] \\
1 & & 0 \arrow[ll]
\end{tikzcd}
\quad .
\]
On voit que la condition $Y'$ lisse sur $\boldsymbol{\Delta}_1$ signifie que pour tout $y'_1 \in \Ob Y'_1$, la catégorie $Y'_0/y'_1$ des objets $y_0$ de $Y'_0$ au dessus de $y'_1$ (\textit{i.e.} munis de $y_0 \to y'_1$) est contractile. On voit que si $Y'_1 \neq \emptyset$ alors $Y'_0 \neq \emptyset$ -- mais un exemple avec $Y'_0 = \emptyset$, où $Y'_1 = \emptyset$, n'est pas intéressant
% foot: Version 8 juillet 2001
\printedpage{48}
% head: A. Grothendieck : Dérivateurs, Ch. I, Édition des Universités de Montpellier II et Paris VII
(car une projection $I \times J \to I$ est toujours une $\D$-fibration cohomologique). Prenons par exemple pour $Y'$
\[
\begin{tikzcd}
& a \arrow[dl] \\
b & \\
& a' \arrow[ul] \\
\underbrace{\qquad}_{\displaystyle Y'_1} & \underbrace{\qquad}_{\displaystyle Y'_0}
\end{tikzcd}
\quad ,
\]
somme amalgamée de deux copies de $Y = \boldsymbol{\Delta}_1$, donc $X' = X \times_Y Y'$ peut se voir comme somme amalgamée $X \amalg_{X_1} X$. Mais ça semble \textit{pas} être un contre-exemple. Il faudrait que $X_1 \to X'$ ne soit \textit{pas} une $\D$-équivalence.

\dgrblock{dgrb48}{Why the first doubled-arrow attempt does not work for spaces}
\begin{dcomment}{examples}{p78:somme amalgamée de deux copies de}{Why the first doubled-arrow attempt does not work for spaces}
The fork $a\to b\leftarrow a'$ is a useful test even though it does not provide the desired counterexample for the space derivator. A functor $X\to[1]$ decomposes $X$ into a sieve $X_0$ and a cosieve $X_1$. Pullback along that fork is two copies of $X$ identified along $X_1$:
\[
 X'=X\amalg_{X_1}X.
\]
Every nonidentity chain in $X'$ lies in one of the copies. Indeed, after it enters the common cosieve $X_1$, it cannot leave it to reach the other copy of $X_0$. Consequently the same pushout description holds degree by degree for nerves:
\[
 NX'=NX\amalg_{NX_1}NX.
\]
The maps $NX_1\to NX$ are inclusions of simplicial subsets. Their realizations are inclusions of CW subcomplexes. If $X\to[1]$ is a cohomological fibration for spaces, both $X_0\to X$ and $X_1\to X$ are weak homotopy equivalences. In particular the latter inclusion is a homotopy equivalence of CW complexes and a cofibration. Gluing a second copy along it preserves its homotopy type: the homotopy extension property, applied to an inverse and its homotopy, makes the inclusion of either copy in the pushout a homotopy equivalence. It follows that $|NX_1|\to|NX'|$ is a weak equivalence.

At $a$ and $a'$ the new fibre and comma category are copies of $X_0$. At $b$ they are $X_1$ and $X'$; the preceding paragraph supplies their comparison. The transport from either of the first two comma categories to the last is a copy of $X_0\to X\to X'$, hence a weak equivalence. Thus this particular base change is still a cohomological fibration. This establishes Grothendieck's expressed suspicion in the space case; it does not assert that arbitrary base changes work.
\end{dcomment}


\msp{page 79}

\textbf{Proposition 2.} \textit{Soit $f : X \to Y$ dans} \Dia. \textit{Conditions équivalentes :}

\hspace*{\parindent}a) \textit{Pour tout $Y'' \overset{p}{\to} Y' \to Y$, d'où}
\[
\begin{tikzcd}
X'' \arrow[r, "q"] \arrow[d] \arrow[dr, phantom, "\text{cart}"] & X' \arrow[r] \arrow[d] \arrow[dr, phantom, "\text{cart}"] & X \arrow[d, "f"] \\
Y'' \arrow[r, "p"'] & Y' \arrow[r] & Y
\end{tikzcd}
\quad ,
\]
\textit{si $p$ est une $\D$-équivalence,} \texttt{[alors]} \textit{$q$ aussi.}\footnote[35]{Il faut supposer $Y', Y''$ \textsf{Hot}-asphériques, cf. page 96, et supposer vérifié l'axiome Der 3 $a'$) de la page 97.}

\hspace*{\parindent}b) \textit{$f$ est une $\D$-fibration cohomologique ``universelle'',} i.e. \textit{est une fibration cohomologique, et le reste après tout changement de base.}

\hspace*{\parindent}b$'$) \textit{Condition duale de} (b) \textit{: $f$ est une $\D$-fibration homologique universelle.}\footnote[36]{b) $\Rightarrow$ a) \underline{canulée}, \textit{cf.} annotation marginale,\footnotemark[37] \textit{cf.} contre-exemple page 89. Mais ça marche pour les dérivateurs $\D_i$ ($i \in \mathbb{N}$) associés aux ``localisateurs fondamentaux'' $W_i$, \textit{i.e.} à la théorie des $i$-types d'homotopie.}\footnotetext[37]{L'implication b) $\Rightarrow$ a) pas prouvée, seulement sous forme plus faible :\par \hspace*{\parindent}a$'$) si $Y''$, $Y'$ sont $\D$-asphériques et 1-connexes, alors $X'' \to X'$ est une $\D$-équivalence. (Il suffit même de le supposer pour $Y'$, $Y''$ ayant des objets finaux (pour impliquer b) ou ayant des objets initiaux (pour impliquer b$'$).)}

\textit{a}) $\Rightarrow$ \textit{b}) : Supposons \textit{a}), soit $Y' \to Y$ changement de base, d'où
\[
\begin{tikzcd}
X' \arrow[r] \arrow[d, "f'"'] & X \arrow[d, "f"] \\
Y' \arrow[r] & Y
\end{tikzcd}
\quad .
\]
% foot: Version 8 juillet 2001

% ---- from Der-ChI-234_p051-055.tex
\printedpage{49}
% head: A. Grothendieck : Dérivateurs, Ch. I, Édition des Universités de Montpellier II et Paris VII
% foot: Version 8 juillet 2001

Pour prouver que $f'$ est une $\DD$-fibration cohomologique, on regarde au dessus de $Y'$ les foncteurs de la forme $\{y'\} \to Y'/y'$ ($y' \in \Ob\, Y'$) et $Y'/y'_1 \to Y'/y'_0$ (si $u' : y'_1 \to y'_0$ flèche de $Y'$), et on voit que les flèches dans $\Cat/Y'$ déduites de ceux-ci par le changement de base $X' \xrightarrow{f'} Y'$ restent des $\DD$-équivalences

\msp{page 80}

\[
\begin{tikzcd}[column sep=large, row sep=large]
X'_1 \arrow[r, "q"] \arrow[d, "f'_1"'] \arrow[dr, phantom, "\text{cart}"] & X'_0 \arrow[r] \arrow[d, "f'_0"] \arrow[dr, phantom, "\text{cart}"] & X' \arrow[r] \arrow[d, "f'"] \arrow[dr, phantom, "\text{cart}"] & X \arrow[d, "f"] \\
Y'_1 \arrow[r, "p"'] & Y'_0 \arrow[r] & Y' \arrow[r] & Y & .
\end{tikzcd}
\]
Comme $p$ est une $\DD$-équivalence, cela résulte bien de a).

\indent\textit{b}) $\Rightarrow$ \textit{a}) : Comme la condition \textit{b}) est évidemment stable par changement de base, on est ramené à prouver que si $f$ satisfait \textit{b}), et si on a
\[
\begin{tikzcd}[column sep=large, row sep=large]
X \arrow[d, "f"'] \arrow[dr, phantom, "\text{cart}"] & X' \arrow[l, "q"'] \arrow[d, "f'"] \\
Y & Y' \arrow[l, "p"]
\end{tikzcd}
\tag*{$(*)$}
\]
si $p$ est une $\DD$-équivalence, alors $q$ aussi, \textit{i.e.}
\[
\begin{tikzcd}[row sep=small]
H^*_{\DD}(X,M) \arrow[r, "\sim"] & H^*_{\DD}(X',M) \\
\wr| & \wr| \\
H^*_{\DD}(Y, f_*(M_X)) & H^*_{\DD}(Y', f'_*(M_{X'}))
\end{tikzcd}
\]
pour tout $M$ dans $\mathcal{A} = \DD(e)$. Mais $F = f_*(M_X)$ est localement constant et se calcule “fibre par fibre”, de même pour $F' = f'_*(M_{X'})$.

\msp{page 81}

L'équivalence \textit{a}) $\Leftrightarrow$ \textit{b}$'$) se déduit de \textit{a}) $\Leftrightarrow$ \textit{b}) par dualité (la condition \textit{a}) étant autoduale), ou en reprenant le même raisonnement pour \textit{a}) $\Rightarrow$ \textit{b}$'$), et pour \textit{b}$'$) $\Rightarrow$ \textit{a}) aussi (en remplaçant la cohomologie par l'homologie).

\dgrblock{dgrb49}{The coefficient lost in the attempted proof}
\begin{dcomment}{exposition}{p81:en remplaçant la cohomologie par l'homologie}{The coefficient lost in the attempted proof}
The proposition is an experiment whose unsuccessful part Grothendieck subsequently identifies himself. Its three conditions must therefore be read together with the marginal corrections. Condition (a) says that pullback along every base change of $f$ preserves the entire class of $\DD$-equivalences. Conditions (b) and (b$'$) say that constant coefficients have universally well-behaved cohomological or homological direct images. The first condition implies the other two. The converse requires information about coefficients which the first definition of a $\DD$-equivalence does not provide.

Here is the precise point. In a pullback square with lower horizontal map $p:Y'\to Y$, universal fibrewise calculation gives an isomorphism
\[
 p^*f_*M_X\simeq f'_*M_{X'}.
\]
To conclude that $q:X'\to X$ is a $\DD$-equivalence, one must compare the cohomology of $Y$ and $Y'$ with coefficient $F=f_*M_X$. This coefficient is locally constant. It need not be constant. The assumption that $p$ is a $\DD$-equivalence compares cohomology only with constant coefficients. There is no applicable hypothesis at the last step of the displayed argument. This is the gap, and the examples on manuscript pages 87--95 make it an actual obstruction.

The surviving implication can be proved without this step. For an object $y$ of a base category, the categories $Y/y$ and $\{y\}$ have final objects. The functors between them, and the functors $Y/y_1\to Y/y_0$ induced by arrows of $Y$, are $\DD$-equivalences under the asphericity axiom. Pulling them back along $f$ gives exactly the inclusions $X_y\to X/y$ and the transition functors between the comma categories. Condition (a), applied also after a preliminary base change, therefore implies (b). Using undercategories gives (b$'$). Nothing in this argument proves the converse.

The correction in footnote 35, restricting the test categories to weakly contractible categories, is adequate once locally constant objects on such categories are constant. The qualification is essential: replacing weak contractibility by $\DD$-acyclicity is insufficient. Nor does connectedness of a base remove monodromy. We shall return below to a two-object covering that makes this visible with set-valued coefficients alone. Thus the sentence in footnote 36 about all truncated homotopy theories cannot be retained under the usual interpretation of $\DD_i$; its case $i=0$ already fails.
\end{dcomment}


\textbf{Définition 2.} Si les conditions équivalentes de la proposition 2 sont satisfaites, on dit que $f$ est une $\DD$-\textit{fibration}.

N.B. C'est une notion \textit{à la fois} locale et “colocale” sur $Y$ (alors que fibration gauche est locale, fibration droite colocale).

\printedpage{50}
% head: A. Grothendieck : Dérivateurs, Ch. I, Édition des Universités de Montpellier II et Paris VII
% foot: Version 8 juillet 2001

\textbf{Proposition 3.} a) \textit{La notion de $\DD$-fibration est stable par changement de base et par composition.}\footnote[38]{La stabilité par composition \textit{n'est} pas claire.}

\indent b) \textit{Soient}
\[
\begin{tikzcd}[column sep=normal, row sep=large]
X \arrow[rr, "f"] \arrow[dr, "p"'] & \phantom{X} & X' \arrow[dl, "p'"] \\
& Y &
\arrow[from=1-2, to=2-2, phantom, "\text{\footnotesize comm}"]
\end{tikzcd}
\]
\textit{avec $p,p'$ des $\DD$-fibrations (il suffit que $p,p'$ soient des $\DD$-fibrations cohomologiques, ou que ce soient des $\DD$-fibrations homologiques, ou que $p,p'$ soient propres, ou $p,p'$ lisses). Si $f$ induit des $\DD$-équivalences dans les fibres, alors $f$ est une $\DD$-équivalence.}\footnote[39]{Peut se préciser considérablement dans le cas où $p$, $p'$ sont des $\DD$-fibrations, \textit{cf.} prop. 6, page 84.}

\indent\textsc{Démonstration.} \textit{a}) est immédiat sous la condition \textit{a}) (prop. 2) de la notion de $\DD$-fibration. Pour \textit{b}), supposons que $f$ induise des $\DD$-équivalences sur les fibres. On veux prouver que % [sic] On veux -> On veut, verb agreement
\[
H^*_{\DD}(X,M) \longrightarrow H^*_{\DD}(X',M)
\]
est un isomorphisme, pour tout $M$ dans $\mathcal{A}$, mais cette flèche s'identifie à
\[
H^*_{\DD}(Y, p_*(M_X)) \longrightarrow H^*_{\DD}(Y, p'_*(M_{X'}))
\]
déduite de
\[
p_*(M_X) \longrightarrow p'_*(M_{X'}) \qquad\qquad (\text{NB}\ \ M_X \simeq f^*M_{X'}\,)
\]
\msp{page 82}
et il suffit de prouver que cette flèche est un isomorphisme. Pour ceci, on regarde fibre par fibre, et on utilise le fait que “les images directes se calculent fibre par fibre”.

\dgrblock{dgrb50}{Fibration terminology during the revision}
\begin{dcomment}{exposition}{p82:les images directes se calculent fibre par fibre}{Fibration terminology during the revision}
The notation temporarily conceals two different classes of functors. If ``$\DD$-fibration'' means the pullback-preservation condition (a), closure under composition is immediate: a $\DD$-equivalence is first pulled back along one factor and then along the other. Base change stability is immediate for the same reason. If the word instead means universal cohomological fibrewise calculation, that proof is unavailable. In a composite, the intermediate coefficient is a locally constant direct image, and the available comparison controls only direct images of coefficients pulled back from the base. The reservation in footnote 38 is consequently a substantive correction to the argument, rather than a doubt about the formal stability of condition (a).

The fibrewise comparison theorem in Proposition 3(b) survives. Write $f:X\to X'$ over $Y$, with projections $p$ and $p'$. The natural cohomological comparison is
\[
 p'_*M_{X'}\longrightarrow p_*M_X,
\]
obtained by applying $p'_*$ to the unit $M_{X'}\to f_*f^*M_{X'}$. If both direct images are computed on strict fibres, its component at $y$ is the cohomological map induced by $f_y:X_y\to X'_y$. A $\DD$-equivalence on each fibre therefore makes this map invertible by the detection axiom, and applying the direct image $Y\to e$ proves the assertion. The three displayed cohomological arrows in the printed proof point in the reverse direction; the theorem does not depend on that slip. The homological proof uses the counit and has the covariant direction.

The lemma on manuscript page 85 strengthens the sufficient criterion. The property that $f_y$ is a $\DD$-equivalence propagates across every arrow of $Y$, by applying two-out-of-three to the square of comma categories. It is therefore enough to test one object in each connected component. The proposed converse, from the total-space equivalence to fibre equivalences, fails because taking global cohomology can annihilate a nonzero local system. This does not contradict pointwise detection in $\DD(Y)$: the functor of global sections is generally not conservative.
\end{dcomment}


\textbf{Question.} Soit $f : X \to Y$ une $\DD$-fibration, et soit $F \in \DD(Y)$. Alors $f_*f^*F$ se calcule-t-il par passage aux fibres? Même question pour $f_!f^*F$. Pour la première propriété, concernant $f_*$, on montre que cela marche si $F$ est de la forme $p_*(M_Y)$, où $p : Y' \to Y$ est $\DD$-\textit{propre,} et $M \in \Ob\, \mathcal{A}$. D'autre part, la sous-catégorie pleine $C$ de $\DD(Y)$ formée des $F$ pour lesquels ça marche est (sauf erreur) “stable par $\DD$-$\mathsf{hot}$-$\varprojlim$” en le sens suivant : si on a une catégorie $I$ avec objet final $e$ formant une sous-catégorie fermée (\textit{i.e.} $e \to x \Rightarrow x = e$), et si on a $\mathcal{F} \in \DD(Y \times I)$, d'où $I^\circ \to \DD(Y)$, $i \mapsto \mathcal{F}_i$ (diagramme squelettique), si les $\mathcal{F}_i$ pour $i \neq e$ sont dans $C$, et si $\mathcal{F}$ est dans l'image essentielle de $j_*$, où $j : Y \times (I \setminus \{e\}) \hookrightarrow Y \times I$ est l'inclusion, alors $\mathcal{F}_e$ est aussi dans $C$.

\dgrblock{dgrb51}{Two useful classes of coefficients}
\begin{dcomment}{exposition}{p82:Pour la première propriété, concernant}{Two useful classes of coefficients}
The question about arbitrary coefficients should be separated from the assertion actually needed in the preceding proof. A universally cohomological fibration already calculates $f_*f^*F$ on strict fibres whenever $F$ is locally constant, under the local constancy axiom. Two further operations on coefficients explain the tentative observations here.

Suppose $p:Y'\to Y$ is $\DD$-proper, $M\in\DD(e)$, and $F=p_*M_{Y'}$. Form $X'=X\times_Y Y'$, with maps $q:X'\to X$ and $f':X'\to Y'$. Proper base change gives
\[
 f^*p_*M_{Y'}\simeq q_*M_{X'},\qquad
 f_*f^*p_*M_{Y'}\simeq p_*f'_*M_{X'}.
\]
The fibre formula for $p_*$ turns evaluation at $y$ into cohomology over $Y'_y$. Universal fibrewise calculation for $f'$ identifies the restriction of $f'_*M_{X'}$ to $Y'_y$ with the direct image along
$X_y\times Y'_y\to Y'_y$. The resulting object is
$H^*_{\DD}(X_y\times Y'_y,M)$.
On the other side, $F_y=H^*_{\DD}(Y'_y,M)$, and the product interchange formula identifies
$H^*_{\DD}(X_y,(F_y)_{X_y})$ with the same object. All identifications are the canonical mates, so this proves the desired comparison, not merely an accidental isomorphism between its endpoints. The printed $p_*(M_Y)$ must here be understood as $p_*(M_{Y'})$.

There is also a formal closure argument. Fix $y$. The two sides of the fibre comparison, as functors of $F\in\DD(Y)$, preserve coherent homotopy limits: restrictions preserve limits computed in the diagram variable, and right Kan extensions compose. Their comparison is a natural transformation between these continuous functors. If it is invertible on a coherent diagram of coefficients, it is invertible on its homotopy limit, by applying the two functors and then the canonical limit comparison. In particular this proves the punctured-cone assertion in the text, provided the indicated diagram shapes belong to $\Dia$. No claim about arbitrary incoherent diagrams in the category $\DD(Y)$ is required. This distinction matters because an arbitrary diagram in a homotopy category need not have a chosen coherent lift.
\end{dcomment}


\indent On peut cependant considérer que si on veux que $f : X \to Y$ satisfasse à toutes les propriétés cohomologiques et homologiques qui s'attachent à la notion d'une “fibration” (“localement triviale”), on devrait mettre dans les conditions que $f$ est \textit{lisse}, et \textit{propre}, pour que les changements de base par $f$ commutent aux foncteurs $g_*$, $g_!$. % [sic] on veux -> on veut, verb agreement

\dgrblock{dgrb52}{The fibre formula for coefficients from the base}
\begin{dcomment}{development}{p82:pour que les changements de base par $f$ commutent aux foncteurs}{The fibre formula for coefficients from the base}
The question on this page has an affirmative answer for the universal notion of fibration retained on manuscript pages 117--118. The restriction on the coefficient is essential to the argument: it may vary arbitrarily on the base, but the coefficient on the total category is its inverse image. We give the argument because it also answers the question about trivial fibrations on page 84. Throughout this block the domain contains the comma categories, categories of factorizations, and their pullbacks used below. The usual domain of all small categories has this property; so does the domain of finite posets when the categories under consideration are finite posets.

\textbf{Proposition (fibre calculation for inverse images).}\ours
Suppose that $\D$ takes admissible disjoint unions to products, has both homotopy Kan extensions and their pointwise formulas, and detects isomorphisms pointwise. All disjoint unions indexed by the hom-sets below are assumed admissible. Let $f:X\to Y$ be a weak cohomological $\D$-fibration after every base change. Then, for every $F\in\D(Y)$ and $y\in Y$, the canonical comparison
\[
 (f_*f^*F)_y\longrightarrow
 H^*_{\D}(X_y,(F_y)_{X_y})
\]
is an isomorphism. These comparisons also imply base change for $f_*f^*F$ along an arbitrary functor into $Y$. The dual assertion holds for $f_!f^*F$ when $f$ is universally a weak homological fibration.

Fix $y$ and write $H=X/y$, $S=X_y$, $v:H\to Y$ for $(x,f(x)\to y)\mapsto f(x)$, and $s:H\to e$, $t:S\to e$ for the projections. The two functors in the comparison are the right adjoints
\[
 T=s_*v^*,\qquad R=t_*t^*y^* :\D(Y)\longrightarrow\D(e).
\]
The pointwise formula identifies $T$ with evaluation of $f_*f^*$ at $y$. Their left adjoints are respectively $v_!s^*$ and $y_!t_!t^*$. It therefore suffices to prove that the left mate of the comparison,
\[
 y_!t_!t^*\longrightarrow v_!s^*,
\]
is invertible. We can test this on $M\in\D(e)$ and then at each object $z\in Y$.

For an arrow $\gamma:z\to y$, let $E_\gamma$ be the category of factorizations
\[
 z\xrightarrow{\beta}c\xrightarrow{\alpha}y,
 \qquad \alpha\beta=\gamma.
\]
A morphism from $(c,\beta,\alpha)$ to $(c',\beta',\alpha')$ is $h:c\to c'$ with $h\beta=\beta'$ and $\alpha'h=\alpha$. This category has initial object $(z,1_z,\gamma)$ and terminal object $(y,\gamma,1_y)$. Put
\[
 K_\gamma=X\times_Y E_\gamma,
\]
where $E_\gamma\to Y$ remembers the middle object. Thus an object of $K_\gamma$ is a factorization $z\to f(x)\to y$ of $\gamma$, and its morphisms are arrows of $X$ making the two triangles commute.

The pullback $K_\gamma\to E_\gamma$ is a weak cohomological fibration by hypothesis. Its strict fibre over the terminal factorization is $X_y$. Since slicing over a terminal object changes nothing, the first clause in the definition of a weak cohomological fibration says precisely that
\[
 X_y\longrightarrow K_\gamma,
 \qquad x\longmapsto(x,\gamma,1_y),
\]
is a $\D$-equivalence. Consequently it induces an isomorphism on homology with every constant coefficient $M$. Here we use the equivalence of the homological and cohomological tests for a $\D$-equivalence established on manuscript pages 61--62; no assertion about arbitrary local systems is being substituted for it.

Now $z\backslash v$ is the disjoint union of the $K_\gamma$ over the discrete set $Y(z,y)$. Indeed, its objects give the two arrows $z\to f(x)\to y$, and every morphism preserves their composite. The pointwise formula for left Kan extension and the coproduct axiom identify the component of the mate at $z$ with
\[
 \coprod_{\gamma\in Y(z,y)}H_*^{\D}(X_y,M)
 \longrightarrow
 \coprod_{\gamma\in Y(z,y)}H_*^{\D}(K_\gamma,M).
\]
Each summand is induced by the displayed inclusion of $X_y$ and is an isomorphism. Hence the mate is an isomorphism by pointwise detection, and so is the original comparison between right adjoints. This proves the fibre formula. For a base change $a:Y'\to Y$, both sides of the base-change map at $y'\in Y'$ identify with the same cohomology of $X_{a(y')}$ with value $F_{a(y')}$. Naturality of the comparisons gives the base-change assertion. Applying the argument with the opposite indexing categories and opposite coefficient categories gives the homological assertion.\hfill$\square$

The mechanism is the relative form of the factorization-category proof of homotopy exactness in \citet[Theorem 3.8]{GrothPontoShulman2014}. The full-faithfulness consequence in the next block has an all-derivator formulation in the homotopical epimorphism criterion of \citet[Theorem 6.10]{GrothStovicek2016Trees}. We have checked these references, as well as \citet{Maltsiniotis2005Asphericite} and \citet{Cisinski2009LocallyConstant}; we did not locate this precise answer to the page-82 question in those texts. The label marks our application and proof here, not a claim of priority. Notice what the proof does not require: neither properness of $f$ nor local constancy of $F$ on $Y$.
\end{dcomment}


\printedpage{51}
% head: A. Grothendieck : Dérivateurs, Ch. I, Édition des Universités de Montpellier II et Paris VII
% foot: Version 8 juillet 2001

\msp{page 83}

Signalons :

\textbf{Proposition 4.} \textit{Soit $f : X \to Y$} propre. \textit{Pour que $f$ soit une $\DD$-fibration, il faut et il suffit que pour tout $y_1 \to y_0$ dans $Y$, le foncteur $X/y_1 \to X/y_0$ induit soit une $\DD$-équivalence} (i.e. \textit{que $f$ soit une $\DD$-fibration} cohomologique). \textit{Dualement, si $f$ est $\DD$-lisse, pour que ce soit une $\DD$-fibration, il faut et il suffit que pour tout $y_1 \to y_0$ dans $Y$, $y_0\backslash X \to y_1\backslash X$ soit une $\DD$-équivalence.}

\indent On pourrait appeler $\DD$-fibrations \textit{parfaites} les foncteurs à la fois propres et lisses qui sont des $\DD$-fibrations cohomologiques (ou ce qui revient au même des $\DD$-fibrations homologiques).

\textbf{Corollaire.} \textit{Supposons $f$ $\Cat$-cofibration (donc propre). Pour que $f$ soit une $\DD$-fibration, il faut et il suffit que les foncteurs $X_{y_0} \to X_{y_1}$ (pour $y_0 \to u_1$ dans $Y$) soient des $\DD$-équivalences. Duallement, pour $\Cat$-fibration (donc lisse).} % [sic] u_1 -> y_1, the arrow is y_0 -> y_1 ; [sic] Duallement -> Dualement, spelling

\textbf{Proposition 5.} \textit{Soit $f : X \to Y$ dans} $\Dia$. \textit{Conditions équivalentes :}
\begin{enumerate}[label={}, leftmargin=2em, labelsep=0pt]
\item a) \textit{$f$ est une $\DD$-équivalence, et le reste par tout changement de base ;}
\item b) \textit{$f$ est une $\DD$-fibration à fibres $\DD$-contractiles.}
\end{enumerate}

\textsc{Démonstration.} \textit{a}) $\Rightarrow$ \textit{b}) : Que \textit{a}) implique que $f$ soit une $\DD$-fibration est immédiat sur la définition de celle-ci (condition \textit{a}) de la proposition 2), et que les fibres sont contractiles est tautologique.

\textit{b}) $\Rightarrow$ \textit{a}) : Il suffit de montrer que \textit{b}) implique que $f$ est une $\DD$-équivalence. Mais c'est un cas particulier de la proposition 3 \textit{b}) (en y prenant $X' = Y$).

\dgrblock{dgrb53}{Properness separates the two tests}
\begin{dcomment}{exposition}{p83:Mais c'est un cas particulier de la proposition 3}{Properness separates the two tests}
For a proper functor the first half of the cohomological-fibration condition is already supplied by proper base change: the inclusion $X_y\to X/y$ is coaspherical, hence a $\DD$-equivalence, and right Kan extensions calculate on strict fibres. What remains is local constancy of $f_*M_X$ for every $M$. This remaining condition is stable under arbitrary change of base. Indeed properness identifies the new direct image with the restriction of the old one, and restriction preserves local constancy. This proves the criterion of Proposition 4 in the universal cohomological sense. Its identification with the revised bilateral fibration notion uses the local-constancy hypotheses explained at manuscript pages 97--100.

For a Grothendieck opfibration, a choice of cocartesian lifts gives transition functors $\alpha_!:X_{y_0}\to X_{y_1}$ for arrows $\alpha:y_0\to y_1$. They are defined up to canonical natural isomorphism. The map on constant-coefficient cohomology between the two strict fibres is the map induced by $\alpha_!$, with the cohomological direction reversed. Consequently local constancy says precisely that every transition functor is a $\DD$-equivalence. The criterion is independent of the chosen cleavage. The dual statement uses cartesian transport in a Grothendieck fibration.

The later adjective ``trivial'' is also justified without the false converse of Proposition 2. A universal $\DD$-equivalence has $\DD$-contractible strict fibres by pulling back to points. Its further pullbacks are again universal equivalences. It has the cohomological and homological fibration tests because the relevant maps form triangles whose other two sides are these universal equivalences. Conversely, a universal cohomological fibration with $\DD$-contractible fibres has $M_Y\to f_*M_X$ invertible at every point, hence invertible in $\DD(Y)$. Taking global cohomology proves that it is a $\DD$-equivalence, and the argument repeats after every base change. Thus Proposition 5 survives the revision. The stronger full-faithfulness question on page 84 requires the separate factorization-category argument given in its commentary.
\end{dcomment}


\msp{page 84}

Quand ces conditions sont satisfaites, on s'attend à ce que le foncteur
\[
f^* : \DD(Y) \longrightarrow \DD(X)
\]
soit pleinement fidèle, \textit{i.e.} que les morphisme d'adjonction % [sic] les morphisme -> les morphismes, plural agreement
\[
F \longrightarrow f_*f^*F \quad , \qquad f_!f^* \longrightarrow F % [sic] f_!f^* -> f_!f^*F, the object F is missing
\]
dans $\DD(Y)$ soient des isomorphismes. (C'est vrai si $f$ est $\DD$-propre ou $\DD$-lisse, cf. cor. 3, p.71.) Cela revient encore à dire que les $f_*f^*F$ se calculent par passage aux fibres, ou encore itou pour les $f_!f^*F$. Cela nous ramène à la question de la page 82.

\dgrblock{dgrb54}{Universal equivalences give fully faithful restriction}
\begin{dcomment}{development}{p84:Cela nous ramène à la question de la page 82.}{Universal equivalences give fully faithful restriction}
The expectation in the preceding paragraph follows from the fibre formula just proved. It is stronger than the assertion that $f$ preserves cohomology of constant objects, because it controls all morphisms between arbitrary coefficients on $Y$.

\textbf{Corollary (the question about trivial fibrations).}\ours
Under the domain and derivator hypotheses of the preceding development, suppose that every pullback of $f:X\to Y$ is a $\D$-equivalence. Then
\[
 f^*:\D(Y)\longrightarrow\D(X)
\]
is fully faithful. Both the unit $F\to f_*f^*F$ and the counit $f_!f^*F\to F$ are isomorphisms, for every $F\in\D(Y)$.

To prove this, first pull back over a point $y\to Y$. The hypothesis says that $X_y$ is $\D$-aspherical. Pulling back over $Y/y$ likewise shows that $X/y$ is $\D$-aspherical, since $Y/y$ has a terminal object. The inclusion $X_y\to X/y$ and all transition functors between these comma categories are therefore $\D$-equivalences by two-out-of-three. The same argument works after any base change. Thus $f$ is universally a weak cohomological fibration. At $y$, compose the adjunction unit with the comparison proved in the preceding block:
\[
 F_y\longrightarrow(f_*f^*F)_y
 \longrightarrow H^*_{\D}(X_y,(F_y)_{X_y}).
\]
The composite is the constant-coefficient unit for $X_y\to e$, which is invertible by asphericity. The second arrow is invertible by the proposition. The first arrow is therefore invertible, and pointwise detection proves that $F\to f_*f^*F$ is an isomorphism. This is equivalent to full faithfulness of $f^*$. Since $f^*$ also has the left adjoint $f_!$, full faithfulness makes the other displayed adjunction counit invertible as well.\hfill$\square$

There is also a useful direct connection with the modern criterion. For each arrow $\gamma:z\to y$, the category $E_\gamma$ of its factorizations has an initial and a terminal object. Universality makes $K_\gamma=X\times_YE_\gamma$ aspherical. These are exactly the factorization categories controlling full faithfulness. In the special case of weak equivalences of nerves, \citet[Theorem 6.10]{GrothStovicek2016Trees} therefore supplies the same conclusion simultaneously for every derivator. The argument above gives the coefficient-relative conclusion asked for in the manuscript and explains why universality, rather than the single global equivalence $X\to Y$, is the relevant hypothesis.
\end{dcomment}


\textbf{Définition 3.} On dit alors que $f$ est une $\DD$-fibration \textit{triviale}.

\textbf{Proposition 6} (précise prop. 3 page 81). \textit{Considérons un triangle commutatif dans} $\Dia$
\[
\begin{tikzcd}[column sep=small, row sep=large]
X \arrow[rr, "f"] \arrow[dr, "p"'] & & X' \arrow[dl, "p'"] \\
& S & & & .
\end{tikzcd}
\]

\printedpage{52}
% head: A. Grothendieck : Dérivateurs, Ch. I, Édition des Universités de Montpellier II et Paris VII
% foot: Version 8 juillet 2001

\indent a) \textit{Supposons que $p$ et $p'$ soient des $\DD$-fibrations gauches, ou que $p$ et $p'$ soient des $\DD$-fibrations droites. Soit $S_0$ une partie de $S$ qui rencontre toute composante connexe de $S$. Pour que $f$ soit une $\DD$-équivalence, il} suffit \textit{que pour tout $s_0 \in S_0$, $f$ induise une $\DD$-équivalence $X_{s_0} \to Y_{s_0}$.}

\indent b) \textit{Cette condition est aussi nécessaire}\footnote[40]{Pas prouvé!\footnotemark[41] Même si $X' = S$, ce n'est pas clair, même si $\DD = \mathsf{HOTAB}$, $f$ propre et lisse. Mais OK si $S$ est $\DD$-contractile.}\footnotetext[41]{\underline{Faux}, \textit{cf.} contre-exemple, p. 87.} \textit{si $p$ et $p'$ sont des $\DD$-fibrations. Donc dans ce cas, $f$ est une $\DD$-équivalence si et seulement si $f$ induit des $\DD$-équivalences fibre par fibre.}

\msp{page 85}

\textsc{Démonstration.} a) Il suffit, en vertu de la proposition 3, de montrer que l'hypothèse implique que pour \textit{tout} $s \in S$, $f$ induit une $\DD$-équivalence $f_s : X_s \to X'_s$. Cela résulte aussitôt du lemme :

\textbf{Lemme.} \textit{Si $p$ et $p'$ sont toutes deux soit des $\DD$-fibrations gauches, soit des $\DD$-fibrations droites, et si $s_1 \to s_0$ est une flèche dans $S$, alors $f_{s_1}$ est une $\DD$-équivalence si et seulement si $f_{s_0}$ l'est. (Donc l'ensemble $S' \subset \Ob\, S$ des $s \in S$ tels que $f_s$ soit une $\DD$-équivalence est une réunion de composantes connexes de $S$.)}

\indent Supposons par exemple que $p,p'$ soient des $\DD$-fibrations gauches. Considérons, pour $s \in S$, le diagramme
\[
\begin{tikzcd}[row sep=large]
X_s \arrow[r, "f_s"] \arrow[d] & X'_s \arrow[d] \\
X/s \arrow[r, "f/s"'] & X'/s & .
\end{tikzcd}
\]
Les flèches verticales sont des $\DD$-équivalences par l'hypothèse sur $p,p'$, donc $f_s$ est une $\DD$-équivalence si et seulement si $f/s$ l'est. Considérons maintenant
\[
\begin{tikzcd}[row sep=large]
X/s_1 \arrow[r, "f/s_1"] \arrow[d] & X'/s_1 \arrow[d] \\
X/s_0 \arrow[r, "f/s_0"'] & X'/s_0 & .
\end{tikzcd}
\]
Ici encore, l'hypothèse sur $p,p'$ implique que les flèches verticales sont des $\DD$-équivalences, donc $f_{s_1}$ est une $\DD$-équivalence si et seulement si $f_{s_0}$ l'est.

\msp{page 86}

Si $Y' = S$, l'énoncé (non prouvé) dit que si $f : X \to S$ est une $\DD$-fibration et une $\DD$-équivalence, alors les fibres sont $\DD$-contractiles. Ce n'est pas sûr même pour une projection $X = T \times S \to S$. Si $\eta = \mathit{hot}_0(S)$, $\xi = \mathit{hot}_0(T)$, la question se réduit à celle-ci : Si pour tout $M$ dans $\mathcal{A}$, on a : l'homomorphisme
\[
\alpha_M : \xi \times_{\DD} M \longrightarrow M
\]
\printedpage{53}
% head: A. Grothendieck : Dérivateurs, Ch. I, Édition des Universités de Montpellier II et Paris VII
% foot: Version 8 juillet 2001
satisfait
\[
\eta \times_{\DD} \alpha_M : \eta \times_{\DD} (\xi \times_{\DD} M) \xrightarrow{\sim} \eta \times_{\DD} M
\]
est un isomorphisme, alors $\alpha_M$ est lui même un isomorphisme.

\dgrblock{dgrb55}{Why a product cannot supply the counterexample}
\begin{dcomment}{development}{p86:est lui même un isomorphisme.}{Why a product cannot supply the counterexample}
The product case singled out on manuscript page 86 has a short resolution that does not use an additive coefficient theory. We give it to isolate the role of twisting in the counterexamples which follow. It is a formal retract argument; we claim no priority for it. The locally constant-coefficient results of \citet{Cisinski2009LocallyConstant} address the general homotopical framework, but we did not find this particular paragraph of the manuscript treated there.

\textbf{Cancellation for a nonempty product base.}\ours\
Assume the action of indexing homotopy types on $\DD(e)$ is functorial and that $S$ has an object. If the product projection $S\times T\to S$ is a $\DD$-equivalence, then $T\to e$ is a $\DD$-equivalence.

\emph{Proof.} Choose $s:e\to S$ and write $p:S\to e$. Then $ps=1_e$. For each coefficient $M$, let $\alpha_M:T\times_{\DD}M\to M$ be the augmentation. Functoriality gives a retract diagram in the arrow category of $\DD(e)$,
\[
 \alpha_M\longrightarrow S\times_{\DD}\alpha_M
                   \longrightarrow\alpha_M,
\]
whose composite is the identity: the first map is induced by $s$, the second by $p$. The middle arrow is an isomorphism by the product hypothesis and Fubini for homology. A retract of an isomorphism is an isomorphism. Explicitly, if an arrow $a$ is a retract of an invertible arrow $b$, with inclusions $i_0,i_1$ and retractions $r_0,r_1$, then $r_0b^{-1}i_1$ is an inverse of $a$, as the two commutative squares show. Hence every $\alpha_M$ is invertible. This is the asserted $\DD$-contractibility of $T$. \hfill$\square$

The nonempty-base hypothesis cannot be omitted: $\varnothing\times T\to\varnothing$ is an isomorphism for every $T$. The argument establishes that the proposed product example cannot work; it does not establish the analogous assertion for a bundle with monodromy. The next example exploits exactly that distinction.
\end{dcomment}


\msp{page 87}

Voici un \textit{contre-exemple}\footnote[42]{Voir un contre-exemple encore plus simple pages 93, 94. (Mais il a l'inconvenient de ne pas être à fibres connexes.) % [sic] inconvenient -> inconvénient, missing accent
} que je donne d'abord dans le contexte des espaces topologiques, au lieu de $\Cat$. Je vais construire une fibration localement triviale et non triviale
\[
\begin{tikzcd}[column sep=0.2em, row sep=large]
X \arrow[d, "f"'] & \supset S^1 & \\
S & = S^1 & \qquad ,
\end{tikzcd}
\]
de fibre $S^1$, telle que $f$ soit une $\DD$-équivalence pour tout $\DD = \mathsf{Der}(k\text{-}\mathsf{Mod})$, $k$ étant un anneau où 2 est inversible. Cela signifie que
\[
H_*(X,k) \longrightarrow H_*(S,k)
\]
est un isomorphisme modulo 2-torsion (noyau et conoyau sont annulés par une puissance de 2 - en fait par 2 ici), ou encore $H_*(X,k) \xrightarrow{\sim} H_*(S,k)$ si 2 inversible dans $k$. Mais la fibre $S^1$ n'est pas $\DD$-contractile (sauf si $k = 0$!).\footnote[43]{Exemple beaucoup plus simple encore : on prend $X = S^1$, et $X \to S$ le revêtement à deux feuillets. Si on prend le revêtement à $n$ feuillets, c'est aussi OK si $k$ anneau où $n$ est inversible.}

\indent La fibration est donnée à isomorphisme unique près, en se fixant un point $s \in S$, et en se donnant un automorphisme $u$ de la fibre $X_s$ (“monodromie”). (La classe d'isomorphie du fibré ne dépend que de la classe d'isotopie de $u$). Je prend pour $u$ la symétrie (inversant l'orientation) $z \mapsto z^{-1}$ du cercle standard $\{z \in \mathbb{C} \mid |z| = 1\}$. Cette fibration “se réalise” par des ensembles ordonés $X \to S$, $S^1$ étant représenté (sur la base, comme sur les fibres) par % [sic] Je prend -> Je prends, verb agreement ; [sic] ordonés -> ordonnés, spelling
\[
\begin{tikzcd}[column sep=normal, row sep=normal]
& & a & & \\
S \ = & 0 \arrow[ur] \arrow[dr] & & 0' \arrow[ul] \arrow[dl] & = \ S' \amalg_R S'' \quad , \\
& & b & &
\end{tikzcd}
\]
où $R$ est la catégorie discrète ayant deux éléments $a$, $b$. On définit $X$ comme $X' \amalg_Y X''$, où $X' = X|S'$, $X'' = X|S''$, $Y = X|R$.\footnote[44]{N.B. $S'$, $S''$, $R = S' \cap S''$ fermés dans $S$, $X'$, $X''$, $Y = X' \cap X''$ fermés dans $X$.}

% ---- from Der-ChI-234_p056-060.tex
\printedpage{54}
% head: A. Grothendieck : Dérivateurs, Ch. I, Édition des Universités de Montpellier II et Paris VII

\msp{page 88}
On pose
\[ X' = S' \times \bar S \qquad\qquad X'' = S'' \times \bar S \]
($\bar S$ copie conforme de $S$, cf ci-dessous) et il faut, pour recoller $X'$ et $X''$, se donner un isomorphisme entre $X'|R \simeq R \times \bar S = \bar S \amalg \bar S$ et $X''|R \simeq R \times \bar S = \bar S \amalg \bar S$ qui soit un $R$-foncteur, \textit{i.e.} qui provient de deux isomorphismes
\[
\begin{aligned}
\bar S &\xrightarrow{u_a} \bar S \\
\bar S &\xrightarrow{u_b} \bar S \qquad .
\end{aligned}
\]
On prend
\[
\begin{aligned}
u_a &= \mathsf{id}_{\bar S} \\
u_b &= \text{symétrie par rapport à l'axe } \bar 0, \bar 0' \ \textit{i.e.} % Doubt resolved: PDF leaf 18 (printed p.54) reads u_b, with plain b.
\\
&\phantom{=}\quad u_b \text{ laisse fixes } \bar 0, \bar 0' \text{ et intervertit } \bar a \text{ et } \bar b.
\end{aligned}
\]
(C'est plus clair si on désidentifie la base et la fibre, en prenant fibre
\[
\bar S \quad = \quad
\begin{tikzcd}
 & \bar a & \\
\bar 0 \arrow[ur] \arrow[dr] & & \bar 0' \arrow[ul] \arrow[dl] \\
 & \bar b &
\end{tikzcd}
\qquad ).
\]

\dgrblock{dgrb56}{The Klein bottle and its sign coefficient}
\begin{dcomment}{examples}{p88:plus clair si on désidentifie la base et la fibre}{The Klein bottle and its sign coefficient}
Take a nonzero commutative ring $k$ in which $2$ is invertible. Let $r:S^1\to S^1$ be complex conjugation and form its mapping torus
\[
 K=(S^1\times[0,1])/((z,1)\sim(r(z),0)).
\]
Projection to $[0,1]/(0\sim1)$ gives the Klein-bottle bundle $f:K\to S^1$. We use the cohomological convention in which $k[-1]$ has its copy of $k$ in degree $1$. Grothendieck's displayed $k[1]$ is explicitly annotated ``en degr\'e 1'' and should be translated accordingly, not corrected as a transcription error.

A cell structure has one vertex, two edges $a,b$, and one two-cell with attaching word $aba^{-1}b$. Here $a$ projects once around the base and $b$ runs around the fibre. The cellular complex is
\[
 0\longrightarrow k\xrightarrow{(0,2)}k a\oplus k b
                         \xrightarrow{0}k\longrightarrow0.
\]
The map to the cellular complex of the base sends $a$ to its generator and kills $b$ and the two-cell. Its kernel is the two-term complex $k\xrightarrow{2}k$, which is contractible. Consequently the projection induces an equivalence with every constant derived $k$-module coefficient, including unbounded complexes: the displayed chain contraction survives applying the appropriate Hom complex. Integrally, one instead gets $H_1(K;\mathbb Z)=\mathbb Z\oplus\mathbb Z/2$, and the kernel of the map to $H_1(S^1;\mathbb Z)$ is the indicated $2$-torsion. The fibre still has $H_1(S^1;k)=k$ and is not $\DD$-contractible.

The same computation in local-coefficient language explains why global cohomology loses the fibre. Let $L$ be the rank-one local system on the base with monodromy $-1$. The fibre cohomology gives
\[
 Rf_*k_K\simeq k_{S^1}\oplus L[-1].
\]
This splitting can be obtained from the cellular circle complex with its reflection action. Alternatively, the fixed point of the reflection gives a section and splits the degree-zero summand. The remaining summand is $L[-1]$. Global sections of $L$ are computed by $k\xrightarrow{-2}k$, in degrees $0,1$; the complex is zero in the derived category. Thus the extra summand is visible over each point and invisible after taking global sections.

The finite categorical model in the text has four objects in the base and four in each fibre. Choose $S'=\{0,a,b\}$ and $S''=\{0',a,b\}$. Their intersection is the discrete category $\{a,b\}$. Gluing the products with the four-object circle by the identity over $a$ and by reflection over $b$ produces sixteen objects. Each arrow in the base transports the fibre by an isomorphism. The projection is therefore both a Grothendieck fibration and an opfibration, and remains so after every base change. Its nerve realizes the same monodromy bundle. This supplies the categorical counterexample required by the manuscript, without appealing only to a topological analogy. Cellular calculations and the general language of local coefficients are reviewed in \citet[\S\S2.2, 3.H]{Hatcher2002}.
\end{dcomment}


Pour voir que $X$ est une $\D$-fibration sur $S$ (quelque soit $\D$) prouvons qu'il est \textit{propre}\footnote[45]{N.B. $f : X \to S$ est également \textit{lisse}, car c'est un fibré trivial (un produit) au dessus de deux ouverts $S - \{a\}$, $S - \{b\}$ qui recouvrent $S$. Mieux : $f$ est à la fois $\Cat$-fibrant et $\Cat$-cofibrant (car cette condition n'utilise que les propriétés des $X \times_S \boldsymbol{\Delta}_2 \to \boldsymbol{\Delta}_2$, pour $\boldsymbol{\Delta}_2 \to S$).} sur $S$ - on est ramenés à voir que $X' \to S'$ et $X'' \to S''$ sont propres, ce qui est trivial. Pour voir que c'est une $\D$-fibration, on est ramenés également au cas de $X' \to S'$ et $X'' \to S''$, où c'est aussi trivial. On a utilisé le : % [sic] quelque soit -> quel que soit, spelling

\textbf{Lemme.} \textit{Soient $f : X \to S$, $(S_i)_{i \in I}$ une famille de sous-catégories fermées recouvrant $S$. Pour tout $i$, soit $X_i = X \times_S S_i \xrightarrow{f_i} S_i$. Pour que $f$ soit propre, il faut et il suffit que les $f_i$ le soient. Pour que $f$ soit une $\D$-fibration propre il faut et il suffit que les $f_i$ le soient.}

\dgrblock{dgrb57}{Why the closed-cover test is sufficient}
\begin{dcomment}{exposition}{p88:Pour que $f$ soit propre, il faut et il suffit que les}{Why the closed-cover test is sufficient}
The gluing lemma has a local proof which does not require a descent theorem for arbitrary closed covers. A closed full subcategory here is a cosieve: whenever it contains the source of an arrow it contains its target. In the lifting-category criterion for properness, fix $x\in X$ and an arrow $f(x)\to s$ in the base. Choose a member $S_i$ of the cover containing $f(x)$. Its cosieve property forces $s$ and every intermediate object relevant to this lift into $S_i$. The entire category of lifts is therefore the same one computed in $X_i\to S_i$. Its $\DD$-contractibility follows from properness there. Necessity follows by base change.

Once properness has been obtained globally, the fibration condition asks whether the fibre-cohomology comparison along each arrow is invertible. Choose $S_i$ containing its source; the arrow and both endpoints lie in that subcategory. The local condition proves the comparison, and the pointwise criterion proves the global condition. This explains both assertions of the lemma. For the displayed four-object circle, the two closed subcategories $S'$ and $S''$ contain respectively $0$ and $0'$, together with both targets $a,b$, so they meet exactly these requirements. The dual open-cover argument explains smoothness.
\end{dcomment}


\msp{page 89}
Voici le petit calcul qui montre que $f : X \to S$ est une $\D$-équivalence. On fait le calcul en cohomologie.\footnote[46]{Mais dans le cas du revêtement à 2 feuillets de $S$, il n'y a pas même de calcul à faire!} On doit prouver que
\[ H^*(S,k) \longrightarrow H^*(X,k) \tag*{$(*)$} \]
\printedpage{55}
% head: A. Grothendieck : Dérivateurs, Ch. I, Édition des Universités de Montpellier II et Paris VII
est un isomorphisme (et de même pour tout coefficient constant $M$ sur $S$). Or
\[ H^*(X,k) \overset{\mathrm{def}}{=} \RR\Gamma_X(k_X) = \RR\Gamma_S(\RR f_*(k_X)) \quad , \]
et $(*)$ est induit par
\[ i : k_S \longrightarrow \RR f_* k_X \quad . \]
Or on a
\[ \RR f_* k_X \simeq \underbrace{k_S}_{= \mathsf{Im}(i)} \quad \oplus \quad \underbrace{k_S[1]}_{k_S \text{ en degré } 1}{}^{\underline{\mathrm{tordu}}} \]
savoir ``tordu'' par l'action de monodromie sur $S = S^1$, laquelle est la symétrie $\lambda \mapsto -\lambda$. On a
\[ \RR\Gamma_S(k_S[1]^{\mathrm{tordu}}) = H^1(\underbrace{\pi_1(S)}_{= \mathbb{Z}}\, , k_S^{\mathrm{tordu}}) = k_{(\mathbb{Z})} = k/(u - \mathsf{id})k = k/2k = 0 \ , \]
d'où le résultat.

Notons que si on regarde
\[ X' = X \times_S X \xrightarrow{f'} S \quad , \]
on a par Künneth relatif
\[ \RR f'_*(k_{X'}) \simeq \RR f_*(k_X) \otimes^{\mathbb{L}} \RR f_*(k_X) \simeq k_S \oplus k_S[1]^t \oplus k_S[1]^t \oplus \underbrace{k_S[1]^t \otimes k_S[1]^t}_{\simeq k_S[2] \text{ non tordu}} \ , \]
d'où
\[ H^*_{\D}(X',k) \simeq \RR\Gamma_{X'}(k_{X'}) \simeq \RR\Gamma_S(\RR f'_*(k_{X'})) \simeq \RR\Gamma_S(k_S) \oplus \RR\Gamma_S(k_S)[2] \]
et
\[ H^*_{\D}(S,k) \longrightarrow H^*_{\D}(X',k) \]
n'est \textit{pas} un isomorphisme.

\msp{page 90}
Ainsi $f : X \to S$ qui est une $\D$ fibration (propre et lisse) satisfait bien les conditions équivalentes $b$), $b'$) de la proposition 2 (page 79), mais pas la condition $a$), puisque le changement de base
\[
\begin{tikzcd}[column sep=large, row sep=large]
X' = X \times_S X \arrow[r, "{q = pr_1}"] \arrow[d, "{f' = pr_2}"'] & X \arrow[d, "f"] \\
S' = X \arrow[r, "{p = f}"'] & S
\end{tikzcd}
\]
est une $\D$-équivalence $S' \to S$, alors que $X' \to X$ ne l'est pas.

\dgrblock{dgrb58}{Self-pullback restores the missing cohomology}
\begin{dcomment}{examples}{p90:alors que $X' \to X$ ne l'est pas.}{Self-pullback restores the missing cohomology}
The self-pullback $K\times_{S^1}K$ is the mapping torus of
$r\times r:T^2\to T^2$. Its fibre has cohomology $k$ in degrees $0$ and $2$, and $k^2$ in degree $1$. The monodromy acts trivially in degrees $0,2$ and as $-1$ on both degree-one generators. The degree-two action is $(-1)(-1)=1$. The two sign summands again have acyclic global sections, whereas the degree-two summand does not. Hence
\[
 H^j(K\times_{S^1}K;k)=
 \begin{cases} k,&j=0,1,2,3,\\0,&\text{otherwise}.\end{cases}
\]
The same conclusion follows directly from the mapping-torus complex: in each fibre degree its extra differential is $1-(r\times r)^*$, which is either zero or multiplication by $2$. This calculation avoids extension ambiguities and is valid over the stated ring.

The map $K\to S^1$ is a constant-coefficient equivalence, but its pullback $K\times_{S^1}K\to K$ cannot be one, because the source has the additional nonzero groups in degrees $2,3$. Thus even a bundle which is both proper and smooth in the categorical sense does not preserve all constant-coefficient equivalences under pullback. Properness controls the base-change comparison for a direct-image functor; it does not assert that the base-change map is an equivalence for whichever weak-equivalence class one has chosen on the base.

The easier covering example in footnote 43 exhibits the same distinction. The degree-$n$ map $S^1\to S^1$ is a $k$-homology equivalence when $n$ is invertible in $k$, since it is the identity in degree $0$ and multiplication by $n$ in degree $1$. Its strict fibre has $n$ points. Its self-pullback is a disjoint union of $n$ circles, so the pulled-back projection fails already on $H_0$ for $n>1$. This example also avoids any question about a derived splitting.
\end{dcomment}


Voici un autre contre-exemple, sur le même principe mais un peu plus simple. Soit $G$ un groupe fini (par exemple $G = \mathbb{Z}/2\mathbb{Z}$), faisons le opérer sur un objet $Z$ de $\Cat$, soit $S = [G] = B_G$ la catégorie réduite à un objet $a$, avec $G = \mathsf{End}(a)$. Soit $E_G$ la catégorie (isomorphe à $(B_G)/a$)
\[ \mathsf{Ob}\ E_G = G \quad , \qquad \mathsf{Hom}_{E_G}(g,g') = \{ h \in G \mid g = g'h \} \quad , \]
composition évidente - alors $E_G$ est un groupoïde trivial (1-connexe). Le goupe $G$ y opère à gauche librement par translations, et $B_G = E_G/G$. % [sic] goupe -> groupe, misprint

\printedpage{56}
% head: A. Grothendieck : Dérivateurs, Ch. I, Édition des Universités de Montpellier II et Paris VII

\msp{page 91}
Si $G$ opère à gauche sur $Z$, on peut alors l'utiliser pour tordre $Z$, par
\[ X = E_G \times_G Z = (E_G \times Z)/\text{opérations ``diagonales'' de G} \quad . \]
Donc on a un diagramme cartésien
\[
\begin{tikzcd}
 & X' \arrow[d, equal] & X'' \arrow[d, equal] \\
X \arrow[d, "f"'] & Z \times E_G \arrow[l] \arrow[d, "pr_2"] & Z \arrow[l, dashed, "{\overset{\D}{\sim}}"'] \arrow[d] \\
B_G \arrow[d, equal] & E_G \arrow[l] \arrow[d, equal] & e \arrow[l, dashed, "{\overset{\D}{\sim}}"'] \arrow[d, equal] \\
S & S' & S''
\end{tikzcd}
\]
Comme le changement de base $S' = E_G = (B_G)/a \to B_G = S$ transforme $X$ sur $B_G$ en un produit (lequel est un $\D$-fibré propre et lisse sur $S'$), on voit que $X$ sur $B_G$ est un $\D$-fibré propre et lisse (car c'est là une condition locale sur $B_G$).

On suppose maintenant que $G$ est fini d'ordre $n > 1$, et soit $k$ un anneau où $n$ est invesible, $\D = \mathsf{Der}(k - \mathsf{Mod})$. Alors $E_G \to B_G$ est une $\D$-équivalence, ou ce qui revient au même $e \to B_G$ l'est.\footnote[47]{En fait, $e$, $E_G$, $B_G$ sont $\D$-asphériques.} Pourtant, après changement de base soit par $S' = E_G$, soit par $S'' = e$, $X' \to X$ (ou ce qui revient au même, $X'' \to X$) n'a guère de raison d'en être aussi. On a en cohomologie % [sic] invesible -> inversible, misprint
\[ \RR\Gamma_X(k_X) \simeq \RR\Gamma_{BG}(f_*(k_X)) \simeq \RR\Gamma_G(G, H^*(Z,k)) \]
\msp{page 92}
donc ($n$ étant inversible dans $k$) :
\[ H^*(X,k) = \RR\Gamma_X(k_X) \simeq H^*(Z,k)^G \longrightarrow H^*(Z,k) \quad , \]
ce n'est pas en général un isomorphisme. On peut prendre par exemple $G = \mathbb{Z}/2\mathbb{Z}$, $Z \sim S^1$ comme avant, et $G$ opérant par symétrie de $S^1$ comme précédemment.

Cet exemple est particulièrement convainquant, car c'est un exemple \textit{avec $S$, $S'$ $\D$-asphériques,} et $S' \to S$ un isomorphisme local (en fait, un revêtement principal de $S$). C'est aussi un exemple d'un $\D$-fibré lisse et propre $X \to S$, sur une base \textit{asphérique} $S$, et tel que l'inclusion d'une fibre dans $X$ % [sic] convainquant -> convaincant, spelling of the adjective
\[ X_s \longrightarrow X \]
ne soit \textit{pas} une $\D$-équivalence.

Le seul reproche qu'on puisse faire à cet exemple, en comparaison avec le précédent, c'est que les catégories $X, S$ ne sont pas des ensembles ordonnés. Et $S$ ne peut être un ensemble ordonné fini, car l'espace classifiant d'un groupe fini a de la cohomologie en toutes dimensions (sauf erreur).

\dgrblock{dgrb59}{Averaging over a finite group}
\begin{dcomment}{examples}{p92:Le seul reproche qu'on puisse faire à cet exemple}{Averaging over a finite group}
The second counterexample places the same phenomenon over a base which is already $\DD$-acyclic. Let $G$ be finite of order invertible in $k$. The category $E_G$ has one arrow between every ordered pair of objects of $G$; in particular each object is both initial and final. The quotient by left translation is $BG$. A $G$-category $Z$ gives an associated bundle
\[
 X=(E_G\times Z)/G\longrightarrow BG.
\]
An arrow in the base acts on the fibre by an isomorphism, so this is again a fibration and an opfibration with universally locally constant fibre transport. Pulling back to the unique object of $BG$ recovers $Z$.

For a $k[G]$-module $V$, the averaging operator
\[
 e_G(v)=|G|^{-1}\sum_{g\in G}gv
\]
is an idempotent with image $V^G$. More than a splitting, it proves exactness of invariants: an invariant element of a quotient can be lifted arbitrarily and its lift averaged. The same argument applies degreewise to complexes, so cohomology commutes with invariants. Thus
\[
 H^*(X;k)=H^*(Z;k)^G,\qquad H^*(BG;k)=k
\]
with the second $k$ in degree zero. The first formula is a formula for cohomology; a functorial complex-level statement uses derived sections of $Z$ with their action, followed by invariants. It does not require choosing a splitting of an arbitrary complex into its cohomology groups.

Take $G=C_2$ and let $Z$ be the four-object circle with reflection. Its cohomology is the trivial representation in degree zero and the sign representation in degree one. Therefore $X$ and $BG$ are $\DD$-acyclic, whereas their common fibre $Z$ is not. This disproves the qualification ``OK si $S$ est $\DD$-contractile'' in footnote 40 as well as the unqualified necessity assertion of Proposition 6(b). The replacement needed there is complete $\DD$-asphericity, which controls locally constant coefficients, or ordinary weak contractibility under the appropriate axiom.

An even smaller choice is the discrete regular $G$-set $Z=G$. Then $X$ is equivalent to the contractible category $E_G$, while its fibre has $|G|$ connected components. This is a wholly categorical example with finitely many objects, although its nerve is not finite-dimensional. A nontrivial finite group cannot have a finite-poset classifying space: if it did, the finite covering corresponding to its universal cover would be a finite contractible complex, so its Euler characteristic would be both $1$ and $|G|$ times an integer. The phrase ``cohomology in all dimensions'' in the manuscript should be read as the obstruction to a bounded-dimensional classifying space, not as a claim of nonvanishing in every positive degree for every coefficient ring.
\end{dcomment}


\printedpage{57}
% head: A. Grothendieck : Dérivateurs, Ch. I, Édition des Universités de Montpellier II et Paris VII

\msp{page 93}
Mais on peut s'inspirer de cet exemple pour le modifier de façon qu'il se transpose à des ensembles ordonnés. Il suffit de réaliser, pour $G$ donné :

\textit{a}) Une opération \textit{libre} de $G$ sur un objet $E$ de $\Cat$, telle que $H^*(E,k)^G \simeq k$, \textit{i.e.} $B = E/G$ soit $\D$-asphérique.

\textit{b}) Une opération de $G$ sur un objet $Z$, telle que l'inclusion
\[ H^*(Z,k)^G \longrightarrow H^*(Z,k)^G \] % [sic] H^*(Z,k)^G \longrightarrow H^*(Z,k)^G -> H^*(Z,k)^G \longrightarrow H^*(Z,k), an inclusion of the invariants (cf. p. 56)
ne soit pas \texttt{[un]} isomorphisme.

\textit{c}) Choisir pour $E$, $Z$ des ensembles ordonnés, et s'assurer que $B = E/G$ est également représenté par une catégorie ordonnée. Sauf erreur, il en sera de même pour $X = E \times Z/G$. On aura en tout les cas % [sic] en tout les cas -> en tous les cas, agreement
\[
\begin{tikzcd}[column sep=large]
 & Z \times E \arrow[d, equal] & \\
X \arrow[d] & X' \arrow[l] \arrow[d] & Z \arrow[l] \arrow[d] \\
S \arrow[d, equal] & S' \arrow[l] \arrow[d, equal] & S'' = \{s\}\, , \quad s \in E \ . \arrow[l] \\
E/G & E & \text{\footnotesize asph} \\
\text{\footnotesize asph} & \substack{\text{asph} \\ \text{si on peut}} &
\end{tikzcd}
\]
Cela donne encore un exemple avec $X$ $\D$-fibré propre et lisse, sur une base $S$ $\D$-asphérique, tel que l'inclusion d'une fibre $X_s$ dans $X$ ne soit pas une $\D$-équivalence. On peut même se débrouiller pour que $H^*(Z,k)^G \simeq k$ ce qui implique
\msp{page 94}
que $X$ est lui aussi asphérique (mais les fibres de $X \to S$ ne sont pas $\D$-asphériques). Cela donne donc un autre contre-exemple au b) de la prop. 6 (p. 84).

Voici la façon la plus économique, je crois, de réaliser la situation que j'ai dite. Topologiquement, je prend % [sic] je prend -> je prends, conjugation
\[
\begin{aligned}
G &= \{\pm 1\} \\
E &= S^2 &\qquad& \text{sphère euclidienne, sur laquelle } \{\pm 1\} \\
  &      &&      \text{opère par antipodisme } x \mapsto -x \\
B &= S^2/\{\pm 1\} && \text{plan projectif réel} \\
Z &= S^1 && \text{où G opère comme précédemment par la symétrie\,.}
\end{aligned}
\]
C'est donc un exemple très proche de celui du début (p. 87), où on pouvait aussi considérer la base ($B = S^1$) comme quotient d'un $E$ ($= S^1$) par $G = \{\pm 1\}$ opérant librement (par
\printedpage{58}
% head: A. Grothendieck : Dérivateurs, Ch. I, Édition des Universités de Montpellier II et Paris VII
l'antipodisme $z \mapsto -z$ dans $\mathbb{C} \simeq \RR^2$). Mais dans cet exemple, la base n'était pas $\D$-asphérique.

La façon la plus économique de réaliser la sphère $S^2$ par un ensemble ordonné est sans doute par sa décomposition ``globulaire'' par trois paires de cellules de dimension 0, 1, 2, donnant l'ensemble ordonné
\[
I \qquad = \qquad
\begin{tikzcd}[row sep=large, column sep=large]
S_2^- & S_2^+ \\
S_1^- \arrow[u] \arrow[ur] & S_1^+ \arrow[u] \arrow[ul] \\
S_0^- \arrow[u] \arrow[ur] & S_0^+ \arrow[u] \arrow[ul]
\end{tikzcd}
\]
et l'antipodisme échange $S_i^-$ et $S_i^+$, mais le quotient n'est pas ordonné, c'est la catégorie (importante pour la théorie des 2-catégories) :
\msp{page 95}
\[
\begin{tikzcd}[row sep=large]
S_2 \\
S_1 \arrow[u, shift left, "p_1^-"] \arrow[u, shift right, "p_1^+"'] \\
S_0 \arrow[u, shift left, "p_0^-"] \arrow[u, shift right, "p_0^+"']
\end{tikzcd}
\qquad\qquad \text{avec} \qquad\qquad
\begin{cases}
p_1^- p_0^- = p_1^- p_0^+ \\
p_1^+ p_0^- = p_1^+ p_0^+
\end{cases}
\]
Mais prenons la subdivision barycentrique de la décomposition cellulaire précédente de $S$ dont les sommets correspondent aux éléments de $\mathsf{Ob}\ I$, les simplexes aux parties totalement ordonnés de $\mathsf{Ob}\ I$. On trouve l'octaèdre, ayant 6 sommets, 12 arêtes, 8 faces, et $G$ opère librement sur l'ensemble des simplexes, donc le quotient : 3 sommets, 6 arêtes, 4 faces, décomposition cellulaire du plan projectif. (Sa décomposition par trois droites.) On présente $S^1$ par l'ensemble ordonné (comme au début) % [sic] parties totalement ordonnés -> parties totalement ordonnées, agreement
\[
\begin{tikzcd}[row sep=large, column sep=large]
S_1^- & S_1^+ \\
S_0^- \arrow[u] \arrow[ur] & S_0^+ \arrow[u] \arrow[ul]
\end{tikzcd}
\qquad\qquad \textit{i.e.} \qquad\qquad
\begin{tikzcd}
 & S_1^+ & \\
0 = S_0^- \arrow[ur] \arrow[dr] & & S_0^+ = 0' \arrow[ul] \arrow[dl] \\
 & S_1^- &
\end{tikzcd}
\]
et $G$ y opère par symétrie échangeant $S_1^-$ et $S_1^+$. Donc l'ensemble ordonné $X$ aura $4 \times 13 = 52$ sommets et c'est modeste \ldots

\dgrblock{dgrb60}{The finite poset and its residual top class}
\begin{dcomment}{development}{p95:Donc l'ensemble ordonné $X$ aura}{The finite poset and its residual top class}
The passage to finite posets deserves a complete computation. The number $52$ in the text is correct. The claimed acyclicity of the resulting total space is not. We give both verifications, retaining the printed construction and its error in the source. The transfer calculation used below is standard; see \citet[\S3.G]{Hatcher2002}. We did not find a published discussion of this particular $52$-object calculation in the sources checked.

Footnote 42 describes the later example as having disconnected fibres, whereas the explicit choice $Z=S^1$ on manuscript page 94 has connected fibres. The covering alternative of footnote 43 and a discrete regular $G$-set instead give disconnected fibres; the intended reference in footnote 42 cannot be fixed unambiguously.

\textbf{The corrected finite model.}\ours\
Let $E$ be the face poset of the boundary of an octahedron. Its vertices are $\{\pm1,\pm2,\pm3\}$; a nonempty face is a subset containing at most one element of each antipodal pair. There are $6+12+8=26$ faces. The involution $\tau$ negates every vertex. Let $Z$ be the four-object circle poset with two vertices $v_0,v_1$ below two edges $e_0,e_1$, and let $\tau$ fix the vertices and exchange the edges. Then $B=E/C_2$ and $X=(E\times Z)/C_2$ are posets with $13$ and $52$ objects. Their realizations satisfy
\[
 |NB|\simeq\mathbb{RP}^2,\qquad
 |NX|\simeq(S^2\times S^1)/((x,z)\sim(-x,\bar z)).
\]
For $2\in k^\times$, their cohomology is
\[
 H^*(B;k)=k,\qquad
 H^j(X;k)=\begin{cases}k,&j=0,3,\\0,&\text{otherwise}.\end{cases}
\]
The fibre is a circle, so its inclusion into $X$ is not a $k$-cohomology equivalence, but $X\to B$ is not a $k$-cohomology equivalence either.

\emph{Proof.} If nonempty faces $\sigma,\rho$ satisfy both $\sigma\subseteq\rho$ and $\sigma\subseteq-\rho$, then $\rho$ contains an antipodal pair, which is impossible. Hence for fixed representatives there is at most one group element transporting a possible upper face to a face containing the lower one. Quotient incidence therefore gives at most one morphism between two orbit objects. Ranks prevent nontrivial cycles. The same argument applied to the first coordinate proves the assertion for $E\times Z$. Every chain in either quotient lifts uniquely once its first object is chosen, so the nerve of the quotient is the quotient of the nerve. The nerve of $E$ is the barycentric subdivision of the octahedral sphere, and the nerve of $Z$ is a circle. Realization commutes with these finite products and these free simplicial quotients. This proves the assertions about realization and object counts.

For a finite regular covering with group $G$, pullback identifies cochains on the quotient with invariant cochains upstairs: every simplex has $|G|$ lifts, permuted freely and transitively. Averaging makes invariants exact when $|G|$ is invertible. Consequently the quotient cohomology is the invariant part of the covering-space cohomology. The antipodal map on $S^2$ acts by $-1$ on $H^2$, whereas reflection on $S^1$ acts by $-1$ on $H^1$. In $H^*(S^2\times S^1;k)$ the actions in degrees $0,1,2,3$ are respectively $+1,-1,-1,+1$. This gives exactly the displayed answer. The product of the two anti-invariant generators is invariant and survives in degree $3$. It is the term omitted by the proposed implication on manuscript pages 93--94. \hfill$\square$

There is an earlier incidence slip on the same page. For the six-object globular poset, label the two orbits of arrows between consecutive levels by the product of the endpoint signs. Composition multiplies these signs. The orbit relations are consequently
\[
 p_1^-p_0^-=p_1^+p_0^+,
 \qquad p_1^-p_0^+=p_1^+p_0^-.
\]
The displayed relations in the edition instead hold $p_1$ fixed while changing $p_0$; they do not describe this antipodal quotient. This does not affect the face-poset construction just proved, which gives an unambiguous replacement.

The broader recipe (a)--(c) also needs care. If $E$ is the discrete regular $G$-set, then $E/G$ is a point and $X=(E\times Z)/G\cong Z$. Taking a nontrivial action on $H^*(Z;k)$ satisfies the proposed invariant-cohomology conditions, yet the fibre inclusion is the identity. Thus those conditions alone do not force a counterexample. For the selected connected $E=S^2$ the fibre inclusion does fail, as calculated above, but invariants of the separate factors do not determine invariants of their tensor product.

As an independent check, the accompanying exact rational chain computation gives simplex counts $(13,36,24)$ for $NB$, with boundary ranks $(12,24)$, and $(52,340,576,288)$ for $NX$, with boundary ranks $(51,289,287)$. The resulting Betti numbers are $(1,0,0)$ and $(1,0,0,1)$. These finite calculations corroborate the proof; they are not used to infer the missing degree-three class.
\end{dcomment}


% ---- from Der-ChI-234_p061-065.tex
\printedpage{59}
% head: A. Grothendieck : Dérivateurs, Ch. I, Édition des Universités de Montpellier II et Paris VII
% foot: Version 8 juillet 2001

\msp{page 96}
Finalement, il semble qu’il ne reste plus rien de la proposition 2 (p.79) introduisant la notion de $\DD$-fibration – pas même l’équivalence des deux conditions $b$) et $b'$) ($\DD$-fibration à gauche universelle, $\DD$-fibration à droite universelle). On songe à leur relation à une condition $a$) modifiée ainsi (pour $f : X \to Y$) :

$a$) Pour tout $Y'' \to Y' \to Y$, avec $Y'$ et $Y''$ \Hot-\underline{asphériques} (et pas seulement $\DD$-asphériques) $X'' = X \times_Y Y'' \to X' = X \times_Y Y'$ est une $\DD$-équivalence. (Quand on suppose que $Y''$ et $Y'$ ont des éléments finaux, c’est la condition que $f$ soit une $\DD$-fibration à gauche universelle; quand on suppose que $Y''$ et $Y'$ ont des éléments initiaux, c’est la condition $\DD$-fibration à droite universelle.)\footnote[48]{Il faut expliciter le raisonnement. Il suffit de tester même sur $Y'' \to Y'$, où $Y''$ est \underline{ponctuelle}, et où $Y'$ a un objet final (resp. initial).}

On a évidemment $a$) $\Rightarrow$ $b$) et $a$) $\Rightarrow$ $b'$) et on aimerait que ce soient des équivalences. Voici un axiome intéressant\footnote[49]{\textit{cf.} \textbf{Der 3} $a'$) page suivante.} sur $\DD$, qui semble vérifié dans tous les cas connus, qui assure l’équivalence souhaitée. Je vais en même temps reformuler \textbf{Der 3} en trois volets (\textit{cf.} p. 67).

\msp{page 97}
\textbf{Der} 3 $a$) Si $Z \in \Dia$ est \Hot-asphérique, alors $p : Z \to e$ est une $\DD$-équivalence, \textit{i.e.} par exemple $p^* : \mathcal{A} = \DD(e) \to \DD(Z)$ est pleinement fidèle.

$b$) Soit $f : Z \to Z'$ une flèche dans \Dia. Si $f$ est une \Hot-équivalence, c’est une $\DD$-équivalence (\textit{i.e.} $f^* : \DD(Z') \to \DD(Z)$ induit une \underline{équivalence} entre les sous-catégories de $Z$, $Z'$ \texttt{[plutôt de} $\DD(Z)$\texttt{,} $\DD(Z')$\texttt{]} formées des objets constants).

$a'$) (Précise considérablement $a$). Si $Z \in \Dia$ est \Hot-asphérique, alors on a $a$), et l’image essentielle de $\DD(e) = \mathcal{A}$ dans $\DD(Z)$ est formée des objets \underline{localement} \underline{constants} de $Z$, \textit{i.e.} des objets tels que le foncteur
\[
Z^{\circ} \longrightarrow \mathcal{A} \qquad z \longmapsto F_z
\]
soit localement constant (\textit{i.e.} transforme toute flèche en isomorphisme) donc isomorphe à un foncteur constant (puisque $Z$ est \Hot-asphérique, donc 1-connexe).

\dgrblock{dgrb61}{Locally constant objects as a homotopy invariant}
\begin{dcomment}{later}{p97:transforme toute flèche en isomorphisme}{Locally constant objects as a homotopy invariant}
The additional axiom sought here became a theorem for the usual derivator framework. \citet[Corollary 3.8]{Cisinski2009LocallyConstant} proves that a weak equivalence of nerves induces an equivalence on locally constant objects for every derivator on small categories. The model-category statement is Theorem 1.3 of the same paper; Theorem 3.13 organizes locally constant objects into a derivator. Thus Der~3~b$'$ is available under those modern hypotheses, rather than remaining an unresolved extra condition. This does not retroactively prove the assertion from only the provisional axioms currently being tried in the manuscript.

The modern interpretation uses the fundamental $\infty$-groupoid $\Pi_\infty|NZ|$. For a coefficient $\infty$-category $\mathcal C$ presenting the homotopy theory, locally constant coherent diagrams are functors from this groupoid to $\mathcal C$, and the category denoted $\DD(Z)_{\mathrm{lc}}$ is their homotopy category. If $|NZ|$ is contractible, evaluation gives the desired equivalence with $\DD(e)$. Constancy of the underlying functor into $\DD(e)$ alone is weaker: it forgets higher coherence. This is the same distinction emphasized in the commentary at manuscript p.~7 (``Coherent squares and what the axioms remember''), now applied to local systems.
\end{dcomment}


Il faut bien noter que dans $a'$), il ne suffit \underline{pas} de supposer que $Z$ soit $\DD$-asphérique, au lieu de \Hot-asphérique. Dès que la $\DD$-équivalence est strictement plus faible que \Hot, il semble que $a'$) devient toujours faux, en tant que condition sur $\DD$, si on y suppose seulement que $Z$ soit $\DD$-asphérique. Par exemple, même si $\DD = \DD_i$ (théorie des $i$-types d’homotopie relative), la catégorie
\msp{page 98}
des objets localement constants sur $Z$ s’exprime en termes du $(i+1)$-type d’homotopie, mais pas en termes de son $i$-type d’homotopie seul.

\dgrblock{dgrb62}{Sets already detect the missing monodromy}
\begin{dcomment}{examples}{p98:mais pas en termes de son $i$-type d’homotopie seul.}{Sets already detect the missing monodromy}
The shift from an $i$-type to an $(i+1)$-type has a concrete beginning. For the represented derivator
$\DD_0(A)=\mathsf{Set}^{A^{\mathrm{op}}}$, natural maps between constant diagrams on a nonempty connected category are simply maps of sets. Hence $e\to BG$ is a $\DD_0$-equivalence for every group $G$. Locally constant set-valued diagrams on $BG$, however, are $G$-sets, and restriction to $e$ forgets the action. For a nontrivial $G$, this functor is not full: there are maps of sets from a one-point set to the regular $G$-set, but no equivariant such map. Thus a $0$-type records connected components, while local systems of $0$-types already record the fundamental group.

This also supplies a direct test of footnote 36. Take $G=C_2$ and the covering $f:E_G\to BG$. Its strict fibres are the discrete two-element set $G$. Since the covering is a discrete fibration and a discrete opfibration, every base change is a covering of the same kind. Right Kan extension of a constant set $M$ has fibre $M^G$, and left Kan extension has fibre $G\times M$; arrows in the base act by permutations of $G$. Both resulting diagrams are locally constant and both Kan extensions calculate on strict fibres. Therefore conditions (b) and (b$'$) hold universally for $\DD_0$.

Now pull back along $e\to BG$. The induced map of total categories is $G_{\mathrm{disc}}\to E_G$. It is not a $\DD_0$-equivalence: for a two-element coefficient set $M$, the constant-coefficient limit on the connected category $E_G$ is $M$, whereas on the discrete fibre it is $M^2$. The comparison is the diagonal map and is not bijective. Thus (b) does not imply (a), even for the derivator of $0$-types. The correct positive assertion for ordinary weak homotopy equivalences uses their invariance of locally constant coefficients, as in the preceding block; it does not replace them by $\DD_0$-equivalences. If $\Dia$ is restricted to finite posets, the connected double cover of the four-object circle gives the same counterexample. Its total poset has eight objects and its self-pullback is the disjoint union of two copies of that total poset. Both bases are connected, while the pulled-back total categories have different numbers of components.
\end{dcomment}


Il semble bien que ce soit l’axiome \textbf{Der} 3 ($a'$) qui assure l’existence d’une bonne notion (“bilatère”) de $\DD$-fibration, qui soit stable par changement de base, et qui se réduise soit
\printedpage{60}
% head: A. Grothendieck : Dérivateurs, Ch. I, Édition des Universités de Montpellier II et Paris VII
% foot: Version 8 juillet 2001
à “$\DD$-fibration à gauche universelle” soit à “$\DD$-fibration à droite universelle” et donc qui soit en plus une condition \underline{vérifiable} pratiquement. Dans le cas où $f$ est propre ou lisse, le critère est particulièrement simple, puisque il se borne à la condition que les $f_* f^* a_Y$, ou les $f_! f^* a_Y$ (pour $a \in \mathcal{A}$), soient des objets localements constants sur $Y$ (donc \underline{constants} si $Y$ est \underline{\underline{\Hot}}-asphérique, en vertu de \textbf{Der} 3 ($a'$). % [sic] localements -> localement, plural ending on the adverb; % [sic] "(donc ... ($a'$)." -> closing parenthesis after "($a'$)" missing, unbalanced parenthesis

Je vais également introduire une variante renforcée de \textbf{Der} 3 $b$) :

\textbf{Der} 3 $b'$) Soit $Z' \to Z$ une \Hot-équivalence dans \Dia. Alors $f^* : \DD(Z) \to \DD(Z')$ induit une équivalence entre les catégories des coefficients \underline{localements} \underline{constants} sur $Z$ et sur $Z'$.\footnote[50]{\textbf{Der} 3 $b'_0$) : le foncteur est fidèle\\
\hspace*{6em}$b'_1$) : le foncteur est pleinement fidèle\\
\hspace*{6em}$b'_2$) : le foncteur est une équivalence} % [sic] localements -> localement, plural ending on the adverb

\msp{page 99}
Dans certaines questions il suffit de savoir que le dit foncteur est pleinement fidèle. Cela implique que
\[
H^*_{\DD}(Z,F) \overset{\sim}{\longrightarrow} H^*_{\DD}(Z', f^*F)
\]
quand $F$ est \underline{localement} \underline{constant} sur $Z$\footnote[51]{On pourra dire que $f : Z' \to Z$ est une $\DD$-équivalence 0-\textit{complète à gauche,} et 0-\textit{complète à droite} dans le cas homologique.} (mais il n’est pas clair que cette condition implique toujours la pleine fidélité\footnote[52]{Mais c’est OK si on a des \underline{\Hom} internes dans $\DD(Z)_{\mathrm{loc.const.}}$ satisfaisant
\[
\Hom(F,G) \simeq \Hom(e_Z, \underline{\Hom}(F,G)) \quad ,
\]
où $e$ est l’élément final de $\mathcal{A}$, les \underline{\Hom} étant compatibles aux images inverses.}) et de même pour l’homologie
\[
H_*^{\DD}(Z', f^*F) \longrightarrow H_*^{\DD}(Z,F) \quad .
\]
Moyennant cette condition, on a la :

\textbf{Proposition :} \textit{On suppose} \textbf{Der} 3 \textup{b}$'$) \textit{satisfait. Soit $f : X \to Y$ dans} \Dia. \textit{Pour que $f$ soit une $\DD$-fibration, il faut et il suffit que pour $Y'' \to Y'$ sur $Y$ qui soit une} \Hot\textit{-équivalence, $X'' = X \times_Y Y'' \to X' = X \times_Y Y'$ soit une $\DD$-équivalence.}

À vrai dire, on peut énoncer quelque chose de raisonnable, sans supposer l’axiome \textbf{Der} 3 $b'$). On doit introduire la :

\textbf{Définition :} Soit $f : Z' \to Z$ dans \Dia. On dit que $f$ et une $\DD$-équivalence 1-\textit{complète} si le foncteur $f^* : \DD(Z) \to \DD(Z')$ est pleinement fidèle sur la sous-catégorie pleine de $\DD(Z)$ formée des coefficients localement constants. On dit que $f$ est une $\DD$-équivalence 2-\textit{complète} si elle induit une équivalence entre les catégories des coefficients localement constants. % [sic] f et une -> f est une, verb "est" misprinted as "et"

\msp{page 100}
On dit que $Z$ est \textit{complètement} $\DD$-\textit{asphérique}, si $Z \to e$ est une $\DD$-équivalence 2-complète.

\printedpage{61}
% head: A. Grothendieck : Dérivateurs, Ch. I, Édition des Universités de Montpellier II et Paris VII
% foot: Version 8 juillet 2001

Ainsi, l’axiome \textbf{Der} 3 $a'$) peut se reformuler, en disant que si $Z$ est \Hot-asphérique, alors $Z$ est complètement $\DD$-asphérique. Et de même, \textbf{Der} 3 $b'$) en disant que si $f : X \to Y$ dans \Dia{} est une \Hot-équivalence, alors $f$ est une $\DD$-équivalence 2-complète. Nous allons donner une variante affaiblie, de \textbf{Der} 3 $a'$), qui suffit pour donner une notion raisonnable de $\DD$-fibration :

\textbf{Der} 3 $a''$) Tout $Z$ dans \Dia{} qui a soit un objet final, soit un objet initial, est complètement $\DD$-asphérique.

\textbf{Définition.} On dit que $f : X \to Y$ dans \Dia{} est une $\DD$-fibration, si pour
\[
Y'' \longrightarrow Y' \longrightarrow Y
\]
avec $Y'' \to Y'$ une $\DD$-équivalence 1-complète (par exemple $Y'$ et $Y''$ complètement $\DD$-asphèriques), $X'' = X \times_Y Y'' \longrightarrow X' = X \times_Y Y'$ est une $\DD$-équivalence. % [sic] asphèriques -> asphériques, grave accent instead of acute

Cela équivaut à chacune des conditions $b$), $b'$) de la prop. 2 (page 79), \textit{i.e.} que $f$ soit universellement une $\DD$ fibration à gauche, ou universellement une $\DD$-fibration à droite (grâce à \textbf{Der 3} $b''$ \texttt{[plutôt} $a''$\texttt{]}). % [sic] $\DD$ fibration -> $\DD$-fibration, hyphen missing

\msp{page 101}
\textbf{Définition 1.} $f : X \to Y$ dans \Dia{} est appelé $\DD$-\textit{équivalence cohomologique} (sous-entendu : \textit{complète}) si pour tout $F \in \Ob\,\DD(Y)$, $F$ \underline{localement} \underline{constant}, on a
\[
H^*_{\DD}(Y,F) \overset{\sim}{\longrightarrow} H^*_{\DD}(X, f^*(F)) \quad .
\]
Notion duale $\DD$-\textit{équivalence homologique}
\[
H_*^{\DD}(X, f^*(F)) \overset{\sim}{\longrightarrow} H_*^{\DD}(Y,F) \quad .
\]
Si $Y$ est complètement $\DD$-asphérique, alors les coefficients localement constants sur $Y$ sont constants, et $f$ est une $\DD$-équivalence cohomologique (resp. homologique) si et seulement si c’est une $\DD$-équivalence \textit{i.e.} si $X$ est $\DD$-asphérique.\footnote[53]{L’axiome \textbf{Der} 3 $b'_1$) implique que toute \Hot-équivalence est une $\DD$-équivalence cohomologique.}

\textbf{Proposition 1} (Soritale). \textit{Soient $X \overset{f}{\to} Y \overset{g}{\to} Z$. Si $f$ et $g$ sont des $\DD$-équivalences cohomologiques, alors $gf$ aussi. Si $gf$ et $f$ sont des $\DD$-équivalences cohomologiques, alors $g$ aussi. Si $gf$ est une $\DD$-équivalence cohomologique, et si $g$ est une $\DD$-équivalence \underline{2-complète}} (donc tout $F$ localement constant sur $Y$ provient d’un $G$ localement constant sur $Z$), \textit{alors $f$ est une $\DD$-équivalence cohomologique.}

\dgrblock{dgrb63}{Complete equivalences repair the base change argument}
\begin{dcomment}{exposition}{p101:provient d’un $G$ localement constant sur $Z$}{Complete equivalences repair the base change argument}
The adjective ``complete'' records which coefficients are being tested. A $1$-complete equivalence is fully faithful on the category of locally constant objects. A $2$-complete equivalence is in addition essentially surjective there. These are levels of categorical information, not degrees of homology. In particular a $1$-complete equivalence $p:Y'\to Y$ compares both homology and cohomology with every locally constant coefficient $F$ on $Y$.

For cohomology, adjunction gives
\[
 \Hom_{\DD(e)}(M,H^*_{\DD}(Y,F))
       =\Hom_{\DD(Y)}(M_Y,F).
\]
Both arguments on the right are locally constant. Full faithfulness identifies this group of morphisms with $\Hom_{\DD(Y')}(M_{Y'},p^*F)$, naturally in $M$. Yoneda gives the asserted cohomology isomorphism. Applying the same argument to morphisms from $F$ to a constant object proves homological invariance. The internal-Hom qualification in footnote 52 is a converse criterion: if locally constant internal Homs exist and commute with pullback, cohomological invariance applied to $\underline{\Hom}(F,G)$ yields full faithfulness. Its unit must be the stated cartesian unit, or an explicitly chosen monoidal unit; the zero object of a stable theory cannot play that role merely by being terminal.

We can now complete the repaired proof. Suppose $f$ is universally cohomological, and let $p$ be $1$-complete. For a constant coefficient $M$, put $F=f_*M_X$. The coefficient $F$ is locally constant, and universal fibre calculation identifies $p^*F$ with $f'_*M_{X'}$. The preceding Yoneda argument compares the global cohomology of these two objects and proves that $X'\to X$ is a $\DD$-equivalence. This establishes the new pullback-preservation condition. Conversely, under Der~3~a$''$, every functor between categories having final objects is $2$-complete, because both locally constant categories are equivalent to $\DD(e)$ and the restriction agrees with the identity on constants. Testing the new condition on $\{y\}\to Y/y$ and on transitions between slices therefore gives universal cohomological fibre calculation. Categories with initial objects give the homological condition. These arguments prove both directions and explain why the weaker endpoint axiom suffices at this stage.

The cancellation statements on manuscript page 101 have the same logic. Composing two comparisons proves closure under composition; cancelling the comparison for $f$ proves the implication from $gf$ and $f$ to $g$. To recover $f$ from $gf$, one must first express an arbitrary locally constant coefficient on $Y$ as the pullback of a coefficient on $Z$. The $2$-complete hypothesis on $g$ supplies exactly that essential-surjectivity step, and its full faithfulness supplies the comparison for $g$. Ordinary two-out-of-three would have omitted the coefficient which must descend.
\end{dcomment}


\textbf{Définition 2.} 1°) $f : X \to Y$ dans \Dia{} est appelé une $\DD$-\textit{fibration cohomologique complète,} je dis bien \underline{complète}, si

$a$) Pour tout $y \in \Ob\, Y$, $X_y \to X/y$ est une $\DD$-équivalence cohomologique (complète).

$b$) Pour tout $y' \to y$ dans $Y$, $X/y' \to X/y$ est une $\DD$-équivalence cohomologique (complète).

\printedpage{62}
% head: A. Grothendieck : Dérivateurs, Ch. I, Édition des Universités de Montpellier II et Paris VII
% foot: Version 8 juillet 2001

\msp{page 102}
$c$) Ces conditions restent satisfaites après tout changement de base $Y' \to Y$.

2°) Notion duale de fibration homologique complète.

\textsc{Exemple.} Si $f$ est \Hot-propre, les flèches des $a$) et $b$) sont des \Hot-équivalences, donc en vertu de \textbf{Der 3} $b'_1$) (qu’il faut admettre), ce sont des $\DD$-équivalences cohomologiques. Cela signifie que $a$) et $b$) sont valables, et $c$) aussi, car la notion de flèche \Hot-propre est stable par changement de base.

\dgrblock{dgrb64}{What the completion requires}
\begin{dcomment}{exposition}{p102:car la notion de flèche \Hot-propre est stable par changement de base.}{What the completion requires}
The quantifier on the coefficients is the decisive change. A weak cohomological fibration is tested by the constant objects $M_X$. A complete cohomological fibration must work with every locally constant object of each category obtained by base change. Such an object need not descend from $X$, and need not become constant on $X$. The monodromy examples just discussed explain why that distinction cannot be suppressed. The role of completion is to make the coefficients produced by one direct image admissible for a second direct image.

We keep the variance of the manuscript: objects of $\D(X)$ behave as coherent contravariant diagrams, and the pointwise formula for a right direct image is
\[
 (f_*F)_y\simeq H^*_{\D}(X/y,F|_{X/y}).
\]
Write $\mathcal L(X)\subseteq\D(X)$ for the full subcategory of locally constant objects. Local constancy means that every transition morphism is invertible; it does not mean that an isomorphism to a constant object has been chosen. For $r:A\to B$, the cohomological condition is the isomorphism
\[
 H^*_{\D}(B,L)\longrightarrow H^*_{\D}(A,r^*L)
 \qquad(L\in\mathcal L(B)).
\]
The conditions on $X_y\to X/y$ therefore compare a true fibre with a comma category for every local system on the latter. The conditions on $X/y\to X/y'$ then control the transport of those cohomology objects along the base. Universal repetition of the first comparison is what makes the resulting fibre formula compatible with arbitrary base changes.

The example immediately preceding Proposition~2 contains a genuine slip: $\Hot$-properness gives the fibre-to-comma comparison, but does not make the comma categories locally constant as $y$ varies. Take the inclusion $\{1\}\hookrightarrow[1]$, where $[1]=\{0\to1\}$. This is a closed immersion and hence proper, but its comma categories at $0$ and $1$ are respectively empty and terminal. The transition $\varnothing\to e$ is not a weak homotopy equivalence. Thus the example becomes correct after adding the transport condition, or by assuming that $f$ is both $\Hot$-proper and a weak $\HOT$-fibration. This is precisely the hypothesis separated out on manuscript page~107 and proved usable in Proposition~6 on page~118. The printed assertion is retained in the source and is audited as false as written.
\end{dcomment}


\textbf{Proposition 2.} \textit{Soient $f : X \to Y$ $\DD$-fibration cohomologique complète, $F \in \Ob\,\DD(X)$ localement constant sur $X$. Alors $f_*(F)$ est localement constant, sa formation commute à tout changement de base, et en particulier se calcule fibre par fibre.}\footnote[54]{Il y a une réciproque \textit{cf.} prop. 3$'$.}

\textbf{Corollaire 1.} \textit{Soit un diagramme cartésien}
\[
\begin{tikzcd}
X & X' \arrow[l, "q"'] \arrow[d, "f'"] \\
Y & Y' \arrow[l, "p"']
\arrow[from=1-1, to=2-1, "f"']
\end{tikzcd}
\quad .
\]
\textit{On suppose $f$ $\DD$-fibration cohomologique et $p$ une $\DD$-équivalence cohomologique. Alors $q$ est une $\DD$-équivalence cohomologique.}

\msp{page 103}
En effet, on a pour $F$ localement constant sur $X$
\[
H^*_{\DD}(X,F) \simeq H^*_{\DD}(Y, f_*F) \overset{\sim}{\longrightarrow} H^*_{\DD}(Y', p^*f_*F) \simeq H^*_{\DD}(Y', f'_*q^*F) \simeq H^*_{\DD}(X', q^*F)
\]
qed.

\textbf{Proposition 3.} \textit{La notion de $\DD$-fibration cohomologique complète est stable par changement de base} (trivial) \textit{et par composition.}

Pour la stabilité par composition, il suffit de noter le corollaire suivant de la proposition 2.

\textbf{Corollaire 2} (de prop. 2). \textit{Soit $f : X \to Y$. Pour que $f$ soit une $\DD$-fibration cohomologique complète, il faut et il suffit que pour tout $Y'' \overset{p}{\to} Y' \to Y$ sur $Y$, avec $p$ une $\DD$-équivalence cohomologique, la flèche correspondante $X'' = X \times_Y Y'' \longrightarrow X' = X \times_Y Y'$ soit aussi une $\DD$-équivalence cohomologique.}

C’est nécessaire par le corollaire 1, compte tenu que $f' : X' = X \times_Y Y' \to Y'$ est aussi une $\DD$-fibration cohomologique. C’est suffisant, car il suffit de l’appliquer aux cas $\{y\} \to Y/\{y\}$, $Y/y \to Y/y'$, qui sont des $\DD$-équivalences cohomologiques en
\msp{page 104}
vertu de \textbf{Der} 3 $b'_1$) (ou seulement \textbf{Der} 3 $a'$)).

\printedpage{63}
% head: A. Grothendieck : Dérivateurs, Ch. I, Édition des Universités de Montpellier II et Paris VII
% foot: Version 8 juillet 2001

\textbf{Proposition 4.} \textit{On suppose que} \textbf{Der} 3 a$'$) \textit{est satisfait} (\textit{i.e.} tout $Z$ dans \Dia{} qui est \Hot-asphérique, est complètement $\DD$-asphérique). \textit{Soit $f : X \to Y$ dans} \Dia. \textit{Les conditions suivantes sont équivalentes}

1°) \textit{$f$ est une $\DD$-fibration cohomologique \underline{complète}} (conditions $a$), $b$), $c$) de \texttt{[la]} définition 2, page 101).

2°) \textit{Pour tout $F \in$} \Ob\,$\DD(X)$\textit{, localement constant, $f_*(F)$ est localement constant, se calcule fibre par fibre; et ça reste vrai après tout changement de base $Y' \to Y$ (avec $F'$ localement constant sur $X' = X \times_Y Y'$). (Alors la formation de $f_*(F)$ commute à \underline{tout} changement de base).}

3°) \textit{Pour tout $Y'' \overset{p}{\to} Y'$ sur $Y$, si $p$ est une $\DD$-équivalence cohomologique, alors $X'' = X \times_Y Y'' \longrightarrow X' = X \times_Y Y'$ aussi.}

3$'$) \textit{Comme 3, en supposant que $Y''$ est la catégorie ponctuelle, et que $Y'$ a un élément final : il faut que $X'' \to X'$ soit une $\DD$-équivalence cohomologique.}

\dgrblock{dgrb65}{The universal criterion, with the quantifiers restored}
\begin{dcomment}{exposition}{p104:la catégorie ponctuelle, et que $Y'$ a un élément final}{The universal criterion, with the quantifiers restored}
Proposition~4 is more than a convenient restatement: it supplies the missing reason that the strengthened notion composes. We give the argument, keeping the hypothesis that every weakly contractible indexing category has only constant local systems, as in \textbf{Der}~3~$a'$. The same argument will apply to Proposition~9, after the terminology is revised.

First assume the comma-category conditions of Definition~2, universally in the base. If $F\in\mathcal L(X)$, its restriction to $X/y$ is locally constant. The comparison $X_y\to X/y$ gives
\[
 (f_*F)_y\simeq H^*_{\D}(X_y,F|_{X_y}).
\]
For an arrow $y\to y'$, the functor $X/y\to X/y'$ is a cohomological equivalence, so the resulting transition $(f_*F)_{y'}\to(f_*F)_y$ is an isomorphism. Thus $f_*F$ is locally constant. In a pullback square with base map $p:Y'\to Y$ and induced maps $f',q$, evaluate the base-change map
\[
 p^*f_*F\longrightarrow f'_*q^*F
\]
at $y'$. Both sides identify with the cohomology of the same strict fibre $X_{p(y')}$ with the same restricted coefficient. The identification of the comparison with the identity is the naturality of the pointwise Kan-extension formula. Pointwise detection of isomorphisms proves base change.

Suppose now that this direct-image property holds after every base change. For a cohomological equivalence $p:Y''\to Y'$, and $L\in\mathcal L(X')$, put $G=f'_*L$. The preceding property makes $G$ locally constant and identifies $p^*G$ with $f''_*q^*L$. The canonical comparison is consequently the composite
\[
 H^*_{\D}(X',L)\simeq H^*_{\D}(Y',G)
 \longrightarrow H^*_{\D}(Y'',p^*G)
 \simeq H^*_{\D}(X'',q^*L),
\]
whose middle map is invertible. This proves the universal preservation criterion. Notice that it tests all $L$ on $X'$, rather than just coefficients pulled back from $X$.

It remains to recover the comma conditions from the apparently much smaller test in~$3'$. Every functor $e\to B$ with $B$ admitting a terminal object is a cohomological equivalence under $a'$: all local systems on $B$ are constant and $B$ is $\D$-aspherical. Apply the test to the base change $Y/y\to Y$, selecting the terminal object $1_y$. Its inverse image is $X_y$, and the total space is $X/y$. This gives the first comma comparison. For $u:y\to y'$, use $Y/y'\to Y$ and select its object $u$. We get a cohomological equivalence $X_y\to X/y'$. It factors as
\[
 X_y\longrightarrow X/y\longrightarrow X/y'.
\]
The first map and the composite are cohomological equivalences. Testing a local system on $X/y'$ in the two comparisons shows that the second map is also one. This is the valid cancellation direction from Proposition~1 on manuscript page~101; no unqualified two-out-of-three assertion for cohomological equivalences is being used. As the point test was required over every base change, all the comparisons are universal.

These arguments also justify the parenthesis in Proposition~9 that only the fibre comparison needs to be imposed universally. Applied to $Y/y'$ and to its arbitrary object $u$, that one comparison already recovers transport. They explain why the converse concerns arbitrary locally constant coefficients on the pullbacks. Merely proving the fibre formula for restrictions of coefficients from the original $X$ would leave a real gap.
\end{dcomment}


\msp{page 105}
\textbf{Définition 3.} On dit que $f : X \to Y$ dans \Dia{} est une $\DD$-\textit{fibration complète} si c’est une $\DD$-fibration cohomologique complète et une $\DD$-fibration homologique complète.

C’est stable par changement de base, par composition.

La définition 3 équivaut à ceci : pour tout $p : Y'' \to Y'$ sur $Y$, avec $Y''$ et $Y'$ \Hot-asphériques, le morphisme correspondant $X'' \to X'$ est \underline{à} \underline{la} \underline{fois} une $\DD$-équivalence cohomologique, et une $\DD$-équivalence homologique. (Et il suffit même de le supposer pour $Y'' = e$ et pour $Y'$ avec objet final ou objet initial.) On peut songer à renforcer encore cette condition en exigeant (sans regarder à gauche ni à droite) que $X'' \to X'$ soit une $\DD$-équivalence 1-complète. Mais il n’est pas clair que cette notion renforcée soit stable par composition – à moins de la renforcer encore, en exigeant que si $Y'' \to Y'$ est une $\DD$-équivalence 1-complète, alors $X'' \to X'$ l’est aussi. L’ennui avec cette définition renforcée, c’est que faute d’avoir une description cohomologique ou homologique de la notion de
\msp{page 106}
$\DD$-équivalence 1-complète, je ne me sens pas capable de donner des exemples intéressants (non totalement triviaux) de $\DD$-fibrations en ce sens très fort. Sauf les \Hot-fibrations triviales (\textit{i.e.} à fibres \Hot-asphériques) – mais ce n’est pas drôle!

Mais quand on a des produits $*$ associatifs et unitaires et des \underline{\Hom} internes correspondants dans les $\DD(X)_{\mathrm{loc.\ const.}}$, compatibles avec \texttt{[les]} images inverses, donc on aura sur $X$ si $F$ et $G$ sont localement constants dans $\DD(X)$, et si $\varepsilon_X$ désigne l’unité pour $*$ dans $\DD(X)$,
\[
\Hom(F,G) \simeq \Hom(\varepsilon_X * F, G) \simeq \Hom(\varepsilon_X, \underline{\Hom}(F,G)) \simeq \Hom_{\mathcal{A}}(\varepsilon_e, p_*\, \underline{\Hom}(F,G)) \ ,
\]
car $\varepsilon_X = p^*(\varepsilon_e)$, où $p : X \to e$, donc dans ce cas une flèche $f : X \to Y$ dans \Dia{} est une $\DD$-équivalence 1-complète si (et seulement si) c’est une $\DD$-équivalence cohomologique. Donc dans ce cas, si $f$ est une $\DD$-fibration cohomologique complète, c’est aussi une $\DD$-fibration homologique complète, donc aussi une $\DD$-fibration complète (définition 3). (Et même au sens soit-disant renforcé de tantôt : si $Y'' \to Y'$ sur $Y$ est une $\DD$-équivalence 1-complète ($\Leftrightarrow$ une $\DD$-équivalence cohomologique!) alors $X'' \to X'$ aussi. % [sic] soit-disant -> soi-disant, spelling; % [sic] "(Et même ... aussi." -> closing parenthesis missing, unbalanced parenthesis

\dgrblock{dgrb66}{What internal Hom does, and what it does not do}
\begin{dcomment}{exposition}{p106:alors $X'' \to X'$ aussi.}{What internal Hom does, and what it does not do}
The internal-Hom argument is sound, provided its compatibility hypotheses are stated at the level where they are used. Suppose that each $\mathcal L(X)$ is a closed monoidal category, that its tensor unit is $\varepsilon_X=p_X^*\varepsilon_e$, and that restriction preserves the internal Hom between locally constant objects. For $F,G\in\mathcal L(Y)$, adjunction gives
\[
 \Hom_{\D(Y)}(F,G)
 \simeq\Hom_{\D(e)}\bigl(\varepsilon_e,
                p_{Y*}\underline{\Hom}_Y(F,G)\bigr).
\]
If $f:X\to Y$ is a cohomological equivalence, apply its defining comparison to the locally constant object $\underline{\Hom}_Y(F,G)$. The assumed comparison
\[
 f^*\underline{\Hom}_Y(F,G)
 \simeq\underline{\Hom}_X(f^*F,f^*G)
\]
then identifies the displayed Hom set with $\Hom_{\D(X)}(f^*F,f^*G)$. Thus $f^*$ is fully faithful on $\mathcal L(Y)$: $f$ is 1-complete.

Conversely, full faithfulness on local systems compares the morphisms from every constant object $M_Y$ to $F\in\mathcal L(Y)$. By adjunction these are the morphisms from $M$ to $p_{Y*}F$. Yoneda gives the cohomological comparison. Comparing morphisms from $F$ to every $M_Y$, instead, gives the homological comparison. This last argument needs no tensor product. It proves the two outgoing implications from 1-completeness in the diagram on page~113, while the internal-Hom hypothesis supplies the converse only on the cohomological side.

There is no need to assume that the unit detects isomorphisms. In the forward direction the cohomological comparison is already an isomorphism before $\Hom(\varepsilon_e,-)$ is applied; in the reverse direction Yoneda uses every $M$. Nor does the argument identify 1-completeness with 2-completeness: essential surjectivity of local systems is an additional demand.

One must also retain the adjective ``cohomological''. On page~122 the printed sentence says that an unqualified $\D$-equivalence is already 1-complete in the presence of internal Homs. That stronger assertion is false. For example, with $\D=\mathsf{Der}(\mathbb Q\text{-}\mathsf{Mod})$ and $G=C_2$, the map $e\to BG$ is a constant-coefficient equivalence, but restriction sends the sign representation to a one-dimensional vector space. There are no $G$-equivariant maps from the trivial representation to the sign representation, whereas after restriction there are nonzero maps. The categories of local systems have the usual derived tensor and internal Hom. The wording on pages~106 and~113, and the diagram on page~123, supply the correct repair: cohomological equivalence, not constant-coefficient equivalence.
\end{dcomment}


% ---- from Der-ChI-234_p066-070.tex
\printedpage{64}
% head: A. Grothendieck : Dérivateurs, Ch. I, Édition des Universités de Montpellier II et Paris VII
% foot: Version 8 juillet 2001

\msp{page 107}

Il se pose la question de trouver des $f$ qui soient des $\DD$-fibrations cohomologiques, ou homologiques. Je ne sais si je serai capable d'en trouver, qui ne soient déjà des \HOT-fibrations cohomologiques, ou homologiques (cela revient maintenant au même\footnote[55]{Les \HOT-équivalences cohomologiques $=$ \HOT-équivalences homologiques $=$ \HOT-équivalences sans plus.} ), car pratiquement je ne sais guère dire qu'une flèche $Z' \to Z$ dans $\Cat$ est une $\DD$-équivalence cohomologique complète, que si cette flèche est une \HOT-équivalence (et en admettant l'axiome \textbf{Der} 3 $b'_1$), ou sa variante affaiblie cohomologique).

\dgrblock{dgrb67}{A complete fibration beyond ordinary homotopy}
\begin{dcomment}{examples}{p107:Je ne sais si je serai capable d'en trouver, qui ne soient déjà des}{A complete fibration beyond ordinary homotopy}
The request for examples that are not already $\HOT$-fibrations has a simple answer with rational coefficients. Let $G=C_2$, let $BG$ denote its one-object category, and take the Grothendieck opfibration
\[
 f:E=\int_{[1]}(BG\longrightarrow e)\longrightarrow[1].
\]
Explicitly, $E$ has objects $a$ and $t$, with $\operatorname{End}(a)=G$, $\operatorname{End}(t)=\{1\}$, one arrow $a\to t$, and no arrow $t\to a$. The unique arrow to $t$ absorbs every $g\in G$. The fibres are $BG$ and $e$, and the transport is the collapse $r:BG\to e$.

For $\D=\mathsf{Der}(\mathbb Q\text{-}\mathsf{Mod})$, the functor $r$ is a cohomological equivalence. All coefficients on its target are constant, and invariants for a finite group over $\mathbb Q$ are exact. Indeed, averaging $v\mapsto |G|^{-1}\sum_{g\in G}gv$ is a natural projection onto invariants, and averaging a lift proves that invariants preserve epimorphisms. Consequently their right derived functor on any complex is computed degreewise. On a trivial $G$-complex $M$, this gives $R\Gamma(BG,M)=M$.

More generally, an opfibration whose transport functors are cohomological equivalences is a complete cohomological fibration. Properness supplies the fibre formula for every coefficient. For a locally constant coefficient $L$ on the total category and an arrow in the base, the cocartesian arrows give an isomorphism between $L$ on the source fibre and the pullback of $L$ on the target fibre. The cohomological equivalence of transport therefore makes the direct-image transition invertible. This argument survives an arbitrary pullback, since the new transport functors are the old ones indexed by the new arrows. Proposition~4, or Proposition~9 in the revised terminology, proves the assertion.

Applying this observation to $f$ shows that it is a complete cohomological fibration. Derived tensor products and internal Homs on rational local systems satisfy the hypotheses explained at page~106, so it is a complete homological fibration as well. It is not a $\HOT$-fibration: its transport $BG\to e$ is not a weak homotopy equivalence, since $\pi_1(BG)=C_2$. Thus stronger coefficient invariance is still relative to the chosen coefficients. Reversing the transport to $e\to BG$ destroys cohomological equivalence, and is exactly the feature exploited by the composition counterexample at manuscript page~118.
\end{dcomment}


Je vais donc chercher les \HOT-fibrations, \textit{i.e.} les $f : X \to Y$ tels que pour tout $Y'' \to Y' \to Y$, avec $Y''$ et $Y'$ asphériques (sous-entendu : pour \HOT), $X'' \to X'$ est une \Hot-équivalence. (Il suffit même de supposer $Y''$ réduit à un point, $Y'$ avec objet final $\ldots$ )

Supposons par exemple $f$ propre (\textit{i.e.} \HOT-propre), donc la condition $a$) de la définition 2, p. 101 satisfaite et de même après tout changement de base. Suffit-il de vérifier de plus la condition $b$), \textit{i.e.} que
\[ X/y \longrightarrow X/y' \quad , \qquad \text{pour } y \to y' \text{ dans } Y \ , \]
sont des \Hot-équivalences? (J'entends en restant sur $Y$, sans supposer aussi d'avoir la même chose après changement de base).
\msp{page 108}
Mais cela reste-il vrai après tout changement de base? % [sic] reste-il -> reste-t-il, forme interrogative

La situation semble ici meilleur pour $f$ lisse (au lieu de propre), en interprétant les conditions $a$) et $b$) sous forme homologique, comme signifiant que $f_!(e_X)$ est un objet \underline{localement} \underline{constant} de $\mathsf{HOT}(Y)$, qui se calcule fibre par fibre. Comme $f$ lisse implique que $f_!$ commute à tout changement de base, cette hypothèse sur $f$ (que $f$ soit une \HOT-fibration cohomologique, peut-être pas complète) est stable par changement de base, donc implique que c'est une \HOT-fibration cohomologique complète. Mais cet argument utilise les propriétés générales des dérivateurs pour \HOT, que je n'ai pas encore établies - surtout pas dans le domaine $\Cat$ tout entier! Donc il est sain de me méfier. Je préfère tâter la difficulté en restant d'abord dans le cadre du formalisme \underline{cohomologique} (non commutatif), dont j'ai une certaine maîtrise. Je suppose donc $f$ propre à nouveau, et que les % [sic] meilleur -> meilleure, accord avec « situation »
\[ X_y \longrightarrow X/y \quad , \qquad X/y \longrightarrow X/y' \]
\msp{page 109}
soient des \HOT-équivalences, je veux prouver la même chose après tout changement de base, au moins que dans la situation
\[
\begin{tikzcd}
X \arrow[r, leftarrow, "q"] \arrow[d, "f"'] & X' \arrow[r, leftarrow, "q'"] \arrow[d, "f'"'] & X'' \arrow[d, "f''"] \\
Y \arrow[r, leftarrow, "p"'] & Y' \arrow[r, leftarrow, "p'"'] & Y''
\end{tikzcd}
\]
\printedpage{65}
% head: A. Grothendieck : Dérivateurs, Ch. I, Édition des Universités de Montpellier II et Paris VII
% foot: Version 8 juillet 2001
si $Y'$, $Y''$ sont asphériques, alors $q'$ est une \Hot-équivalence. Si je veux l'établir par critère cohomologique, je dois commencer par prendre sur $X'$ un faisceau d'\textit{ensembles} localement constant $F$, et prouver que $\Gamma(X',F) \overset{\sim}{\to} \Gamma(X'',{q'}^{*}F)$. Comme ce faisceau ne provient pas de $X$, je suis ennuyé pour appliquer l'hypothèse sur $f$, qui concerne à priori des faisceaux d'ensembles constants sur $X$, ou sur les $X/y$. Mais on devrait avoir ceci :

\textbf{Lemme.} \textit{Soit $f : X'' \to X'$. Alors $f$ est une} \HOT-\textit{équivalence} i.e.
\[ H^{*}_{\mathsf{HOT}}(X', a_{X'}) \xrightarrow{\ \sim\ } H^{*}_{\mathsf{HOT}}(X'', a_{X''}) \]
\textit{est} \texttt{[un]} \textit{isomorphisme pour tout $a \in \mathsf{Hot} = \mathsf{HOT}(e)$, si et seulement si $f$ est une} \Hot-\textit{équivalence,}\footnote[56]{Attention ici \texttt{[mot illisible]} la distinction entre \HOT-équivalence et \Hot-équivalence.} \textit{et alors $f$ est même une} \HOT-\textit{équivalence} $2$-\textit{complète} (j'ai besoin seulement de : \HOT-équivalence cohomologique et homologique complète) i.e. \textit{pour tout $\xi \in \DD(X')$ \underline{localement} \underline{constant} sur $X'$, on a}\footnote[57]{Est-ce vrai aussi pour les dérivateurs $\DD_i$?}
\[ H^{*}_{\mathsf{HOT}}(X',\xi) \xrightarrow{\ \sim\ } H^{*}_{\mathsf{HOT}}(X'',{q'}^{*}(\xi)) \quad . \]
\msp{page 110}

La première assertion doit résulter des formules
\[ H^{*}_{\mathsf{HOT}}(Z,a_Z) \simeq \underline{\mathsf{Hom}}_{\mathsf{Hot}}(\mathsf{hot}(Z),a) \tag*{$(*)$} \]
\[ \text{(et aussi} \qquad\qquad H_{*}^{\mathsf{HOT}}(Z,a_Z) \simeq \mathsf{hot}(Z) \times a \qquad\qquad ) \]
donc si
\[ \mathsf{hot}(X'') \xrightarrow{\ \sim\ } \mathsf{hot}(X') \tag*{$(**)$} \]
est \texttt{[un]} isomorphisme, alors
\[ H^{*}_{\mathsf{HOT}}(X', a_{X'}) \longrightarrow H^{*}_{\mathsf{HOT}}(X'', a_{X''}) \]
est \texttt{[un]} isomorphisme, puisque c'est déduit de $(*)$ en prenant $\underline{\mathsf{Hom}}_{\mathsf{Hot}}(\,.\,,a)$. Inversement, si $(*)$ est toujours \texttt{[un]} isomorphisme, exprimant les deux membres \texttt{[mot illisible]} de $\underline{\mathsf{Hom}}_{\mathsf{Hot}}(\,.\,,a)$ et prenant les $\pi_0$, on trouve que
\[ \mathsf{Hom}_{\mathsf{Hot}}(\mathsf{hot}(X'),a) \xrightarrow{\ \sim\ } \mathsf{Hom}_{\mathsf{Hot}}(\mathsf{hot}(X''),a) \]
est un isomorphisme pour tout $a$, donc $(**)$ est un isomorphisme.

Mais le fait que ça implique que $f''$ soit une \HOT-équivalence cohomologique complète, \textit{i.e.} que la relation $(*)$ reste vraie pour un $\xi \in \DD(X')$ \underline{localement} \underline{constant} sur $X'$, me paraît une propriété délicate du dérivateur \HOT, sur laquelle il me faudra revenir.

Donc le lemme est loin d'être établi, faute d'avoir seulement construit le dérivateur \HOT, sans même parler d'en établir les propriétés spéciales délicates! Mais je n'ai aucune raison de douter du lemme, il fait partie de la kyrielle des intuitions de base que je voudrais fonder solidement.

\dgrblock{dgrb68}{The lemma for spaces}
\begin{dcomment}{later}{p110:Donc le lemme est loin d'être établi, faute d'avoir seulement construit le dérivateur}{The lemma for spaces}
The lemma on pages~109--110 is proved by the modern theory of locally constant diagrams. We can separate its two assertions, which need different arguments. Put $BX=|NX|$, and let $\mathcal S$ denote the homotopy theory of spaces. In the homotopy derivator $\HOT$, the constant-diagram formulas are
\[
 H^*_{\HOT}(X,M_X)\simeq\operatorname{Map}(BX,M),
 \qquad
 H_*^{\HOT}(X,M_X)\simeq BX\times M.
\]
These are derived mapping spaces and derived products, viewed in $\Hot$. Thus a weak equivalence $BX\to BY$ induces both constant-coefficient comparisons. Conversely, if all the cohomological comparisons are isomorphisms, their sets of components give bijections
\[
 [BY,M]\longrightarrow[BX,M]
\]
for every homotopy type $M$. Yoneda in $\Hot$ makes $BX\to BY$ an isomorphism. Homologically one can simply take $M=e$. This proves the identification of the manuscript's uppercase $\HOT$-equivalences with lowercase $\Hot$-equivalences, without conflating the homotopy derivator with its value at a point.

The local-coefficient assertion has another content: all coherence in a locally constant diagram must be transportable along a weak homotopy equivalence of the base. An invariant modern formulation identifies locally constant space-valued diagrams on $X$ with space-valued functors on the homotopy groupoid $\Pi_\infty(BX)$. Restriction along an equivalence of these groupoids is an equivalence of functor categories, and remains so on passing to homotopy categories. Hence a $\Hot$-equivalence is 2-complete for $\HOT$. It preserves cohomology and homology with every local coefficient by the adjunction argument in the preceding block.

A precise derivator reference is \citet[Theorem~1.3 and Corollary~3.8]{Cisinski2009LocallyConstant}: a nerve weak equivalence induces an equivalence of categories of locally constant objects, for model categories and, in the latter result, for arbitrary Grothendieck derivators on small categories. The construction of those local systems as a derivator, with both adjoints to its inclusion into the diagram derivator, is Theorem~3.13. Thus the strongest $b$-axiom proposed below is available in the usual modern setting; it is not a currently unproved property of spaces.

The higher-categorical description is compatible with straightening and unstraightening \citep[Theorem~2.2.1.2]{Lurie2009HTT}; the homotopy limit of a space-valued local system is its space of sections \citep[Corollary~3.3.3.4]{Lurie2009HTT}. This gives a geometric reading of the two comparisons: transport the family, then compare its sections or its total homotopy colimit. It also indicates exactly what is absent from the ordinary diagram in the homotopy category: the coherence of transport around paths and higher simplices.

Footnote~57 requires a qualification. If its question asks whether a genuine $\Hot$-equivalence induces 2-completeness for the derivators of truncated homotopy types, the answer is yes, by the same invariance theorem. If it asks whether a $\D_i$-equivalence alone has the same consequence, the answer is no already for $i=0$. The example in ``Sets already detect the missing monodromy'' at manuscript page~98 has $e\to BC_2$ a constant-coefficient equivalence for sets, while restriction of $C_2$-sets is not full: a map from the one-point trivial action to the regular action exists after forgetting the action and does not exist before. Truncating the coefficients can lose information about the base that the category of local systems still retains.
\end{dcomment}


\printedpage{66}
% head: A. Grothendieck : Dérivateurs, Ch. I, Édition des Universités de Montpellier II et Paris VII
% foot: Version 8 juillet 2001

\msp{page 111}

Le lemme implique qu'une $\DD = \mathsf{HOT}$-fibration (au sens de la définition 2 page 81, revue et corrigée à la page 96-97) est en fait déjà une \HOT-fibration cohomologique et homologique complète. Or on a vu que si $f$ est \underline{propre}, $f$ est une $\DD$-fibration si et seulement si les conditions précédentes \underline{sur $Y$} (sans changement de base) sont satisfaites. Et dualement pour $f$ lisse.

\begin{flushright}
11.11.90
\end{flushright}

Je sens le besoin de revoir les notions de $\DD$-fibrations en tous genres (p. 76 ff). D'abord les diverses variantes utiles de l'axiome \textbf{Der} 3 sur les dérivateurs :

\textbf{Der} 3 $a$)\footnote[58]{\textit{cf.} aussi la série \textbf{Der 3} $c \ldots$ , page 124 ff.} Si $Z \in \mathsf{Dia}$ est \Hot-asphérique, il est $\DD$-asphérique \textit{i.e.} $\mathcal{A} = \DD(e) \to \DD(Z)$ est pleinement fidèle. % Doubt resolved: attached PDF leaf 30 (printed p.66), footnote 58 has bold Der 3.

$a'$) Comme $a$), avec en plus la caractérisation de l'image essentielle de $\mathcal{A}$ dans $\DD(Z)$, comme formée des $\DD$-coefficients localements constants. % [sic] localements -> localement, adverbe invariable

\textbf{Der} 3 $b$) Si $f : Z' \to Z$ dans \Dia{} est une \Hot-équivalence, alors c'est une $\DD$-équivalence, \textit{i.e.} le foncteur $f^{*} : \DD(Z) \to \DD(Z')$ induit une équivalence entre les sous-catégories formées des coefficients constants (ou encore un foncteur pleinement fidèle $\DD(Z)_{\text{loc. const.}} \longrightarrow \DD(Z')$).

$b'_1$) Si $f$ comme ci-dessus, alors $f$ est une $\DD$-équivalence 1-complète, \textit{i.e.} le foncteur $f^{*} : \DD(Z) \to \DD(Z')$ induit un foncteur pleinement fidèle
\msp{page 112}
entre les sous-catégories pleines formées des objets localement constants.

$b'_2$) Comme $b'_1$), mais avec $\DD$-équivalence 2-\underline{complète} au lieu de 1-complète, \textit{i.e.} $f^{*}$ induit une équivalence entre les sous-catégories formées des coefficients localement constants.

$b''_{\mathrm{coh}}$) Si $f : Z' \to Z$ dans \Dia{} est une \Hot-équivalence, alors c'est une $\DD$-équivalence cohomologique complète, \textit{i.e.} pour tout $F \in \mathsf{Ob}\ \DD(Z)$ qui est \underline{localement} \underline{constant}, on a
\[ H^{*}_{\DD}(Z,F) \xrightarrow{\ \sim\ } H^{*}_{\DD}(Z',f^{*}F) \quad . \]

$b''_{\mathrm{hom}}$) Si $f : Z' \to Z$ dans \Dia{} est une \Hot-équivalence, c'est une $\DD$-équivalence homologique complète, \textit{i.e.} pour tout $F \in \mathsf{Ob}\ \DD(Z)$ qui est \underline{localement} \underline{constant}, on a\footnote[59]{\textit{cf.} variante $\pm$ spéciale pour \HOT{} pages 124-125.}
\[ H_{*}^{\DD}(Z',f^{*}F) \xrightarrow{\ \sim\ } H_{*}^{\DD}(Z,F) \quad . \]

\printedpage{67}
% head: A. Grothendieck : Dérivateurs, Ch. I, Édition des Universités de Montpellier II et Paris VII
% foot: Version 8 juillet 2001

On a le diagramme d'implications tautologiques pour les variantes de \textbf{Der} 3 :
\[
\begin{tikzcd}
 & b'_2 \arrow[d, Rightarrow] \arrow[drr, Rightarrow] & & & \\
 & b'_1 \arrow[dl, Rightarrow] \arrow[dr, Rightarrow] \arrow[dd, Rightarrow] & & a' \arrow[ddd, Rightarrow] & \\
b''_{\mathrm{coh}} \arrow[dr, Rightarrow] & & b''_{\mathrm{hom}} \arrow[dl, Rightarrow] & & \\
 & b \arrow[drr, Rightarrow] & & & \\
 & & & a & .
\end{tikzcd}
\]
Dans tous les cas à ma connaissance, les dérivateurs satisfont la condition la plus forte \textbf{Der} 3 $b'_2$.
\msp{page 113}

\dgrblock{dgrb69}{Two families of axioms with different antecedents}
\begin{dcomment}{exposition}{p112:Dans tous les cas à ma connaissance, les dérivateurs satisfont la condition la plus forte}{Two families of axioms with different antecedents}
The $a$- and $b$-conditions should be read as assertions about the strength of the coefficient theory, not as successive definitions of the same phrase. Condition $a$ says that a weakly contractible base does not create new morphisms between constant coefficients. Condition $a'$ adds that it does not create new locally constant objects either. The $b$-conditions ask for analogous invariance under arbitrary weak homotopy equivalences, rather than just maps to a point.

Here is a direct check of the diagram. If $f^*$ is an equivalence on local systems, it is fully faithful there. Full faithfulness gives cohomological and homological invariance by testing morphisms from and to all constant objects. Each of these implies invariance for constant coefficients. Applying the latter to $Z\to e$ proves $a$. Applying the equivalence on local systems to this same map proves $a'$, because every object at $e$ is constant. Finally $a'$ contains $a$ by definition. These are the arrows drawn in the source; the diagram does not claim that full faithfulness alone gives essential surjectivity.

There is a small but consequential inconsistency in the parenthesis defining $b$ on page~111. Its first wording concerns constant coefficients, but the parenthesis writes $\D(Z)_{\mathrm{loc.const.}}$. Taken literally this is $b'_1$, the stronger condition introduced on the next line. For the displayed implication diagram and the earlier definition of a $\D$-equivalence, the subscript in that parenthesis must be read as ``constant''. We retain the edition's wording and record the distinction in the audit.

Modern homotopy invariance supplies $b'_2$ for ordinary Grothendieck derivators with all small categories as shapes. A restricted diagram domain still requires care: the statement must only quantify over shapes in that domain, and a proof that passes through simplicial resolutions or larger indexing categories cannot be imported without an extension hypothesis. This is the reason the manuscript continues to list these properties as separate assumptions. Its present calculation needs the explicitly stated hypothesis, not an unannounced change of domain.
\end{dcomment}


Je rappelle aussi la :

\textbf{Définition 1 :} Soit $f : X \to Y$ dans \Dia. On dit que $f$ est une $\DD$-\textit{équivalence cohomologique} (sous-entendu : “\textit{complète}”) si pour tout $\DD$-coefficient $F$ \underline{localement} \underline{constant} sur $Y$,
\[ H^{*}_{\DD}(Y,F) \xrightarrow{\ \sim\ } H^{*}_{\DD}(X,f^{*}F) \quad ; \]
on dit que $f$ est une $\DD$-\textit{équivalence homologique} (“\textit{complète}”) si pour tout $F$ comme \texttt{[ci-]}dessus,
\[ H_{*}^{\DD}(X,f^{*}F) \xrightarrow{\ \sim\ } H_{*}^{\DD}(Y,F) \quad . \]
On dit que $f$ est une $\DD$-\textit{équivalence} 1-\textit{complète} (resp. 2-\textit{complète}) si le foncteur
\[ f^{*} : \DD(Y)_{\text{loc. const.}} \longrightarrow \DD(X)_{\text{loc. const.}} \]
est pleinement fidèle (resp. une équivalence).
\begingroup\small
\[
\begin{tikzcd}[column sep=small]
 & \text{$\DD$-équivalence 2-complète} \arrow[d, Rightarrow] & & \\
 & \text{$\DD$-équivalence 1-complète} \arrow[dl, Rightarrow] \arrow[dr, Rightarrow] & & \\
\text{$\DD$-équivalence cohomologique} \arrow[dr, Rightarrow] & & \text{$\DD$-équivalence homologique} \arrow[dl, Rightarrow] & \\
 & \text{$\DD$-équivalence} & & .
\end{tikzcd}
\]
\endgroup
Si dans les sous-catégories $\DD(X)_{\text{loc. const.}}$ il y a une loi $*$ (associative, unitaire, commutative) et des \underline{\textsf{Hom}} internes associés, préservés par les images inverses, alors les $\DD$-équivalences cohomologiques sont déjà 1-complètes et sont donc des $\DD$-équivalences homologiques, donc dans ce cas particulier
\printedpage{68}
% head: A. Grothendieck : Dérivateurs, Ch. I, Édition des Universités de Montpellier II et Paris VII
% foot: Version 8 juillet 2001
\msp{page 114 vide}
\msp{page 115}
\[
\begin{tikzcd}
\text{$\DD$-équivalence 2-complète} \arrow[d, Rightarrow] & & \\
\text{$\DD$-équivalence 1-complète} \arrow[r, Leftrightarrow] & \text{$\DD$-équivalence cohomologique} \arrow[d, Rightarrow] & \\
 & \text{$\DD$-équivalence homologique} \arrow[d, Rightarrow] & \\
 & \text{$\DD$-équivalence} & .
\end{tikzcd}
\]
Dans $\DD = \mathsf{HOT}$, toutes ces notions devraient être équivalentes, et équivalentes à la condition encore plus faible en apparence, que $f : X \to Y$ soit une \Hot-équivalence. Mais je ne connais que \HOT{} comme dérivateur où la $\DD$-équivalence d'une flèche $f : X \to Y$ dans \Dia{} implique déjà la $\DD$-équivalence cohomologique (complète), ou la $\DD$-équivalence homologique (complète), voire même la $\DD$-équivalence 1-complète. (Je reviendrai là-dessus tantôt.)

\textbf{Proposition 1.} \textit{Soit $f : X \to Y$ dans} \Dia. \textit{Conditions équivalentes :}

a) \textit{Pour tout $y \in \mathsf{Ob}\, Y$, l'inclusion $X_y \to X/y$ est une $\DD$-équivalence, et pour toute flèche $y' \to y$ dans $Y$, $X/y' \to X/y$ est une $\DD$-équivalence.}

b) \textit{Pour tout coefficient \underline{constant} $F$ sur $Y$, $f_{*}f^{*}F$ est localement constant, et se calcule “fibre par fibre”.}

(C'est la prop. 1, p. 76).

\textbf{Proposition 1$'$.} Dual.

\msp{page 116}

\textbf{Définition 2.} Sous les conditions équivalentes de la proposition 1, on dit que $f$ est une $\DD$-\textit{fibration cohomologique faible.} Sous les conditions duales, que $f$ est une $\DD$-\textit{fibration homologique faible.}

C'est ce que dans la définition 1, p. 77 j'appelais “$\DD$-fibration cohomologique” ou “$\DD$-fibration à gauche”, et dualement. Maintenant la notion me paraît trop faible pour mériter ce nom, que je réserve pour le cas “complet”.

\textbf{Proposition 2.} \textit{Les $\DD$-fibrations cohomologiques faibles sont stables par changements de base lisses, les $\DD$-fibrations homologiques faibles stables par changements de base propres.}

Je devrais construire un exemple où $f$ est une $\DD$-fibration cohomologique faible, et ne l'est plus après changement de base.

\dgrblock{dgrb70}{A failure already among finite posets}
\begin{dcomment}{development}{p116:Je devrais construire un exemple où}{A failure already among finite posets}
The example at manuscript page~77 uses isomorphisms in the base to expose the failure of arbitrary base change. The same failure occurs for finite posets; it is therefore relevant even to the restricted domains repeatedly considered in these pages.

\textbf{A four-object example.}\ours\quad
Let $Y=[2]=\{0<1<2\}$ and let $X$ be the poset whose distinct nonidentity relations are
\[
 x<a,\qquad b<a,\qquad b<z.
\]
Define $f:X\to Y$ by $f(x)=0$, $f(a)=f(b)=1$, and $f(z)=2$. Then $f$ is both a weak cohomological and a weak homological $\HOT$-fibration. Its pullback along the long edge $[1]\to[2]$, selecting $0$ and $2$, is neither.

To prove the first assertion, the strict fibres are $\{x\}$, $\{b<a\}$, and $\{z\}$, all contractible. Since $Y$ is a poset, the overcategories $X/y$ are the subposets on the elements mapping to $\{0,\ldots,y\}$. They are respectively
\[
 \{x\},\qquad \{x\to a\leftarrow b\},\qquad X.
\]
Their nerves realize to a point, an interval subdivided into two edges, and an interval subdivided into three edges. They are all contractible. Every fibre-to-overcategory inclusion and every induced transition between these categories is consequently a weak homotopy equivalence. This is exactly the criterion for a weak cohomological fibration. Dually, the undercategories $y\backslash X$ are $X$, the poset $\{a\leftarrow b\to z\}$, and $\{z\}$; the same argument proves the weak homological condition.

In the pullback, only $x$ and $z$ remain. There is no morphism $x\to z$ in $X$: the connecting path in its nerve traverses $b\to a$ backwards. Thus the pullback category $X'$ is the discrete two-point category, mapping to the two endpoints of $[1]$. The overcategory at the terminal endpoint is all of $X'$, whereas that strict fibre is a point. Its inclusion is not an equivalence on components. The dual failure occurs at the initial endpoint. Even ordinary degree-zero cohomology detects the defect: for a nonzero coefficient field $k$, the relevant comparison is $k\times k\to k$, not an isomorphism.

This settles the request for an example on page~116 without any infinite category or nontrivial automorphism. The original weak conditions hold for every coefficient derivator in which weakly contractible shapes are completely aspherical; the failure is visible for both spaces and ordinary derived coefficients. The construction is supplied here with no priority claim. We checked the local-system and smoothness treatments cited in this commentary and did not locate this precise four-object example. In particular, the ``weakly smooth'' functors of \citet[Definition~4.1 and Proposition~4.3]{Maltsiniotis2005Asphericite} have a different, comma-asphericity condition; their terminology must not be substituted for the weak cohomological condition used here.
\end{dcomment}


% ---- from Der-ChI-234_p071-075.tex
\printedpage{69}
% head: A. Grothendieck : Dérivateurs, Ch. I, Édition des Universités de Montpellier II et Paris VII

\textbf{Proposition 3.} \textit{Soit}
\[
\begin{tikzcd}
X \arrow[rr, "f"] \arrow[dr, "p"'] & & X' \arrow[dl, "p'"] \\
& S &
\end{tikzcd}
\]
\textit{dans \Dia, avec $p$ et $p'$ des $\D$-fibrations cohomologiques faibles} (resp. \textit{des $\D$-fibrations homologiques faibles}), \textit{soit $S_0 \subset S$ tel que $\pi_0(S_0) \to \pi_0(S)$ soit surjective. Pour que $f$ soit une $\D$-équivalence, il suffit que les $f_s : X_s \to X'_s$, pour $s \in S_0$, le soient.}

\textbf{Proposition 4.} \textit{Si $f : X \to Y$ dans \Dia{} est $\D$-propre} (resp. $\D$-\textit{lisse}), \textit{alors $f$ est une $\D$-fibration cohomologique} (resp. \textit{homologique}) \textit{faible si et seulement si pour tout $F$ constant sur $Y$, $f_*f^*F$} (resp. $f_!f^*F$) \textit{est localement constant sur $Y$. Et s'il en est ainsi, pour tout $Y' \to Y$ dans \Dia, $f' : X' = X \times_Y Y' \to Y'$ est encore une $\D$-fibration cohomologique} (resp. \textit{homologique}) \textit{faible.}

Immédiat sur le critère \textit{b}) de la proposition 1.

\msp{page 117}\textbf{Proposition 5.} \textit{Soit $f : X \to Y$ dans \Dia. Considérons les conditions suivantes :}

a) \textit{Pour tout $Y'' \to Y'$ sur $Y$ dans \Dia, si $Y''$ et $Y'$ sont \Hot-asphériques, alors}
\[
X'' = X \times_Y Y'' \longrightarrow X' = X \times_Y Y'
\]
\textit{est une $\D$-équivalence.}
\[
\begin{tikzcd}
X & X' \arrow[l] & X'' \arrow[l] \\
Y \arrow[from=u] & Y' \arrow[l] \arrow[from=u] & Y'' \arrow[l, "p"] \arrow[from=u]
\end{tikzcd}
\]

a$'$) \textit{Comme} a), \textit{mais en supposant $p : Y'' \to Y'$ une \Hot-équivalence (au lieu de supposer $Y'$, $Y''$ \Hot-asphériques).}

a$_1$) \textit{Comme} a), \textit{mais en supposant (au lieu de l'asphéricité) que $Y'$ a un objet final, et $Y''$ est la catégorie unité.}

a$_2$) \textit{Comme} a), \textit{mais en supposant que $Y'$ a un objet cofinal, et $Y''$ est la catégorie unité.}

b$_1$) \textit{$f$ est une $\D$-fibration cohomologique faible, et le reste après tout changement de base.}

b$_2$) \textit{$f$ est une $\D$-fibration homologique faible, et le reste après tout changement de base.}

\printedpage{70}
% head: A. Grothendieck : Dérivateurs, Ch. I, Édition des Universités de Montpellier II et Paris VII
\textit{On a alors le diagramme d'implications tautologiques}
\[
\begin{tikzcd}[column sep=small]
& & & a' \arrow[d, Rightarrow] & & & & \\
& & & a \arrow[dlll, Rightarrow] \arrow[drrr, Rightarrow] & & & & \\
a_1 & \Longleftrightarrow & b_1 & & b_2 & \Longleftrightarrow & a_2 & .
\end{tikzcd}
\]
\textit{Lorsque $\D$ satisfait à l'axiome} \textbf{Der} 3 a$'$ (p. 111), \textit{alors toutes ces conditions, sauf} a$'$), \textit{sont équivalentes. Lorsque $\D$ satisfait à l'axiome} \textbf{Der} 3 b$'_1$ \textit{(ou seulement} \textbf{Der} 3 b$''_{\mathrm{coh}}$ ou b$''_{\mathrm{hom}}$\textit{),} \textit{alors toutes les conditions (y inclus} a$'$\textit{) sont équivalentes.}

Pour simplifier, on va par la suite supposer que l'axiome \textbf{Der} 3 $a'$) est satisfait.

\textbf{Définition 3.} Sous les conditions équivalentes $a$, $a_1$, $b_1$, $a_2$, $b_2$ de la proposition 4, on dit que $f$ est une $\D$-\textit{fibration}. % [sic] proposition 4 -> proposition 5, the conditions a, a_1, b_1, a_2, b_2 are those of proposition 5

\dgrblock{dgrb71}{Why universality identifies the two weak conditions}
\begin{dcomment}{exposition}{p117:Sous les conditions équivalentes $a$, $a_1$, $b_1$, $a_2$, $b_2$ de la proposition 4}{Why universality identifies the two weak conditions}
Proposition~5 organizes the earlier failed argument by changing its input. Under $a'$, one may pull back between weakly contractible bases. The preceding example shows why this is an additional requirement; it is not implied by the fibre criterion on the original base.

Suppose first that the weak cohomological condition holds after every base change. Over a category $B$ with a terminal object, the direct image of a constant coefficient is locally constant and hence constant by $a'$. Its cohomology over $B$ is therefore its value at any selected object $b$. The pointwise formula identifies that value with the cohomology of the strict fibre. Hence every point inclusion into such a base produces a $\D$-equivalence on total categories. Conversely, these point tests applied over $Y/y$, and then at its objects $u:y\to y'$, give the fibre and transport comparisons, exactly as in the proof of Proposition~4 above, now with constant coefficients only. This proves $a_1\Leftrightarrow b_1$.

More generally, take a pullback $p:B'\to B$ with both bases weakly contractible. Put $G=f_*M_X$ and $G'=f'_*M_{X'}$. They are locally constant. Their pointwise comparison is an isomorphism because both universal weak fibre formulas use the same strict fibre; thus $G'\simeq p^*G$. Since $G$ is constant by $a'$, the comparison of its cohomology over the two contractible bases is invertible. This proves $b_1\Rightarrow a$. The dual proof gives $b_2\Rightarrow a$, and the remaining implications are immediate. In $a_2$, the printed ``objet cofinal'' must carry the dual, initial-object meaning required by this argument.

If the stronger $b''_{\mathrm{coh}}$ is assumed, the same calculation works with an arbitrary $\Hot$-equivalence $p$: $G$ need not become constant, because $p$ preserves cohomology of every locally constant coefficient. That proves $a'$. The homological version is dual. Thus the two different uses of the prime on $a'$ must be kept apart: \textbf{Der}~3~$a'$ is an axiom about coefficients on contractible categories, whereas condition~$a'$ of Proposition~5 is a universal pullback test using all weak homotopy equivalences.
\end{dcomment}


\msp{page 118}\textbf{Proposition 6.} \textit{Supposons que $f : X \to Y$ dans \Dia{} soit $\D$-propre} (resp. $\D$-\textit{lisse}). \textit{Conditions équivalentes :}

a) \textit{$f$ est une $\D$-fibration.}

b) \textit{$f$ est une $\D$-fibration cohomologique} (resp. \textit{homologique}) \textit{faible.}

c) \textit{Pour tout $y' \to y$ dans $Y$, $X/y' \to X/y$} (resp. $y\backslash X \to y'\backslash X$) \textit{est une $\D$-équivalence.}

d) \textit{Pour tout $F$ constant sur $Y$, $f_*f^*F$} (resp. $f_!f^*F$) \textit{est localement constant sur $Y$.}

d$'$) \textit{Pour tout $F$ localement constant sur $Y$, $f_*f^*F$} (resp. $f_!f^*F$) \textit{est localement constant sur $Y$.}

Cela résulte formellement de Prop. 4 et Prop. 5.

\textbf{Corollaire.} \textit{Supposons $f : X \to Y$} \Cat\textit{-fibrant} (resp. \Cat\textit{-cofibrant}). \textit{Pour que $f$ soit une $\D$-fibration, il faut et il suffit que pour tout $u : y \to y'$ dans $Y$, $u^* : X_{y'} \to X_y$} (resp. $u_* : X_y \to X_{y'}$) \textit{soit une $\D$-équivalence.}

Il est immédiat que la notion de $\D$-fibration est stable par changement de base, mais il est douteux qu'elle le soit par composition.\footnote[60]{N.B. Si les $\D$-équivalences sont déjà des \Hot-équivalences, alors la notion coïncide avec celle de \Hot-fibration, et elle est stable par composition. Également, si les $\D$-équivalences sont au sens des $\D$-équivalences cohomologiques, ou homologiques, alors la notion coïncide avec celle des $\D$-fibrations cohomologiques (resp. homologiques) ci-dessous, et est également stable par composition.} (Ici encore, j'ai omis de faire un contre-exemple, avec des dérivateurs $\D = \sfup{Der}(k\text{-}\sfup{Mod})$ disons.) Autre lacune de la notion : il n'est pas clair, si $f$ est une $\D$-fibration et $F$ un $\D$-coefficient sur $Y$, si
\[
f_*f^*F \qquad \text{et} \qquad f_!f^*F
\]
\printedpage{71}
% head: A. Grothendieck : Dérivateurs, Ch. I, Édition des Universités de Montpellier II et Paris VII
se calculent “fibre par fibre”, comme on s'attendrait dans le cas d'un “fibré”.

Par contre, voici un résultat sympa :

\dgrblock{dgrb72}{Two fibrations whose composite is not a fibration}
\begin{dcomment}{development}{p118:Par contre, voici un résultat sympa :}{Two fibrations whose composite is not a fibration}
The doubt about composition can be settled by a small categorical example. It uses the same loss of monodromy information as the preceding counterexamples, but places that loss in a transition functor. We work with $\D=\mathsf{Der}(\mathbb Q\text{-}\mathsf{Mod})$ on all small categories and with the final definition of $\D$-fibration on page 117, which includes universality under base change.

\textbf{Proposition (failure of composition).}\ours
There are $\D$-fibrations $p:X\to Y$ and $q:Y\to[1]$ such that $qp$ is not even a weak cohomological $\D$-fibration. The first map can be chosen to be a covering functor, hence both proper and smooth; the second is a Grothendieck opfibration.

Let $G=C_2=\{1,\sigma\}$. Define $Y$ to have two objects $a,b$, with
\[
 Y(a,a)=\{1_a\},\quad Y(b,b)=G,\quad
 Y(a,b)=G,\quad Y(b,a)=\varnothing.
\]
Composition on $Y(b,b)$ is multiplication, and postcomposition of the arrow $g:a\to b$ by $h\in G$ is the arrow $hg$. Send $a$ to $0$ and $b$ to $1$. This defines $q:Y\to[1]$. It is the Grothendieck construction of the functor $e\to BG$. More explicitly, any arrow $g:a\to b$ is opcartesian: for every arrow $k:a\to b$, the unique factorization through $g$ is multiplication by $kg^{-1}$. Its endpoint fibres are $e$ and $BG$, and its transition functor is $e\to BG$.

This transition is a $\D$-equivalence. For a complex $M$ of rational vector spaces with trivial $G$-action, averaging by $(1+\sigma)/2$ makes invariants exact, and coinvariants are exact for the same reason. Thus
\[
 \mathbb R\Gamma(G,M)=M,
 \qquad \mathbb Q\mathbin{\otimes^{\mathbb L}_{\mathbb Q[G]}}M=M.
\]
These identities are natural for arbitrary, including unbounded, complexes: the trivial module is projective over $\mathbb Q[G]$, and the averaging contraction is degreewise compatible with the differential. The opfibration criterion on page 118 now makes $q$ a $\D$-fibration, with all base changes covered by that criterion.

Let $Z$ be the discrete category on the set $G$. Define a functor $L:Y\to\Cat$ by $L(a)=L(b)=Z$, and let every arrow labeled $g$ act by left translation $z\mapsto gz$. Put $X=\int_Y L$ and let $p$ be the projection. All its transport functors are isomorphisms. Every arrow of $Y$ therefore has a unique lift with either prescribed source or prescribed target. In particular $p$ is both a discrete fibration and a discrete opfibration, and the transport criterion makes it a $\D$-fibration (indeed a $\HOT$-fibration).

The composite $qp$ is an opfibration whose fibre at $0$ is the discrete two-point category $Z$. Its fibre at $1$ is the action groupoid $G\ltimes G$ for the left regular action. There is exactly one arrow between any ordered pair of objects of this groupoid, so it is equivalent to the terminal category. The transition functor is the inclusion
\[
 Z_{\mathrm{disc}}\longrightarrow G\ltimes G.
\]
On cohomology in degree zero it induces
\[
 \mathbb Q\longrightarrow\mathbb Q^2,
 \qquad r\longmapsto(r,r),
\]
which is not an isomorphism. Since $qp$ is proper, the criterion of Proposition 6 identifies this failure with failure to be a weak cohomological fibration. Equivalently, $(qp)_*\mathbb Q_X$ has values $\mathbb Q^2$ at $0$ and $\mathbb Q$ at $1$, and its transition is the displayed diagonal. This completes the construction and the proof.\hfill$\square$

The example explains why the positive answer concerning $f_*f^*F$ does not force closure under composition. The coefficient $p_*\mathbb Q_X$ on $Y$ is the permutation local system attached to the regular $G$-set. It is not pulled back from a coefficient on $[1]$: every such pullback has trivial $G$-action at $b$. The constant-coefficient test for $q$ cannot control this local system. The stronger cohomological fibrations introduced on page 121 impose exactly the missing condition. Our literature check covered the smooth-functor formalism of \citet{Maltsiniotis2005Asphericite} and the locally constant coefficient theory of \citet{Cisinski2009LocallyConstant}; we did not find this particular composition counterexample there. It is an elementary development of the mechanism already present in Grothendieck's finite-group example.
\end{dcomment}


\msp{page 119}\textbf{Proposition 7.} \textit{Soit $f : X \to Y$ dans \Dia. Conditions équivalentes :}

a) \textit{$f$ est une $\D$-équivalence “universelle”,} i.e. \textit{pour tout $Y' \to Y$ dans \Dia, $f' : X' = X \times_Y Y' \to Y'$ est une $\D$-équivalence.}

b) \textit{$f$ est une $\D$-fibration à fibres $\D$-asphériques.}

\textbf{Définition 4.} Sous les conditions équivalentes de la proposition 7, on dit que $f$ est une $\D$-fibration \textit{triviale}.

Mais attention, $f$ peut-être une $\D$-fibration et une $\D$-équivalence, sans que la $\D$-fibration soit triviale \textit{i.e.} sans que les fibres soient $\D$-asphériques, cf. p. 87 et suivantes. % [sic] peut-être -> peut être, verb "pouvoir" followed by infinitive
Et également (devrait remonter) même si $\D$ satisfait \textbf{Der} 3 $b_2$ (la plus forte des variantes de \textbf{Der} 3), donc quand \underline{toutes} les conditions de la proposition 5, y inclus $a'$), sont équivalentes, celles-ci (\textit{i.e.} que $f$ soit une $\D$-fibration) n'impliquent pourtant pas, en général, que dans la condition $a'$) on puisse se borner à $p : Y'' \to Y'$ une $\D$-équivalence (au lieu d'une \Hot-équivalence). En d'autres termes, si on a un diagramme cartésien dans \Dia
\[
\begin{tikzcd}
X \arrow[d, "f"'] & X' \arrow[l, "q"'] \arrow[d, "f'"] \\
Y & Y' \arrow[l, "p"] & ,
\end{tikzcd}
\]
avec $f$ une $\D$-fibration, et $p$ une $\D$-équivalence, il ne s'ensuit pas forcément que $q$ soit une $\D$-équivalence. (Mais ce sera pourtant le cas, “bien sûr”, si $\D = \HOT$.)

\msp{page 120}Par contre, on a :

\textbf{Proposition 8.} \textit{Soit $f : X \to Y$ dans \Dia. Conditions équivalentes :}

a) \textit{$f$ est une $\D$-fibration.}

b) \textit{Pour tout $p : Y'' \to Y'$ sur $Y$, si $p$ est une $\D$-équivalence} 1-\textit{complète} (déf. 1, p. 113), \textit{alors $q : X'' \to X'$ est une $\D$-équivalence.}
\[
\begin{tikzcd}
X \arrow[d, "f"'] & X' \arrow[l] \arrow[d, "f'"'] & X'' \arrow[l, "q"'] \arrow[d, "f''"] \\
Y & Y' \arrow[l] & Y'' \arrow[l, "p"]
\end{tikzcd}
\]
\textbf{Corollaire.} \textit{Supposons que dans \Dia, toute $\D$-équivalence soit} 1-\textit{complète.} (N.B. Je ne connais que le dérivateur \HOT{} pour satisfaire à cette condition, qui paraît quasiment caractéristique de \HOT.) \textit{Alors $f : X \to Y$ est une $\D$-fibration si et seulement si pour tout $p : Y'' \to Y'$ sur $Y$, si $p$ est une $\D$-équivalence, alors $q : X'' \to X'$ aussi. Dans ce cas exceptionnel parmi les dérivateurs, la notion de $\D$-fibration est stable par composition} (ce qui est tautologique sur le critère que je viens de donner.)

\dgrblock{dgrb73}{Universal equivalence and the 1-complete test}
\begin{dcomment}{exposition}{p120:ce qui est tautologique sur le critère que je viens de donner.}{Universal equivalence and the 1-complete test}
Two different universal assertions occur here. A universal $\D$-equivalence is a functor whose every base change is a $\D$-equivalence. A $\D$-fibration preserves certain equivalences in diagrams over its base; it need not itself be a $\D$-equivalence. Proposition~7 identifies their intersection with the condition that all strict fibres be $\D$-aspherical. For the forward implication, base change to each point gives aspherical fibres, while the two-out-of-three property applied to the two universally equivalent projections gives the fibration criterion. For the reverse implication, the universal weak fibre formula makes $f_*M_X$ a constant object $M_Y$: its unit is an isomorphism at every $y$ because $X_y$ is $\D$-aspherical. Taking global cohomology proves the equivalence, and the same argument applies after base change.

This result does not permit replacement of ``universal'' by ``ordinary''. The rational monodromy examples already show that a fibration which is a $\D$-equivalence can have a fibre that is not $\D$-aspherical. The issue is not a failure of the fibre formula, but a loss of information upon taking global cohomology of a nonconstant local system.

Proposition~8 repairs the pullback assertion by strengthening the equivalence in the base. If $p:Y''\to Y'$ is 1-complete, it preserves cohomology of every local system. For a constant $M$, universality of the weak fibration condition makes $f'_*M_{X'}$ locally constant and identifies its pullback with $f''_*M_{X''}$. Taking cohomology proves that $X''\to X'$ is a $\D$-equivalence. Conversely, apply the preservation assertion to $e\to B$ with $B$ weakly contractible. Under $a'$, restriction on local systems is an equivalence and hence 1-complete, so the point tests of Proposition~5 follow.

There is a useful distinction between the result here and the stronger definition below. Proposition~8 asks that pulling back a 1-complete equivalence produce a constant-coefficient equivalence. It does not assert that the pullback is itself 1-complete. Demanding that stronger conclusion supplies a composition-stable class by definition, but it is an extra condition unless the internal-Hom argument identifies the relevant equivalence classes.
\end{dcomment}


\printedpage{72}
% head: A. Grothendieck : Dérivateurs, Ch. I, Édition des Universités de Montpellier II et Paris VII

On va renforcer la notion de $\D$-fibration, pour la rendre plus souple, et notamment stable par composition.

\textbf{Proposition 9.} \textit{Soit $f : X \to Y$ dans \Dia. Conditions équivalentes :}

a) \textit{Pour tout $y \in Y$, $X_y \to X/y$ est une $\D$-\underline{équivalence} \underline{cohomologique}} (déf. 1, p. 113), \textit{et de même pour $X/y \to X/y'$ pour toute flèche $y \to y'$ dans $Y$. Et ces conditions sur $f$ restent valables après tout changement de base} (N.B. il suffit que la première de ces conditions le reste).

\msp{page 121}b) \textit{pour tout $F \in \Ob\ \D(X)$ \underline{localement} \underline{constant}, $f_*F$ est localement constant sur $Y$ et se calcule “fibre par fibre” -- et itou après tout changement de base.} % Doubt resolved: attached PDF leaf 36 (printed p.72) has an en dash, encoded --.

c) \textit{Pour tout $p : Y'' \to Y'$ sur $Y$, si $p$ est une $\D$-équivalence cohomologique, alors il en est de même de $q : X'' \to X'$.}

c$'$) \textit{Comme} c, \textit{mais en supposant comme hypothèse sur $p$ que $Y''$ soit la catégorie finale, et $Y'$ ait un objet final.}

\textsc{Démonstration}. Dans l'ordre ($\pm$ tautologique)
\[
\mathrm{a} \Longrightarrow \mathrm{b} \Longrightarrow \mathrm{c} \Longrightarrow \mathrm{c}' \Longrightarrow \mathrm{a} \quad .
\]

\textbf{Définition 5.} Sous les conditions équivalentes de la proposition 9, on dit que $f$ est une $\D$-\textit{fibration cohomologique}. Sous les conditions duales, on dit que $f$ est une $\D$-\textit{fibration homologique}. On dit que $f$ est une $\D$-\textit{fibration forte}, si c'est à la fois une $\D$-fibration cohomologique et une $\D$-fibration homologique.

Ces notions renforcent la notion de $\D$-fibration. Elles ont l'avantage de plus d'être \underline{stables} \underline{par} \underline{composition}, non seulement par changement de base.

\dgrblock{dgrb74}{The reason the strong notion composes}
\begin{dcomment}{exposition}{p121:Elles ont l'avantage de plus d'être}{The reason the strong notion composes}
We may now prove the composition assertion without an appeal to intuition about bundles. Let $X\xrightarrow{f}Y\xrightarrow{g}Z$ be cohomological fibrations in the revised sense. For $L\in\mathcal L(X)$, the object $f_*L$ is locally constant on $Y$, and then $g_*f_*L$ is locally constant on $Z$. The canonical composition isomorphism identifies it with $(gf)_*L$. This is the coefficient condition that was missing for the earlier notion.

For the universal base-change assertion, pull back the composite along any $Z'\to Z$, and factor the resulting rectangle into the two cartesian squares for $g$ and $f$. The comparison for $g$, applied to the locally constant object $f_*L$, is invertible. The comparison for $f$, applied to $L$, is invertible. Pasting gives the comparison for $gf$. The same reasoning works for every locally constant coefficient on the new total category, because the hypotheses on $f$ and $g$ were universal. Specializing to a point in $Z$ gives the fibre formula. Alternatively, condition~$c$ of Proposition~9 proves composition by applying two successive pullbacks to a cohomological equivalence. The two arguments express the same closure property at the levels of coefficients and indexing categories.

The dual reasoning applies to homological fibrations, and hence to strong fibrations. It does not prove that every $\D$-fibration composes: a weak fibration controls direct images of constant coefficients, whereas the first direct image in a composite is usually only locally constant. The explicit counterexample at page~118 realizes exactly this obstruction. The questions about $f_*f^*F$ and $f_!f^*F$, for an arbitrary coefficient on the base, are different again; the factorization-category argument at page~82 proves the required fibre formulas under the stated domain hypotheses. There is no conflict, because a locally constant object produced on the intermediate category need not be pulled back from the final base.

For the proposed 1-complete and 2-complete fibrations of page~122, the phrase ``not stable under composition (as far as I know)'' supplies no proof. In the compatible closed monoidal situation, the 1-complete variant reduces to the cohomological one and therefore composes. The 2-complete variant nevertheless fails to compose, even for the represented set-valued derivator: the development ``The 2-complete variant also fails to compose'' constructs two functors between finite categories from a normalized $C_2$-valued cocycle on a triangulated two-sphere. Their composite fails a permitted contractible-base test because restriction gives only trivial $C_2$-actions. This restriction is fully faithful but not essentially surjective. An internal Hom controls morphisms between objects already descended; it does not supply descent of every local system. The counterexample thus settles the 2-complete question without extending its negative conclusion to the 1-complete class.
\end{dcomment}


\msp{page 122}La notion de $\D$-fibration forte a été choisie pour coiffer les deux notions duales de $\D$-fibration cohomologique et homologique (à défaut d'une meilleur notion) tout en restant stable par composition. On pourrait exiger que pour tout % [sic] meilleur -> meilleure, agreement with "notion"
\[
Y'' \longrightarrow Y'
\]
sur $Y$, avec $Y'$, $Y''$ \Hot-asphériques, $q : X'' \to X'$ soit une $\D$-équivalence 1-complète (voire même 2-complète), appeler ça les $\D$-\textit{fibrations} 1-\textit{complètes} (resp. 2-\textit{complètes}). Ça renforce de façon assez élégante la notion de $\D$-fibration forte, mais a le grand tort de ne pas être stable par composition (que je sache!), sauf dans le cas particulier où on a des \underline{\Hom} internes sympas (cf. bas de la page 113), auquel cas la notion de $\D$-équivalence implique déjà celle de $\D$-équivalence 1-complète,
\printedpage{73}
% head: A. Grothendieck : Dérivateurs, Ch. I, Édition des Universités de Montpellier II et Paris VII
\msp{page 123}donc aussi celle de $\D$-équivalence homologique. Donc dans ce cas
\[
\begin{tikzcd}
\text{$\D$-fibration cohomologique} \arrow[r, Leftrightarrow] \arrow[d, Rightarrow] & \text{$\D$-fibration 1-complète} \\
\text{$\D$-fibration homologique} \quad , &
\end{tikzcd}
\]
et il n'y a pas à chercher plus loin que les $\D$-fibrations cohomologiques.

\dgrblock{dgrb75}{The 2-complete variant also fails to compose}
\begin{dcomment}{development}{p122:La notion de $\D$-fibration forte a été choisie}{The 2-complete variant also fails to compose}
The qualification ``que je sache'' is not a proof of failure. Nevertheless, the proposed 2-complete class really need not be closed under composition, even for a represented derivator with all small categories as shapes. We give a categorical counterexample, including the coherence responsible for the failure. The existence of the local-systems derivator in \citet[Theorem~3.13]{Cisinski2009LocallyConstant} does not supply the missing composition property. We checked the local-systems treatment just cited and did not find this particular construction there. The construction below is a verification given here, without a claim of priority.

\textbf{Proposition.}\ours\quad
Let $\D_{\mathrm{set}}(A)=\operatorname{Fun}(A^{\mathrm{op}},\mathsf{Set})$. There are functors
\[
 X\xrightarrow{p}Y\xrightarrow{q}[1]
\]
which are both 2-complete fibrations in the precise sense proposed on manuscript page~122, while $qp$ is not one. The categories $X$ and $Y$ are finite, and the functor $p$ is even a $\HOT$-fibration. Thus the failure persists in the usual domain of all small categories, with cartesian closed coefficients.

Write $\Pi_1(A)$ for the fundamental groupoid of $|NA|$, or equivalently the groupoid obtained by inverting all arrows of $A$. A locally constant set-valued presheaf factors through this groupoid. Consequently
\[
 \mathcal L_{\mathrm{set}}(A)
 \simeq\operatorname{Fun}(\Pi_1(A)^{\mathrm{op}},\mathsf{Set}).
\]
In particular, an equivalence of fundamental groupoids induces a 2-complete equivalence. We will verify this condition for the two factors, and exhibit an object which cannot descend for the composite.

Let $B$ be the poset of nonempty proper subsets of $\{0,1,2,3\}$, ordered by inclusion. Its nerve is the barycentric subdivision of the boundary of a tetrahedron, so $|NB|$ is a two-sphere. Fix the maximal chain
\[
 \sigma=(\{0\}<\{0,1\}<\{0,1,2\}).
\]
Write $C_2=\mathbb Z/2\mathbb Z$ additively. For composable arrows $u:b\to b'$ and $v:b'\to b''$, set $c(u,v)=1$ when they constitute this chain, and set $c(u,v)=0$ otherwise. In particular it is zero if either arrow is an identity. This is a normalized $C_2$-valued 2-cocycle: there is no chain of three nonidentity arrows in $B$, and the cocycle identity with an identity arrow cancels term by term.

Define $E_c$ to have the same objects as $B$, with one arrow $(u,n)$ for each arrow $u$ of $B$ and label $n\in C_2$. Put
\[
 (v,n)\circ(u,m)=(vu,m+n+c(u,v)),
 \qquad 1_b=(1_b,0).
\]
The cocycle identity is precisely associativity. The projection $E_c\to B$ is both a Grothendieck fibration and an opfibration. Indeed, after fixing a lift of one base arrow, any required factorization is obtained by uniquely solving for the $C_2$-label of the other lift. Its fibre at each object is the one-object groupoid $BC_2$, and its transport functors act identically on $C_2$. The coherence is nontrivial, although the individual transport functors are identities.

We need the fundamental group of $E_c$, not the full homotopy type of its nerve. Put $t_b=(1_b,1)$ and $h_{bb'}=(b\leq b',0)$. Its groupoid completion has these generators and the relations
\[
 t_b^2=1,\qquad t_{b'}h_{bb'}=h_{bb'}t_b,
 \qquad
 h_{b'b''}h_{bb'}=h_{bb''}t_b^{c(b,b',b'')}.
\]
Choose a spanning tree of the one-skeleton of $NB$ and use it to identify all objects. The first relations identify the $t_b$ with a single central generator $t$, with $t^2=1$. Every triangle other than $\sigma$ gives the ordinary, untwisted edge relation. These relations kill all remaining edge generators, because the complex obtained from the triangulated sphere by removing the interior of $\sigma$ is a disk, with the same one-skeleton. The remaining triangle relation then says $1=t$. This proves that $E_c$ is connected and that $\Pi_1(E_c)$ is equivalent to the terminal groupoid.

Now choose $b_0=\{0\}$ and let
\[
 q:Y=\int_{[1]}(e\xrightarrow{b_0}B)\longrightarrow[1]
\]
be the Grothendieck construction. Concretely, $Y$ consists of $B$ and one new object $a$, with
\[
 Y(a,b)=B(b_0,b),\quad Y(b,a)=\varnothing,
 \quad Y(a,a)=\{1_a\}.
\]
The functor $q$ sends $a$ to $0$ and $B$ to $1$. There is a retraction $r:Y\to B$ sending $a$ to $b_0$. If $i:B\hookrightarrow Y$ is the inclusion, then $ri=1$ and the arrow $a\to b_0$ gives a natural transformation $1_Y\to ir$. Finally define
\[
 X=Y\times_B E_c,\qquad p:X\longrightarrow Y.
\]
The retraction and its natural transformation lift to $X\to E_c$, using the arrow of $C_2$-label zero over $a\to b_0$. Naturality follows from the normalization of $c$, since this arrow maps to $1_{b_0}$ under $r$. Hence $|NX|\simeq|NE_c|$, and $X$ is simply connected. The strict fibre of $qp$ at $0$ is $BC_2$.

We now check the universal quantifiers in the definition. First, $p$ is a pullback of the opfibration $E_c\to B$; all its transports are equivalences of groupoids. The opfibration criterion already established on page~118 therefore makes $p$ a $\HOT$-fibration, after every base change. A functor between weakly contractible bases is a nerve weak equivalence. Pullback along $p$ consequently gives a nerve weak equivalence, and hence a 2-complete equivalence for sets.

There is also a direct verification of this last assertion. For any weakly contractible category $A\to Y$, the pulled-back cocycle on $A$ is a coboundary, since $H^2(|NA|;C_2)=0$. Writing it as $\delta d$ and replacing the label $n$ of an arrow $u$ by $n+d(u)$ identifies the pulled-back category with $A\times BC_2$. Choose $d$ normalized, so that this identification preserves the vertical $C_2$-loops. A comparison induced by a functor between two such bases preserves their vertical generator, and hence induces an isomorphism of their fundamental groups, both equal to $C_2$. Both total categories are connected. Thus their fundamental groupoids are equivalent, as required.

For $q$, we prove a stronger assertion. For every $A\to[1]$, with no contractibility assumption, the projection
\[
 q_A:Y_A=Y\times_{[1]}A\longrightarrow A
\]
induces an equivalence of fundamental groupoids. Its strict fibres are $e$ or $B$, and it has a section selecting $a$ over $0$ and $b_0$ over $1$. The transport functors preserve these selected objects. Within each fibre there is a unique morphism in its fundamental groupoid from the selected object to any given object, because the fibre is connected and simply connected. These fibre paths give a natural isomorphism, in $\Pi_1(Y_A)$, from the section followed by the projection to the identity. For naturality, factor an arrow into its opcartesian lift and a vertical arrow; transport of the chosen path and composition with the vertical arrow give the same path in the target fibre, by uniqueness. The section is therefore a quasi-inverse to $\Pi_1(q_A)$.

If $A'\to A$ is a functor over $[1]$ with both bases weakly contractible, the two projections just considered and the equivalence $\Pi_1(A')\to\Pi_1(A)$ show that $Y_{A'}\to Y_A$ is a 2-complete equivalence. This proves exactly the page-122 condition for $q$.

For the composite, use the permitted test $\{0\}\to[1]$. Both bases are contractible, but the upstairs comparison is
\[
 BC_2=X_0\longrightarrow X.
\]
Since $X$ is connected and simply connected, all its locally constant set-valued presheaves are constant. Restriction to $BC_2$ therefore has only trivial $C_2$-actions in its essential image. A two-element set on which $1\in C_2$ exchanges the two elements is a locally constant coefficient on $BC_2$ which does not extend to $X$. Restriction is not essentially surjective, so the comparison is not 2-complete. This proves the proposition.\hfill$\square$

The failure is precisely in descent of objects. In this final test the functor from sets to $C_2$-sets with trivial action is fully faithful; the obstruction does not concern morphisms between objects already descended. Thus the example is consistent with closure of the 1-complete class under the internal-Hom hypotheses, and explains why that argument cannot settle the 2-complete variant. Theorem~3.13 of Cisinski constructs coherent diagrams of local systems for a fixed indexing category; it does not assert that the restricted pullback tests defining the proposed class preserve essential surjectivity under a composite.
\end{dcomment}


On pourrait songer, il est vrai, à définir (dans le cas général) une notion plus forte, en exigeant que pour tout $p : Y'' \to Y'$ sur $Y$ qui soit une $\D$-équivalence 1-complète, $q : X'' \to X'$ soit aussi une $\D$-équivalence 1-complète. Cette fois-ci, la notion est stable par composition, mais elle me paraît néanmoins inutilisable (en dehors du cas qu'on vient de dire, où elle coïncide avec la notion de $\D$-fibration cohomologique, ou encore de $\D$-fibration forte), car n'étant pas en mesure de décrire les $\D$-équivalences 1-complètes, je suis aussi incapable de donner des exemples non triviaux de $\D$-fibrations “strictement 1-complètes”.

Pour vérifier qu'une $f : X \to Y$ dans \Dia{} est une $\D$-fibration cohomologique, il faut
\msp{page 124}savoir vérifier que des flèches $q : Z' \to Z$ dans \Dia{} sont des $\D$-équivalences cohomologiques, donc que $H^*_{\D}(Z,F) \overset{\sim}{\rightarrow} H^*_{\D}(Z',q^*(F))$ est \texttt{[un]} isomorphisme pour tout $F \in \Ob\ \D(Z)$ qui est localement constant sur $Z$. Or je ne sais vérifier un énoncé de cette nature (en dehors du cas où je sais d'avance que $Z$ est complètement $\D$-asphérique, \textit{i.e.} que les coefficients localement constants sur $Z$ sont en fait constant) \underline{que} si $q$ est déjà une \Hot-équivalence. % [sic] constant) -> constants), agreement with "coefficients"
C'est dire que je ne sais exhiber des $\D$-fibrations cohomologiques, que quand ce sont déjà des \HOT-fibrations cohomologiques. Je veux, dans la suite, donner des critères utilisables pour des \HOT-fibrations. Mais on va axiomatiser ces développements, en gardant un $\D$ général, sur lequel je vais faire une hypothèse si draconienne que pratiquement, je ne connais que le dérivateur \HOT{} qui y satisfait. Il s'agit de $c_1$ ci dessous.

\textbf{Der} 3 $c_1$) Toute $\D$-équivalence est une $\D$-équivalence cohomologique.

$c_2$) Toute $\D$-équivalence est une $\D$-équivalence homologique.

\msp{page 125}$c'$) Toute $\D$-équivalence est 1-complète.

$c''$) Toute $\D$-équivalence est 2-complète.

N.B. On a
\[
\begin{tikzcd}
& c'' \arrow[d, Rightarrow] & \\
& c' \arrow[dl, Rightarrow] \arrow[dr, Rightarrow] & \\
c_1 & & c_2
\end{tikzcd}
\]

% ---- from Der-ChI-234_p076-076.tex
\printedpage{74}
% head: A. Grothendieck : Dérivateurs, Ch. I, Édition des Universités de Montpellier II et Paris VII
% foot: Version 8 juillet 2001
tautologiquement. Je ne connais que \HOT{} pour satisfaire l'ensemble des deux plus faibles conditions $c_1$, $c_2$. Et en fait, sous réserve d'écrire la théorie du dérivateur \HOT, il doit être vrai qu'elle satisfait même la forme la plus forte $c''$. (Et les \HOT-équivalences doivent être simplement les \Hot-équivalences.) (Cf. le “lemme” p.~109.)

La proposition suivante est une tautologie.

\noindent\textbf{Proposition 10.} \textit{Supposons que $\D$ satisfasse la condition \textup{\textbf{Der} 3 $\mathrm{c}_1$ (p.~124),} par exemple $\D = \text{\HOT}$. Alors la notion de $\D$-fibration coïncide avec celle de $\D$-fibration cohomologique. Elle est donc stable par composition. Pour que $f : X \to Y$ soit une $\D$-fibration, il faut et il suffit que pour $p : Y'' \to Y'$ sur $Y$, si $p$ est une $\D$-équivalence, il en est de même de $q : X'' \to X'$.}

\noindent\textbf{Corollaire 1.} \textit{Si $f : X \to Y$ dans \textup{\Dia} est $\D$-propre, alors $f$ est une $\D$-fibration cohomologique si et seulement si pour toute flèche $y' \to y$ dans $Y$, $X/y' \to X/y$ est une $\D$-équivalence \textup{(i.e.} si et seulement si $f$ est une $\D$-fibration cohomologique faible\textup{).}}

Appliquer le critère $c$) de \texttt{[la]} prop.~6.

\msp{page 126}

\noindent\textbf{Corollaire 2.} \textit{Supposons que $\D$ satisfasse \textup{\textbf{Der} 3 $(\mathrm{c}_1, \mathrm{c}_2)$ (par exemple $\D = \text{\HOT}$).} Alors pour une flèche $f : X \to Y$ dans \textup{\Dia}, il revient au même que $f$ soit une $\D$-fibration, une $\D$-fibration cohomologique, ou une $\D$-fibration homologique. Si $f$ est propre \textup{(resp.} lisse\textup{),} il revient au même de dire que $f$ est une $\D$-fibration cohomologique \textup{(resp.} homologique\textup{),} faible.}

\dgrblock{dgrb76}{The closing axioms and the theorem they isolate}
\begin{dcomment}{status}{p126:Alors pour une flèche $f : X \to Y$ dans}{The closing axioms and the theorem they isolate}
The $c$-conditions have an antecedent different from that of the $b$-conditions. They start with a $\D$-equivalence, which may be much weaker than a weak homotopy equivalence of nerves. Thus the modern validity of $b'_2$ does not make the $c$-conditions automatic. Rational coefficients and the map $e\to BC_2$ already separate them: all the homotopy-invariance axioms hold, while invariance for constant coefficients does not imply invariance for the sign local system.

For spaces, the argument at page~110 proves all of $c_1,c_2,c',c''$. Propositions~5 and~9 therefore describe one class, and the criterion says that its members preserve weak homotopy equivalences under arbitrary pullback over their base. This is the intrinsic categorical fibration property being isolated here. It should not be identified with the class of fibrations in every model structure on categories, or with an assertion that the realization of the displayed functor is a literal Serre fibration. The invariant content is the comparison with homotopy fibres and the stability of the appropriate pullback squares.

The transition criterion has a classical source. In the dual form of \citet[\S1, Theorem~B and its corollary]{Quillen1973}, weak equivalences between the comma categories along every arrow identify these comma categories with homotopy fibres of the induced map of classifying spaces. For an opfibration, transport functors between strict fibres provide the corresponding test. Properness then gives the strict-fibre-to-comma comparison required in the manuscript. The resulting criterion is the content of Corollary~1 here; its dual gives Corollary~2. These statements were not awaiting a new homotopy group calculation: they were awaiting a coefficient formalism in which the calculation could be made functorial and self-dual.

The scope of the later theorems is worth stating exactly. The ordinary homotopy derivator exists on all small categories \citep[Theorem~6.11]{Cisinski2003}, and its locally constant objects have the invariance needed in the lemma. Within that setting, the question on pages~107--111 and the final conditional assertions are settled affirmatively. The weaker notions for other coefficient derivators retain the distinctions exposed by the examples in these pages. The manuscript's remark that $\HOT$ is the only example known to him is autobiographical, not a uniqueness theorem: for instance a product of $\HOT$ with another derivator satisfying $b'_2$ still detects nerve weak equivalences and satisfies $c''$.

The last statement follows directly. A map is an equivalence for the product derivator exactly when it is one for both factors. The $\HOT$ factor forces it to be a nerve weak equivalence. The $b'_2$ property of each factor then makes restriction an equivalence on their local systems, and local systems in a product are computed factorwise. This proves $c''$ for the product without identifying the coefficient category with spaces. What the axiom singles out is the strength of detection and transport, not a unique presentation of the theory.
\end{dcomment}


On voit donc que dans $\Cat$, on doit avoir une seule notion raisonnable de \HOT-fibration, qui coïncidera avec celle de \HOT-fibration cohomologique ou homologique. Mais pour que cette notion soit dès à présent disponible, il faudrait que le dérivateur \HOT{} soit déjà construit sur le domaine $\text{\Dia} = \Cat$ tout entier. Or ce n'est pas chose faite encore (et il n'a donc pas pu être question de prouver les propriétés de \HOT{} dont il vient d'être question, disons la plus forte \textbf{Der} 3 $c''$).

Néanmoins, je voudrais développer de façon directe (“élémentaire”) la notion de \Hot-fibration, sans me réferer à la théorie du dérivateur \HOT, et aux développements théoriques précédents. % [sic] réferer -> référer, missing accent on the second e


\dgrblock{dgrb77}{The later letter and the meaning of a fibration}
\begin{dcomment}{context}{p126}{The later letter and the meaning of a fibration}
There is a useful chronological comparison at this stopping point. In his letter to Thomason of 2 April 1991, Section~1, Grothendieck starts with a category equipped with a class $W$ of weak equivalences and defines its $W$-fibrations by requiring pullback to carry $W$-equivalences to $W$-equivalences \citep{Grothendieck1991Letter}. Closure under composition then follows by composing pullback functors. This is the strong universal preservation condition proposed on manuscript page~79, and it becomes the final condition in this chapter when the supplementary assumptions on $\D$ identify the appropriate equivalence classes.

The comparison explains why the word ``fibration'' must always be accompanied by its test class. A $W$-fibration defined by preservation of $W$ has a formal closure property. A $\D$-fibration defined by preservation of weak equivalences of nerves only as $\D$-equivalences has two different test classes, so that formal argument does not apply. The composition example on page~118 makes the difference explicit. In the case $\D=\HOT$, the classes agree and the difficulty disappears.

The earlier motivation in Section~69 of \emph{Pursuing Stacks} was to retain coherent diagram categories and the two Kan extensions as part of the coefficient theory \citep{Grothendieck1983}. The present pages show why that enrichment matters even before stability or triangulated categories enter: the coefficient produced by one direct image can have nontrivial transport, and the next direct image must retain it. The mathematical work of these pages is the progressive adjustment of the admissible coefficients and the admissible base changes until this operation can be iterated. It should be read with those successive adjustments intact.
\end{dcomment}
\begin{otherlanguage}{english}
\dgroverview{I-b}{Chapter I \textperiodcentered\ \textup{I-b} \textperiodcentered\ Overview}
\begin{otherlanguage}{english}
\noindent\textbf{Coefficients, monodromy, and the search for a usable fibration.}

These pages of Chapter~I begin with a correction to an overly strong use of cofinality and end with a programme for a theory of fibrations in the category of small categories. Their continuity lies in a single question: which coefficients must a comparison preserve if that comparison is to remain useful after another operation? Preservation of constant coefficients is enough to define a class of weak equivalences. It is not, in general, enough to carry out base change or to compose fibrations. The passage from constants to locally constant coherent coefficients is therefore the principal mathematical movement of the text.

These are pages~34--74 of the edition, manuscript pages~60--126. The text follows the editors' electronic edition \citep{Grothendieck1990Derivateurs}. Manuscript pages~60--66 also close the first part of the chapter, where their commentary, claims and readings are given. Printed and manuscript numbers are kept distinct throughout. No assertion in the commentary silently corrects a mathematical mistake in the French source: genuine transcription errors are corrected, whereas errors already present in the edition remain visible and are explained where they affect the argument.

We use the manuscript's convention that an object of $\D(X)$ has underlying diagram $X^{\mathrm{op}}\to\D(e)$. Thus the value of a right Kan extension $f_*F$ at $y$ is calculated on $X/y$, and the value of the left Kan extension $f_!F$ is calculated on $y\backslash X$. The symbols $H^*_{\D}$ and $H_*^{\D}$ denote objects in the coefficient category, not merely lists of cohomology or homology groups. This convention is responsible for several apparently reversed statements when one compares the manuscript with modern covariant-diagram treatments; the variance distinction is explained explicitly in \citet[Warning~2.2]{GrothPontoShulman2014}.

\medskip\noindent\textbf{From constant equivalences to base change.}

For $f:X\to Y$, with terminal projections $p$ and $q$, a $\D$-equivalence means that restriction induces bijections
\[
 \operatorname{Hom}_{\D(Y)}(q^*M,q^*N)
 \longrightarrow
 \operatorname{Hom}_{\D(X)}(p^*M,p^*N)
\]
for all $M,N\in\D(e)$. Adjunction and Yoneda identify this with invariance of constant-coefficient homology or cohomology whenever the relevant Kan extensions exist. It does not say that $f^*$ is fully faithful on all of $\D(Y)$. The distinction is already decisive for an endpoint inclusion in the walking arrow.

Theorem~1 on manuscript page~67 gives the stronger cofinality test actually needed: cohomological invariance for every coefficient on $Y$ is equivalent to asphericity of the comma categories $y\backslash X$. The proof uses the counit $f_!p^*M\to q^*M$. One adjunction reads it as a cohomology comparison; pointwise evaluation reads it as constant homology of a comma category. This explains the apparent use of homology inside a cohomological theorem without adding a missing hypothesis to a valid argument.

The proper and smooth base-change results then have a precise local content. Properness requires the categories of lifts of every arrow $f(x)\to y$ starting at $x$ to be $\D$-aspherical. The requirement involves both objects and their morphisms, rather than mere nonemptiness of the set of lifts. A Grothendieck opfibration satisfies it because an opcartesian lift is an initial object of the corresponding lift category. The dual account gives smoothness. The commentary supplies the omitted implications among the endpoint, lifting, and base-change tests in Theorem~2, while preserving the edition's mistyped objects and cross-references as source readings.

This vocabulary retains the geometry of proper and smooth base change, but defines a property of functors relative to a coefficient theory. The later theory of asphericity structures provides an appropriate comparison, rather than an identification with topological proper maps \citep[Theorems~1.2, 4.13, and~5.4]{Maltsiniotis2005Asphericite}. The intervening construction $\Phi(X)$, now called the twisted-arrow category, gives natural comparisons between $X$ and $X^{\mathrm{op}}$. Both endpoint projections are opfibrations with contractible fibres. Their universal equivalence proves that taking opposites preserves the equivalence class detected by $\D$. This is a comparison of homotopy information, not an equivalence of the original categories.

\medskip\noindent\textbf{What the counterexamples distinguish.}

From manuscript page~76 onward, the word ``fibration'' changes meaning as the argument develops. Its first version asks that the actual fibre and its comma-category model have the same constant-coefficient theory and that transport between comma categories preserve that theory. Requiring this after every base change is an additional condition. We exhibit both a two-object example and an example involving only finite posets where the first condition holds but fails after a strict base change.

The more substantial failure concerns the proposed implication that a universal fibration must preserve every $\D$-equivalence of bases. For rational complexes, a base can be cohomologically indistinguishable from a point while still carrying nontrivial local systems. Grothendieck's reflected-circle bundle over the circle makes the point concretely. The degree-one cohomology of its fibres is a sign local system on the base circle; its derived cohomology has differential multiplication by $-2$, hence is zero over $\mathbb Q$. The projection from the Klein bottle is consequently a rational constant-coefficient equivalence although its circle fibres are not rationally aspherical. In a fibre square the tensor square of the orientation local system becomes untwisted, and the extra cohomology reappears. The commentary calculates the full complexes and the comparison maps, including the degree-zero term omitted from the source's abbreviated group-cohomology calculation.

The finite-group version replaces this computation by the averaging idempotent. It also shows why the proposed extension to truncated homotopy theories needs care: already the point mapping into $BC_2$ is an equivalence for the represented set-valued derivator, whereas locally constant $C_2$-sets retain the action. The corresponding assertions about complete equivalences cannot be inferred from connectedness.

One further correction is necessary in the proposed finite-poset model. The action on $S^2$ is antipodal and the action on $S^1$ is a reflection. Their product acts trivially on the top class of degree three. Hence
\[
 H^*((S^2\times S^1)/C_2;\mathbb Q)
 \cong\mathbb Q\quad\text{in degrees }0,3,
\]
and vanishes in other degrees. The base $\mathbb{RP}^2$ is rationally aspherical, but this total space is not. The example still disproves invariance under inclusion of a fibre; it does not supply the additional aspherical-total-space claim. An exact rational chain calculation for the thirteen-object base and fifty-two-object total poset independently confirms these Betti numbers. The orbit-category relations printed before this construction also require correction, which is explained separately in the commentary.

\medskip\noindent\textbf{The final distinctions and their present interpretation.}

The later definitions separate constant-coefficient equivalence, invariance of cohomology with every locally constant coefficient, and full faithfulness or equivalence on the entire category of locally constant objects. The last two are called $1$- and $2$-complete equivalence. A closed internal-Hom formalism identifies cohomological equivalence with $1$-complete equivalence, under the stated compatibility assumptions. It does not identify ordinary constant-coefficient equivalence with either of them. The missing word ``cohomological'' in the sentence on manuscript page~122 is therefore a material slip.

The nerve-invariance assertions eventually formulated as the $b'$ variants of \textbf{Der~3} have a precise later answer. \citet[Theorem~1.3 and Corollary~3.8]{Cisinski2009LocallyConstant} proves invariance of locally constant coefficients under weak equivalences of indexing nerves, first in the model-category setting and then for derivators. In the space-valued case these coefficients can be read as coherent functors out of the fundamental infinity-groupoid. A weak equivalence changes neither that groupoid nor its coefficient categories. This establishes the strong form of the lemma about $\HOT$ that Grothendieck had deferred; it does not collapse the distinctions for a fixed linear coefficient theory. The minimal-localizer theorem places the ordinary weak-equivalence class inside the broader framework that motivated these pages \citep[Theorem~2.2.11]{Cisinski2004Minimal}.

Two developments in the commentary on these pages address unfinished questions directly. First, for a universally weak cohomological fibration, the formula for $f_*f^*F$ holds for every coefficient $F$ on the base, with no local-constancy assumption. A factorization category for each arrow $z\to y$ reduces the left mate of the desired comparison to constant homology. The dual proof applies to $f_!f^*F$. As a consequence, a universal $\D$-equivalence induces fully faithful restriction on the whole coefficient category. The required closure of the diagram domain is stated explicitly. These arguments use the factorization method underlying homotopy exactness \citep[Theorem~3.8]{GrothPontoShulman2014}, and the full-faithfulness conclusion for universal weak equivalences of nerves also follows from the modern homotopical epimorphism criterion \citep[Theorem~6.10]{GrothStovicek2016Trees}.

Second, the intermediate $\D$-fibrations are not closed under composition. Over $\mathbb Q$, take the opfibration over $[1]$ with transition $e\to BC_2$ and then the covering determined by the regular $C_2$-set. Both factors satisfy the universal constant-coefficient condition. Their composite has transition from two discrete points to a contractible action groupoid, and fails already in degree-zero cohomology. The stronger cohomological and homological fibrations on manuscript page~121 control locally constant coefficients on the total category and are closed under composition. A separate example with transition $BC_2\to e$ shows that this stronger relative theory has examples beyond $\HOT$-fibrations. The proposed $2$-complete variant also fails to compose for the represented set-valued derivator: a $C_2$-valued cocycle on the tetrahedral boundary face poset gives two functors between finite categories whose composite fails the required test. The obstructed restriction includes sets as trivial $C_2$-sets; its failure of essential surjectivity distinguishes this counterexample from the $1$-complete case, which composes under the internal-Hom hypotheses.

The mathematical status is thus stated at the level of the actual assertions. The main cofinality and base-change theorems hold with their conventions and domain hypotheses made explicit; the false implications have concrete counterexamples; several stated questions admit the supplied proofs or the cited later theorems. A failure to locate a particular construction in the sources checked is recorded as such, never converted into a claim that it is an open problem or a new result. The label \ours\ marks proofs and constructions given in the commentary; it carries no priority claim.
\end{otherlanguage}
\end{otherlanguage}
\dgrchunkend
